\documentclass[11pt,reqno]{amsart}
\usepackage[margin=2.5cm]{geometry}
\usepackage{amssymb,pict2e}
\usepackage[noadjust]{cite}
\usepackage[colorlinks=true,linkcolor=magenta,urlcolor=magenta,citecolor=magenta]{hyperref}
\newtheorem{thm}{Theorem}[section]
\newtheorem{prop}[thm]{Proposition}
\newtheorem{lemma}[thm]{Lemma}
\newtheorem{cor}[thm]{Corollary}
\newtheorem{defi}[thm]{Definition}
\newtheorem{quest}[thm]{Question}
\newtheorem*{thmA}{Theorem A}
\newtheorem*{thmB}{Theorem B}
\newtheorem*{thmC}{Theorem C}
\theoremstyle{definition}
\newtheorem{remark}[thm]{Remark}
\newtheorem{example}[thm]{Example}
\numberwithin{equation}{section}
\title[Spectral measure blowup]{Spectral measure blowup for Hadamard subfactors}
\author{Banica, Codex and Fable}
\begin{document}
\subjclass[2020]{46L37 (05B20, 46L65)}
\keywords{Hadamard matrix, subfactor, spectral measure}
\begin{abstract}
We study the spectral measures of the Hadamard subfactors coming from the deformed Fourier matrices $F_X\otimes_QF_Y$. After preliminaries on small-index graphs and exceptional root measures, we express the law through a Gram matrix over a compact abelian group, interpret this as an abelian covering, and compute the moment growth. At index 6 the generic law comes from the honeycomb and kagome lattices, and finite cases correspond to Eisenstein ideals. Circulant families at square index give planar random walk laws. We locate the first Jacobi defect and construct finite-depth laws with arbitrarily long generic moment sequences.
\end{abstract}
\maketitle
\setcounter{tocdepth}{1}
\tableofcontents

\section*{Introduction}

A spectral measure records moments, but often conceals the structure which produces them. For a subfactor these moments are the dimensions of the box spaces, or equivalently the numbers of closed walks of even length at the root of the principal graph. The question of spectral measure blowup is to realize the same numbers on a space where the measure becomes simpler. The circular measures of \cite{BB07,Ban09} provide the basic examples, with roots of unity and trigonometric densities replacing the atoms of the rooted law. Representation graphs and joint spectral measures offer further examples in the work of Evans and Pugh \cite{EP10,EP11,EP15}. The free orthogonal coordinate laws of \cite{BCZ09} provide another example, with an $SU_2^q$ model and an explicit circular density.

For deformed Fourier subfactors of index $MN$, with $M,N$ the orders of the 2 finite abelian groups, the models of \cite{BBi15} suggest Gram matrices of phases as blowup variables. We extend the generic formula over independent circles to arbitrary parameters, identify the compact group replacing these circles, and study its geometry and law.

We begin with 2 preliminary calculations explaining this choice of viewpoint. First, at small index $>4$, the rooted laws can still be computed entirely from finite graphs. We organize the Haagerup, Asaeda--Haagerup, $3311$ and $2221$ examples \cite{AH99,Haa94,JMS14} through their moment recurrences, and then recover their atoms and annular multiplicities. The scalar lift associated to the annular theory of Jones \cite{Jon01} remains available, but its support now includes a reciprocal pair on the positive half-line. This is a useful computation, without being a uniform model for subfactors of larger index.

Second, the exceptional diagrams lead to a different kind of measure. Their ternary indices $2,3,5$ are the inverses of the Deligne parameters of $E_6,E_7,E_8$ \cite{Del96}. We first examine how the root affects the circular law, then turn to the signed root measures of Deligne and de Man \cite{DDM96}. For Cartan powers of the adjoint representation, their atoms admit a uniform description along the exceptional series. This illustrates the usefulness of retaining root or representation data behind a scalar measure. The question of obtaining such a description for the circular laws themselves remains open.

For the main results, let $X,Y$ be finite abelian groups of orders $M,N\geq2$, and take $Q\in\mathbb T^{X\times Y}$. The deformed Fourier matrix is
$$H=F_X\otimes_QF_Y,\quad H_{ia,jb}=Q_{ib}(F_X)_{ij}(F_Y)_{ab}$$
Row and column phases of $Q$ do not change its subfactor law, so we work with the dephasing $Q_{0b}=Q_{i0}=1$. The generic case means that all the other entries are multiplicatively independent.

\begin{thmA}
Form, for each $k\in X$ and $e\in Y$, the phase matrix
$$W^{k,e}_{ij}=\frac{Q_{i+k,j+e}Q_{ij}}{Q_{i+k,j}Q_{i,j+e}}$$
and let $K_Q$ be the closed subgroup of $\mathbb T^{X\times Y}$ generated by the dephasings $(W_{ij}W_{00}/W_{i0}W_{0j})_{ij}$ of these matrices. With Haar measure on $K_Q$ and normalized matrix trace $M^{-1}Tr$, we have
$$\mu_H=\left(1-\frac1N\right)\delta_0+\frac1N law(ZZ^*),\quad Z\in K_Q$$
Consequently the box spaces satisfy
$$\dim P_n=\frac1{MN}\int_{K_Q}Tr((ZZ^*)^n)dZ$$
for every $n\geq1$.
\end{thmA}

Thus the blowup has 2 ingredients: a compact abelian group and the Gram matrix of its coordinate phases. The group is the whole dephased torus $\mathbb T^{(M-1)(N-1)}$ for generic $Q$, and is finite when all entries of $Q$ are roots of unity. At a bicharacter all the dephased second differences are trivial, giving back the Fourier law. Theorem \ref{thm:gram-all} and Proposition \ref{prop:gram-generators} establish the formula using the crossed coproduct descriptions of \cite{BBi15,Bic15,BO26}.

There is also a graph interpretation. The lattice $I_X\otimes I_Y$, formed from the augmentation ideals of $X$ and $Y$, is the first homology of the complete bipartite graph $K_{X,Y}$. The quotient determined by $Q$ gives an abelian covering, maximal in the generic case. Fourier transform turns its incidence operator into the phase matrix $Z$, and identifies $K_Q$ with the flat connections which are trivial on the relations. This is the content of Theorem \ref{thm:covering}.

\begin{thmB}
The generic moments are lower bounds for all the rooted moments. Equality at every order holds precisely when the quantum permutation group is the generic one. Otherwise the first strict inequality occurs at the length of the shortest nontrivial relation in the crossed coproduct generators. Moreover, writing $d=\dim K_Q$, we have
$$\dim P_n\sim C(MN)^nn^{-d/2}$$
with an explicit constant $C>0$, so the rooted law determines the dimension of $K_Q$.
\end{thmB}

Theorems \ref{thm:rigidity}, \ref{thm:degenerate} and \ref{thm:growth} describe these ways of recovering information from the law. There is also information which it loses: transposition preserves the rooted law of any complex Hadamard matrix \cite{BO26}, and here it transposes $K_Q$. For generic parameters with $M\neq N$, however, it changes the isomorphism class of the quantum permutation group, as shown in Theorem \ref{thm:transpose-all}.

The index 6 case gives a concrete picture of both the finite and the infinite possibilities:

\begin{thmC}
For $H=F_2\otimes_QF_3$, the relation subgroups are exactly the ideals of the ring of Eisenstein integers $\mathbb Z[\zeta]$, where $\zeta=e^{\pi i/3}$. A nonzero ideal generated by $b+c\zeta$ gives a quantum group of order $6(b^2+bc+c^2)$ and the coset graph of the torus map $\{6,3\}_{b,c}$. The zero ideal gives the honeycomb lattice and the law
$$\mu_H=\frac12\delta_0+\frac12 law(2+A)=\frac23\delta_0+\frac16p_3(|x-3|)dx$$
where $A$ is the adjacency operator of the kagome lattice and $p_3$ is the density of the length of the 3-step planar random walk, expressible by a complete elliptic integral.
\end{thmC}

These are Theorems \ref{thm:eisenstein}, \ref{thm:eisenstein-finite} and \ref{thm:density}. The finite graphs include the cube, Heawood and Pappus graphs; the infinite law explains the kagome distribution of \cite{ABF26}. At index 8 we find the diamond lattice and arbitrarily delayed exceptional moments. The triangular lattice at index 9 extends to circulant families at every square index. Following \cite{BP26}, we locate the first Jacobi defect and construct finite-depth laws with arbitrarily long generic truncations.

Sections \ref{sec:prelim} and \ref{sec:deligne} contain the preliminaries, Section \ref{sec:hadamard} recalls the quantum groups, and Section \ref{sec:fourier-blowups} proves the Gram formula. Section \ref{sec:lattices} treats coverings and index 6, Section \ref{sec:growth} treats growth and indices 8 and 9, and Section \ref{sec:circulant-jacobi} contains the circulant families and Jacobi comparisons.

\medskip

\noindent{\bf Acknowledgements.} We would like to thank C. Arhancet and S. Palcoux for useful discussions.

\section{Rooted laws, lifts and loop counts}\label{sec:prelim}

We begin with the data one can read directly from a principal graph. Walk counts determine a probability measure, the choice of a root determines its weights, and a change of variable reorganizes these weights into a lift. Our examples below index 5 will serve as a point of comparison for the matrix-valued blowups later in the paper. We write $Tr$ for the ordinary trace, displaying the dimension factor whenever a normalized trace is intended.

\subsection{From walks to a measure}

Let $(\Gamma,o)$ be a bipartite rooted graph with bounded adjacency operator $A$ on $l^2(V)$; edge multiplicities are included in $A$. Choose the even part to contain $o$ and write
$$A=\begin{pmatrix}0&\Lambda\\\Lambda^*&0\end{pmatrix},\quad K=\Lambda\Lambda^*$$
The positive operator $K$ describes 2 successive steps on the graph. This is the operator whose measure we use throughout.

\begin{defi}
The rooted law of $(\Gamma,o)$ is the probability measure $\mu$ characterized by
$$c_n=\int x^nd\mu(x)=\langle K^n\delta_o,\delta_o\rangle=\langle A^{2n}\delta_o,\delta_o\rangle$$
for $n\geq0$. Thus $c_n$ counts the walks of length $2n$ starting and ending at $o$.
\end{defi}

Equivalently, this is the image by $x\to x^2$ of the spectral measure of $A$ at the root. For a finite graph, the mass at an eigenvalue $x$ of $K$ is $\|P_x\delta_o\|^2$, where $P_x$ is its spectral projection. Eigenvalues alone therefore do not determine the rooted law.

\begin{prop}\label{prop:midpoint}
Let $p_n(v)$ count the walks of length $n$ from $o$ to $v$. The moments and their generating function satisfy
$$c_n=\sum_vp_n(v)^2,\quad \sum_{n\geq0}c_nz^n=\langle(Id-zK)^{-1}\delta_o,\delta_o\rangle$$
where the series converges for $|z|\|K\|<1$.
\end{prop}
\begin{proof}
For each possible midpoint $v$, choose the outward walk and the reversed return walk independently. There are $p_n(v)^2$ choices, giving the first formula. The second is the Neumann series for $Id-zK$, evaluated on the vector at the root.
\end{proof}

Finite graphs need only finitely many such comparisons. We record the elementary observation which will replace the diagonalization of the 4 graphs below.

\begin{lemma}\label{lem:finite-moments}
Suppose that $\Gamma$ has $m$ even vertices and that $\eta$ is a finite signed measure with at most $r$ atoms. If
$$\int x^nd\eta(x)=c_n\quad(0\leq n<m+r)$$
then $\eta$ is the rooted law of $\Gamma$.
\end{lemma}
\begin{proof}
The spectral measure of $K$ has at most $m$ atoms. Subtracting $\eta$ gives a signed measure supported at, say, $t\leq m+r$ distinct points. Its first $t$ moments vanish. The equations for its masses form a Vandermonde system at these points, which is invertible, so every mass is zero.
\end{proof}

For a finite-index irreducible inclusion $N\subset M$ of type ${\rm II}_1$ factors, the principal graph is rooted at the trivial $N$-$N$ bimodule. If $M_{-1}=N$, $M_0=M$ and $M_1,M_2,\ldots$ is the Jones tower, its box spaces and their dimensions are
$$P_n=N'\cap M_{n-1},\quad \dim P_n=c_n=\dim Hom(1,\theta^{\otimes n})$$
where $\theta$ is the dual canonical object. Multiplication by $\theta$ corresponds to 2 steps on the graph. The dual principal graph gives the same numbers, so either graph can be used to compute the subfactor law; see \cite{Jon83,Jon99,JS97,JMS14} for this background.

The positive bimodule dimension vector $w$, normalized by $w_o=1$, satisfies $Kw=[M:N]w$. In the finite case, Perron--Frobenius theory identifies the top atom and its mass as
$$x_{\max}=[M:N],\quad \mu(\{x_{\max}\})=D^{-1},\quad D=\sum_{v\ {\rm even}}w_v^2$$
Here $D$ is the global dimension of the even part. This gives a useful check on a rooted law computed from loops.

\subsection{The circle and the position of the root}\label{sec:circle}

For a law on $[0,4]$ we lift through the map
$$\sigma(u)=(u+u^{-1})^2,\quad u\in\mathbb T$$
by giving equal weights to the points of each fiber. The resulting circular measure $\varepsilon$ is uniquely specified by $\sigma_*\varepsilon=\mu$ and invariance under inversion and multiplication by $-1$. In the notation of \cite{BB07,Ban09}, let $d_n$ be uniform on the $2n$-th roots of unity and let $d'_n=2d_{2n}-d_n$, uniform on the solutions of $u^{2n}=-1$. We use the densities
$$h_j(u)=1-Re(u^{2j})$$
and write $h_jd_n$ for multiplication of $d_n$ by this density.

\begin{prop}\label{prop:path}
For the path $A_{n-1}$ on vertices $1,\ldots,n-1$, the circular measure at vertex $j$ is
$$\varepsilon_j=h_jd_n$$
For the half-line rooted at $j$, the same formula holds with Haar measure in place of $d_n$.
\end{prop}
\begin{proof}
The adjacency eigenvalues of the path and their weights at vertex $j$ are
$$2\cos(k\pi/n),\quad \frac2n\sin^2(jk\pi/n),\quad 1\leq k<n$$
On the circle the points $e^{\pm ik\pi/n}$ each have $d_n$-mass $1/(2n)$, and their density is $h_j=2\sin^2(jk\pi/n)$. Their combined weight is therefore exactly the one above. The density vanishes at $\pm1$, and invariance under $u\to-u$ gives the required measure after squaring the adjacency eigenvalue.

For a fixed moment order, a sufficiently distant truncation of the half-line cannot affect a walk from $j$ and back. Passing to the limit in the finite paths proves the Haar formula.
\end{proof}

The harmonic is therefore read from the distance of the root to the missing endpoint: moving the root changes the density, even on the same path. At $j=1$ on the half-line we obtain the Catalan moments and the free Poisson law. For comparison, the remaining $A,D$ formulas of \cite[Theorems 2.1, 3.1, 3.2 and 4.1]{BB07} are
$$\begin{array}{|c|c|c|c|c|}\hline
\Gamma&A_{n-1},\ n\geq2&\widetilde A_{2n-1},\ n\geq2&D_{n+1},\ n\geq3&\widetilde D_{n+2},\ n\geq2\\\hline
\varepsilon_\Gamma&h_1d_n&d_n&h_1d'_n&(d_n+d'_1)/2\\\hline
\end{array}$$
Here $A_{n-1}$ is rooted at an endpoint, the cycle $\widetilde A_{2n-1}$ anywhere, $D_{n+1}$ at the end of its long arm, and $\widetilde D_{n+2}$ at an extending vertex. The cycle formula follows directly by diagonalizing the cyclic shift: every Fourier vector has squared modulus $1/(2n)$ at the root. The program \texttt{zabel.py} compares all 4 formulas with walk counts for $3\leq n\leq8$. We return to the exceptional diagrams in Section \ref{sec:deligne}.

\subsection{Annular coefficients beyond the circle}

To pass the threshold 4, it is convenient to use $z=u^2$ and the map $\Psi(z)=2+z+z^{-1}$. Its fibers are still reciprocal pairs, although they need not lie on the circle.

\begin{prop}\label{prop:lift}
Every compactly supported probability measure $\mu$ on $[0,\infty)$ lifts uniquely to an inversion-invariant measure $\nu$ on $\mathbb T\cup(0,\infty)$ under $\Psi$. If $f(t)=\sum_{n\geq0}c_nt^n$ and $b_k=\int z^kd\nu(z)$, then
$$1+2\sum_{k\geq1}b_kq^k=\frac{1-q}{1+q}f\left(\frac q{(1+q)^2}\right)$$
as formal series. The fibers over $[0,4]$ lie on $\mathbb T$, and those over $(4,\infty)$ lie on the positive half-line.
\end{prop}
\begin{proof}
Solving $\Psi(z)=x$ amounts to solving $z^2-(x-2)z+1=0$. The 2 roots are inverse to each other: they are conjugate points of modulus 1 for $0\leq x\leq4$, and positive real points for $x>4$. Uniform measure on each fiber constructs $\nu$, including the singleton fibers at the endpoints. Inversion invariance forces exactly these weights.

For the moment identity, use $x=2+z+z^{-1}$ to write
$$\frac{1-q}{1+q}\cdot\frac1{1-qx/(1+q)^2}=\frac{1-q^2}{(1-qz)(1-qz^{-1})}=1+\sum_{k\geq1}(z^k+z^{-k})q^k$$
Integration and inversion invariance give the claimed coefficients. This is the transform used in \cite{BB07,Ban09}.
\end{proof}

On the circle, $\Psi(u^2)=\sigma(u)$, so this construction is compatible with the preceding circular measures. For a subfactor of index $>4$, its significance is given by the annular Temperley--Lieb decomposition of \cite{Jon01}:
\begin{equation}\label{eq:annular}
c_n=C_n+\sum_{k=1}^n\binom{2n}{n-k}a_k,\quad C_n=\frac1{n+1}\binom{2n}n
\end{equation}
The coefficient $a_k$ is the total multiplicity of lowest weight $k$, summing over rotation eigenvalues, while the weight zero term is Temperley--Lieb. Comparing this triangular expansion with the transform yields
$$a_1=2b_1+1,\quad a_k=2b_k\quad(k\geq2)$$
In particular, the moments of the lift have a representation-theoretic meaning. The first nonzero annular coefficient requires no spectral computation.

\begin{lemma}\label{prop:first-branch}
Suppose that a connected locally finite bipartite graph has no multiple edges, has exactly 1 vertex at each distance less than $n$ from its root, and has $m\geq1$ vertices at distance $n$. Then
$$c_k=C_k\quad(k<n),\quad c_n-C_n=m-1$$
For a principal graph of index $>4$, the annular coefficients satisfy $a_k=0$ for $k<n$ and $a_n=m-1$.
\end{lemma}
\begin{proof}
A closed walk which reaches distance $n$ must use at least $2n$ steps. Before this order, every walk is a Dyck path on the initial chain, giving $C_k$. At order $n$, the only new walks travel straight to a vertex at distance $n$ and straight back. There are $m$ such choices, in place of the single choice on the half-line. All other walks remain unchanged, proving the difference $m-1$. Substitution in \eqref{eq:annular}, starting at $k=1$, gives the annular assertion.
\end{proof}

\subsection{Finite computations between indices 4 and 5}

The classification in \cite{JMS14} lists 5 non-Temperley--Lieb principal graph pairs in the interval $(4,5)$, up to duality. We use 4 of them: Haagerup and Asaeda--Haagerup \cite{AH99,Haa94}, the $3311$ example of Goodman, de la Harpe and Jones \cite{GHJ89}, and the $2221$ example of Izumi \cite{Izu01}. Figure \ref{fig:graphs} fixes the roots. In the last 2 pairs, the graphs coincide and their 4 arms have the indicated lengths.

\begin{figure}[htbp]
\centering
\setlength{\unitlength}{5mm}
\begin{tabular}{lll}
&principal graph&dual principal graph\\
Haagerup&\begin{picture}(6.3,2.2)(-0.15,-0.6)
\put(0,0){\line(1,0){3}}\put(3,0){\line(2,1){1}}\put(3,0){\line(2,-1){1}}\put(4,0.5){\line(1,0){2}}\put(4,-0.5){\line(1,0){2}}
\multiput(0,0)(1,0){4}{\circle*{0.2}}\multiput(4,0.5)(1,0){3}{\circle*{0.2}}\multiput(4,-0.5)(1,0){3}{\circle*{0.2}}
\end{picture}&\begin{picture}(5.3,2.2)(-0.15,-0.6)
\put(0,0){\line(1,0){3}}\put(3,0){\line(2,1){1}}\put(3,0){\line(2,-1){1}}\put(4,0.5){\line(1,0){1}}\put(4,0.5){\line(1,-1){1}}
\multiput(0,0)(1,0){4}{\circle*{0.2}}\multiput(4,0.5)(1,0){2}{\circle*{0.2}}\multiput(4,-0.5)(1,0){2}{\circle*{0.2}}
\end{picture}\\
Asaeda--Haagerup&\begin{picture}(10.3,2.2)(-0.15,-0.6)
\put(0,0){\line(1,0){5}}\put(5,0){\line(2,1){1}}\put(5,0){\line(2,-1){1}}\put(6,0.5){\line(1,0){1}}\put(7,0.5){\line(2,1){1}}\put(7,0.5){\line(2,-1){1}}\put(6,-0.5){\line(1,0){4}}
\multiput(0,0)(1,0){6}{\circle*{0.2}}\multiput(6,0.5)(1,0){2}{\circle*{0.2}}\put(8,1){\circle*{0.2}}\put(8,0){\circle*{0.2}}\multiput(6,-0.5)(1,0){5}{\circle*{0.2}}
\end{picture}&\begin{picture}(9.3,2.2)(-0.15,-0.6)
\put(0,0){\line(1,0){5}}\put(5,0){\line(2,1){1}}\put(5,0){\line(2,-1){1}}\put(6,0.5){\line(1,0){3}}\put(6,0.5){\line(1,-1){1}}
\multiput(0,0)(1,0){6}{\circle*{0.2}}\multiput(6,0.5)(1,0){4}{\circle*{0.2}}\multiput(6,-0.5)(1,0){2}{\circle*{0.2}}
\end{picture}\\
$3311$&\begin{picture}(6.3,2.2)(-0.15,-0.6)
\put(0,0){\line(1,0){6}}\put(3,0){\line(2,1){1}}\put(3,0){\line(2,-1){1}}
\multiput(0,0)(1,0){7}{\circle*{0.2}}\put(4,0.5){\circle*{0.2}}\put(4,-0.5){\circle*{0.2}}
\end{picture}&\begin{picture}(6.3,2.2)(-0.15,-0.6)
\put(0,0){\line(1,0){6}}\put(3,0){\line(2,1){1}}\put(3,0){\line(2,-1){1}}
\multiput(0,0)(1,0){7}{\circle*{0.2}}\put(4,0.5){\circle*{0.2}}\put(4,-0.5){\circle*{0.2}}
\end{picture}\\
$2221$&\begin{picture}(4.3,2.2)(-0.15,-0.6)
\put(0,0){\line(1,0){4}}\put(2,0){\line(2,1){1}}\put(3,0.5){\line(1,0){1}}\put(2,0){\line(2,-1){1}}
\multiput(0,0)(1,0){5}{\circle*{0.2}}\multiput(3,0.5)(1,0){2}{\circle*{0.2}}\put(3,-0.5){\circle*{0.2}}
\end{picture}&\begin{picture}(4.3,2.2)(-0.15,-0.6)
\put(0,0){\line(1,0){4}}\put(2,0){\line(2,1){1}}\put(3,0.5){\line(1,0){1}}\put(2,0){\line(2,-1){1}}
\multiput(0,0)(1,0){5}{\circle*{0.2}}\multiput(3,0.5)(1,0){2}{\circle*{0.2}}\put(3,-0.5){\circle*{0.2}}
\end{picture}
\end{tabular}
\caption{The 4 pairs of principal graphs, with the root at the left.}
\label{fig:graphs}
\end{figure}

We start with integer sequences. This keeps the computation rational until the final extraction of the atoms, and treats the 4 examples in the same way.

\begin{prop}\label{prop:four-moments}
For each graph in the following table, define $r_0,r_1$ as indicated and put
$$r_{n+2}=sr_{n+1}-pr_n$$
Its rooted moments, with $c_0=1$, are given for $n\geq1$ by
$$\begin{array}{|c|c|c|c|}\hline
\Gamma&(s,p)&(r_0,r_1)&c_n\\\hline
{\rm Haagerup}&(5,3)&(2,2)&(2^{n+1}+r_n)/6\\\hline
\mbox{Asaeda--Haagerup}&(5,2)&(2,1)&(2+2^n+3^n+r_n)/8\\\hline
3311&(6,6)&(4,6)&(3\cdot2^n+r_n)/12\\\hline
2221&(5,1)&(2,2)&(4+r_n)/6\\\hline
\end{array}$$
Thus, apart from the displayed elementary powers, only a recurrence of order 2 is needed.
\end{prop}
\begin{proof}
Let $x_\pm$ be the roots of $t^2-st+p$. They are distinct and positive in all 4 cases. The recurrence has the solution
$$r_n=\frac{r_1-r_0x_-}{x_+-x_-}x_+^n+\frac{r_0x_+-r_1}{x_+-x_-}x_-^n$$
Hence each proposed sequence is the moment sequence of a finite signed measure: the displayed powers give its rational atoms, and $x_\pm$ give 2 more atoms. Add an atom at zero if needed to make the zeroth moment equal to 1. The resulting support sizes are $r=4,6,4,3$.

Use the dual graphs for Haagerup and Asaeda--Haagerup, and the principal graphs for $3311$ and $2221$. They have respectively $m=4,6,6,4$ even vertices. Lemma \ref{lem:finite-moments} reduces the proof to walk counts through orders $7,11,9,6$, respectively. Starting with $p_0=\delta_o$, the updates $p_{n+1}=Ap_n$ and $c_n=\sum_vp_n(v)^2$ give exactly the proposed values at these orders. The program \texttt{riis.py} performs these finite comparisons in integer arithmetic on both graphs of each pair, and checks the recurrences through order 40.
\end{proof}

The radicals enter only when we solve these 4 quadratic recurrences. They give the following description of the laws.

\begin{thm}\label{thm:four-laws}
The rooted laws of the Haagerup, Asaeda--Haagerup, $3311$ and $2221$ subfactors are, in this order,
\begin{eqnarray*}
\mu&=&\frac13\delta_0+\frac13\delta_2+\sum_\pm\frac16\left(1\mp\frac3{\sqrt{13}}\right)\delta_{(5\pm\sqrt{13})/2}\\
\mu&=&\frac14\delta_0+\frac14\delta_1+\frac18\delta_2+\frac18\delta_3+\sum_\pm\frac18\left(1\mp\frac4{\sqrt{17}}\right)\delta_{(5\pm\sqrt{17})/2}\\
\mu&=&\frac5{12}\delta_0+\frac14\delta_2+\sum_\pm\frac{2\mp\sqrt3}{12}\delta_{3\pm\sqrt3}\\
\mu&=&\frac23\delta_1+\sum_\pm\frac16\left(1\mp\frac3{\sqrt{21}}\right)\delta_{(5\pm\sqrt{21})/2}
\end{eqnarray*}
where the sum over $\pm$ includes both conjugate terms.
\end{thm}
\begin{proof}
In the first row of Proposition \ref{prop:four-moments}, the coefficients of $x_+^n$ and $x_-^n$ in $r_n$ are $1-3/\sqrt{13}$ and $1+3/\sqrt{13}$. Dividing by 6 gives the quadratic masses, while $2^{n+1}/6$ gives mass $1/3$ at 2. Their total mass is $2/3$, leaving mass $1/3$ at zero. The other rows work in exactly the same way: solve the recurrence using $r_0,r_1$, divide by the denominator in the table, and restore total mass 1 at zero. In the $2221$ case no mass is left at zero.

For Asaeda--Haagerup, the quadratic contribution before division by 8 is
$$r_n=\left(1-\frac4{\sqrt{17}}\right)\left(\frac{5+\sqrt{17}}2\right)^n+\left(1+\frac4{\sqrt{17}}\right)\left(\frac{5-\sqrt{17}}2\right)^n$$
The atoms at $1,2,3$ have total mass $1/2$ and the quadratic pair has mass $1/4$, leaving mass $1/4$ at zero. Lemma \ref{lem:finite-moments} identifies these measures with the rooted laws, and \texttt{riis.py} checks their radical weights by exact arithmetic.
\end{proof}

Only the top atom lies above 4. It is the index, with weight $D^{-1}$, and Proposition \ref{prop:lift} replaces it by 2 atoms of weight $1/(2D)$ at $z,z^{-1}$, where $z>1$ and $z+z^{-1}=[M:N]-2$. All remaining mass lifts to the circle. Combining this with Lemma \ref{prop:first-branch} gives the following data:
$$\begin{array}{|c|c|c|c|c|}\hline
&{\rm Haagerup}&\mbox{Asaeda--Haagerup}&3311&2221\\\hline
{\rm index}&(5+\sqrt{13})/2&(5+\sqrt{17})/2&3+\sqrt3&(5+\sqrt{21})/2\\\hline
D&(39+9\sqrt{13})/2&136+32\sqrt{17}&24+12\sqrt3&(21+3\sqrt{21})/2\\\hline
z&1.7221&2.0810&2.2966&2.3692\\\hline
(n,a_n)&(4,1)&(6,1)&(4,2)&(3,2)\\\hline
\end{array}$$
The moments, and the annular multiplicities given by \eqref{eq:annular}, start as follows:
$$\begin{array}{|l|c|r|r|r|r|r|r|r|r|r|r|}\hline
&n&1&2&3&4&5&6&7&8&9&10\\\hline
{\rm Haagerup}&c_n&1&2&5&15&52&199&807&3374&14321&61227\\\cline{2-12}
&a_n&0&0&0&1&0&1&0&4&3&6\\\hline
\mbox{Asaeda--Haagerup}&c_n&1&2&5&14&42&133&443&1552&5716&22067\\\cline{2-12}
&a_n&0&0&0&0&0&1&0&2&2&6\\\hline
3311&c_n&1&2&5&16&62&268&1220&5680&26696&125968\\\cline{2-12}
&a_n&0&0&0&2&0&4&7&18&39&92\\\hline
2221&c_n&1&2&7&31&146&697&3337&15986&76591&366967\\\cline{2-12}
&a_n&0&0&2&1&4&11&24&56&137&320\\\hline
\end{array}$$
The computations therefore separate 2 questions. The scalar lift and its annular coefficients remain explicit above index 4; the choice of a useful model at arbitrary index is another matter. The Hadamard examples in the following sections replace the varying scalar support by a compact phase group and a Gram matrix.

\section{Exceptional graphs and root measures}\label{sec:deligne}

The ternary indices $2,3,5$ of the exceptional graphs are the inverses of their Deligne parameters. We examine the graph weights first, then compute a uniform family of signed root measures. This separates the 2 calculations underlying the question at the end of the section.

\subsection{Root choices and ternary grids}

Retain $d_n,d'_n,h_j$ from Section \ref{sec:prelim}. Let $d''_l$ be uniform on the points $e^{2\pi ij/(12l)}$ with $j\equiv\pm1$ mod $6$, and let $d'''_1=\frac12(3d_3-d_1)$ be uniform on the sixth roots other than $\pm1$. The smallest exceptional graph illustrates the dependence on the root.

\begin{prop}\label{prop:e6-roots}
In $E_6$, let $s$ be the tip of the short arm and $a$ the tip of an arm of length 2. Their rooted laws satisfy
$$\mu_s=\frac12\left(1+\frac1{\sqrt3}\right)\delta_{2-\sqrt3}+\frac12\left(1-\frac1{\sqrt3}\right)\delta_{2+\sqrt3},\quad \mu_a=\frac12\mu_s+\frac12\delta_1$$
Their circular lifts are
$$\varepsilon_s=\frac13(1+2h_1)d''_2,\quad \varepsilon_a=\frac12\varepsilon_s+\frac12d'''_1$$
so changing the tip removes one component of the circular measure and rescales the other.
\end{prop}
\begin{proof}
Exchange the 2 long arms and decompose the vertex space into the symmetric and antisymmetric subspaces. The antisymmetric part is a 2-vertex path, so the square of its adjacency is the identity. The symmetric part is the weighted path with edge weights $1,\sqrt2,1$. Reflection of this path exchanges $e_s$ and the normalized sum of the 2 long-arm tips.

The vector $e_a$ has squared norm $1/2$ in each subspace, giving $\mu_a=(\mu_s+\delta_1)/2$. The square of the weighted path, restricted to the parity class of $s$, has diagonal entries $1,3$ and off-diagonal entries $\sqrt2$. Its eigenvalues are $2\pm\sqrt3$, whose weights follow from total mass 1 and first moment 1.

On the 2 fibers of $d''_2$, the density $(1+2h_1)/3$ is $1\mp1/\sqrt3$ at $\sigma(u)=2\pm\sqrt3$. Each fiber has $d''_2$-mass $1/2$, and $d'''_1$ lifts $\delta_1$. This gives the circular formulas.
\end{proof}

We now use the long-arm tip in $E_6$, and the tip of the longest arm in $E_7,E_8$. The Coxeter numbers $h$, indices and initial moments at these tips are
$$\begin{array}{|c|c|c|c|l|}\hline
&h&l&{\rm index}&c_1,c_2,\ldots,c_8\\\hline
E_6&12&2&4\cos^2(\pi/12)&1,2,6,21,77,286,1066,3977\\\hline
E_7&18&3&4\cos^2(\pi/18)&1,2,5,15,51,186,702,2691\\\hline
E_8&30&5&4\cos^2(\pi/30)&1,2,5,14,43,143,506,1870\\\hline
\end{array}$$
The eigenvalues $2\cos(\pi m/h)$, with $m$ running over the exponents, place the circular support on the grid of $2h=12l$ roots of unity. This explains $l=h/6$; the weights still depend on the root. The formulas of \cite{Ban09} are as follows.

\begin{thm}[\cite{Ban09}]\label{thm:egrids}
The circular measures at these roots are
$$\varepsilon_{E_6}=\frac16(d''_2+2h_1d''_2+3d'''_1),\quad \varepsilon_{E_7}=\frac13(2h_2d''_3+d'_1),\quad \varepsilon_{E_8}=\frac13(2(h_1+h_3)d''_5-d''_1)$$
In terms of the binary grids, the same measures are
$$\varepsilon_{E_6}=h_1d_{12}+\frac12(d_{12}-d_6-d_4+d_3),\quad \varepsilon_{E_7}=h_2d'_9+\frac12(d'_1-d'_3),\quad \varepsilon_{E_8}=(h_1+h_3)d'_{15}-\frac12(d'_5+d'_3)$$
The respective ternary indices are $l=2,3,5$.
\end{thm}
\begin{proof}
The first $E_6$ formula follows from Proposition \ref{prop:e6-roots}. The remaining identities are \cite[Theorems 7.1 and 8.7]{Ban09}, with the binary $E_6$ expression also given in \cite[Theorem 6.2]{BB07}. Passing between the 2 presentations amounts to separating root-of-unity grids into their congruence classes. All 6 expressions are checked against graph walks through order 12 by \texttt{zabel.py}.
\end{proof}

\subsection{Signed root measures and their moments}

The exceptional series of \cite{Del96} uses the following parameter $\mu$, expressed through the dual Coxeter number $h^\vee$:
$$\begin{array}{|c|c|c|c|c|c|c|c|c|}\hline
\mathfrak g&A_1&A_2&G_2&D_4&F_4&E_6&E_7&E_8\\\hline
h^\vee&2&3&4&6&9&12&18&30\\\hline
\mu=6/h^\vee&3&2&3/2&1&2/3&1/2&1/3&1/5\\\hline
\dim\mathfrak g&3&8&14&28&52&78&133&248\\\hline
\end{array}$$
For example, the adjoint dimension is given throughout the table by
\begin{equation}\label{eq:deligne-dim}
\dim\mathfrak g=\frac{2(5-\mu)(6+\mu)}{\mu(1+\mu)}
\end{equation}
Similar formulas describe small tensor powers \cite{Del96,LM02}. We recover this one from the signed root measures of Deligne and de Man \cite{DDM96}.

Let $\Delta_+$ be the positive roots, $\rho$ their half-sum and $\theta$ the highest root. Normalize the pairing on weights by the Killing form, so $(\theta,\theta)=1/h^\vee$. For a dominant integral weight $\lambda$, define
$$L_{\mathfrak g}(\lambda)=\sum_{\alpha\in\Delta_+}(\delta_{12(\lambda+\rho,\alpha)}-\delta_{12(\rho,\alpha)})$$
Each root moves a unit atom to the right. The resulting measure has total mass zero and support in $(0,\infty)$. In the normalization $(\cdot,\cdot)_2$ where long roots have squared length 2, the atom locations are $\mu(\lambda+\rho,\alpha)_2$ and $\mu(\rho,\alpha)_2$.

\begin{prop}\label{prop:shift}
For $L=L_{\mathfrak g}(\lambda)$, the function $\xi(x)=-L((-\infty,x])$ is a sum of interval indicators, and
$$\xi(x)=\sum_{\alpha\in\Delta_+}1_{[12(\rho,\alpha),12(\lambda+\rho,\alpha))}(x),\quad \int fdL=\int f'(x)\xi(x)dx$$
for $f\in C^1(0,\infty)$. Its logarithmic and first 2 ordinary moments are
$$\int\log xdL(x)=\log\dim V_\lambda,\quad \int xdL(x)=24(\lambda,\rho),\quad \int x^2dL(x)=72c_\lambda$$
where $c_\lambda=(\lambda,\lambda+2\rho)$ is the quadratic Casimir eigenvalue, equal to 1 on the adjoint representation.
\end{prop}
\begin{proof}
Dominance gives $12(\lambda+\rho,\alpha)\geq12(\rho,\alpha)>0$. Each difference $\delta_b-\delta_a$ therefore contributes the interval $[a,b)$ to $\xi$, and $f(b)-f(a)$ is the integral of $f'$ on that interval. This proves the integral formula and $\xi\geq0$.

The logarithmic identity is the Weyl dimension formula: $\dim V_\lambda=\prod_{\alpha\in\Delta_+}(\lambda+\rho,\alpha)/(\rho,\alpha)$. The linear identity uses $\sum_{\alpha\in\Delta_+}\alpha=2\rho$. For the quadratic identity, the Killing normalization gives $\sum_{\alpha\in\Delta_+}(v,\alpha)^2=(v,v)/2$, and hence
$$72((\lambda+\rho,\lambda+\rho)-(\rho,\rho))=72(\lambda,\lambda+2\rho)$$
This proves the remaining identities.
\end{proof}

Thus $L$ compares 2 ordered positive diagonal operators. For $\lambda\neq0$, normalizing $\xi(x)dx$ gives a probability measure; retaining the signs instead makes dimension and Casimir simple integrals.

One can also retain the Casimir through the heat measure of the compact group. If $V$ is the adjoint representation, $D=\dim V$, and $V^{\otimes n}=\bigoplus_\lambda m_{n,\lambda}V_\lambda$, the image $\nu_t$ of heat measure under $\chi_V/D$ satisfies
$$\int x^nd\nu_t(x)=D^{-n}\sum_\lambda m_{n,\lambda}e^{-tc_\lambda}\dim V_\lambda$$
The limit $t\to\infty$ gives the Haar character moments. These representation measures provide models for the graph question, without identifying them with the circular laws.

\subsection{A uniform calculation for Cartan powers}

Take $\lambda=k\theta$, so that $V_{k\theta}$ is the $k$-th Cartan power of the adjoint representation. A grading by the highest coroot reduces the root calculation to one repeated translation. Write $x_\alpha=12(\rho,\alpha)$ and let $\tau_a$ denote translation by $a$.

\begin{lemma}\label{lem:levels}
For any simple Lie algebra, the levels $\langle\alpha,\theta^\vee\rangle$ of positive roots are $0,1,2$. Only $\theta$ has level 2, and $2h^\vee-4$ roots have level 1. Put
$$U_{\mathfrak g}=\sum_{\langle\alpha,\theta^\vee\rangle=1}(\delta_{x_\alpha+\mu}-\delta_{x_\alpha}),\quad \mu=6/h^\vee$$
Then the Cartan measures satisfy
$$L_{\mathfrak g}(k\theta)=\sum_{j=0}^{k-1}\tau_{j\mu}U_{\mathfrak g}+\delta_{6+(2k-1)\mu}-\delta_{6-\mu}$$
for every integer $k\geq1$.
\end{lemma}
\begin{proof}
The highest root is dominant and long. Hence the integer $2(\alpha,\theta)/(\theta,\theta)$ is at least zero, and Cauchy--Schwarz bounds it by 2, with equality only at $\alpha=\theta$. Summing these levels over $\Delta_+$ gives $2\langle\rho,\theta^\vee\rangle=2h^\vee-2$. Subtracting the contribution 2 of $\theta$ gives the asserted count at level 1.

Adding $k\theta$ moves an atom by $k\mu$ times its level. Level zero contributes nothing. At level 1, split the displacement into $k$ successive steps of size $\mu$: the intermediate atoms cancel, leaving the translated copies of $U_{\mathfrak g}$. At level 2, the starting point is $x_\theta=\mu(h^\vee-1)=6-\mu$ and the displacement is $2k\mu$.
\end{proof}

The level-one measure is particularly short on the exceptional series. Although the number of roots increases, most consecutive endpoints cancel.

\begin{prop}\label{prop:level-one}
For the Lie algebras in the exceptional series, we have
$$U_{\mathfrak g}=\delta_4+\delta_{5-\mu}+\delta_{6-\mu}-\delta_\mu-\delta_{1+\mu}-\delta_2$$
with coincident atoms combined at each value of $\mu$.
\end{prop}
\begin{proof}
List the starting points $x_\alpha$ of the level-one shifts. For $A_1$ there are none and the displayed expression cancels to zero. For $G_2$ the list is $3/2,2,5/2,3$, with step $\mu=3/2$. For the remaining types, the values of $x_\alpha/\mu$ are given below; exponents denote multiplicities.
$$\begin{array}{|c|l|}\hline
A_2&1^2\\\hline
D_4&1,2^3,3^3,4\\\hline
F_4&1,2,5/2,3^2,7/2,4^2,9/2,5^2,11/2,6,7\\\hline
E_6&1,2,3^2,4^3,5^3,6^3,7^3,8^2,9,10\\\hline
E_7&1,2,3,4^2,5^2,6^3,7^3,\ldots,11^3,12^2,13^2,14,15,16\\\hline
E_8&1,\ldots,5,6^2,\ldots,9^2,10^3,\ldots,19^3,20^2,\ldots,23^2,24,\ldots,28\\\hline
\end{array}$$
For example, at $A_2$ the 2 starting points are both 2, so $U_{\mathfrak g}=2\delta_4-2\delta_2$, as claimed. In each other row, shifting the list by one and subtracting the original list leaves exactly the 6 endpoints in the statement. The lists can be generated from the Cartan matrices by successive root reflections; \texttt{zabel.py} performs this enumeration and the cancellations in exact rational arithmetic.
\end{proof}

Except at $A_1,G_2$, the lists are 3 progressions of step $\mu$, from $\mu$ to $4-\mu$, from $1+\mu$ to $5-2\mu$, and from $2$ to $6-2\mu$. Their starts and the points just after their ends give the 6 surviving atoms. Translation now gives the Cartan formulas.

\begin{thm}\label{thm:cartan}
For $\mathfrak g$ in the exceptional series and $k\geq1$, the Cartan dimension is
$$\dim V_{k\theta}=\frac{6+(2k-1)\mu}{6+(k-1)\mu}\prod_{j=1}^k\frac{(4+(j-1)\mu)(5+(j-2)\mu)(6+(j-1)\mu)}{j\mu(1+j\mu)(2+(j-1)\mu)}$$
The corresponding signed measure is
$$L_{\mathfrak g}(k\theta)=\sum_{j=1}^k(\delta_{4+(j-1)\mu}+\delta_{5+(j-2)\mu}-\delta_{j\mu}-\delta_{1+j\mu}-\delta_{2+(j-1)\mu})+\sum_{j=1}^{k-1}\delta_{6+(j-1)\mu}+\delta_{6+(2k-1)\mu}$$
Its first moment is $2k(6-\mu)$ and its second moment is $12k(6+(k-1)\mu)$. At $k=1$ the dimension formula is \eqref{eq:deligne-dim}.
\end{thm}
\begin{proof}
Translate the 6 atoms of Proposition \ref{prop:level-one} by $j\mu$, for $0\leq j<k$, and add the highest-root contribution from Lemma \ref{lem:levels}. The positive atom $\delta_{6-\mu}$ in the first translate cancels the negative highest-root atom. Reindexing the remaining terms gives the stated signed measure.

Its logarithmic integral is the logarithm of the product of the positive locations divided by the product of the negative locations, with multiplicities. Proposition \ref{prop:shift} identifies this ratio with the dimension. To write it as the displayed product, include the missing factor $6+(k-1)\mu$ in the product and divide by the same factor outside it.

For the ordinary moments, use $x_\theta=6-\mu$ and
$$24(k\theta,\rho)=2k(6-\mu),\quad c_{k\theta}=\frac{k^2}{h^\vee}+\frac{k(6-\mu)}6$$
Proposition \ref{prop:shift} gives the 2 formulas immediately.
\end{proof}

The product is the Cartan dimension formula of Landsberg and Manivel \cite{LM02}; the signed measure is the restriction of the Deligne--de Man interpolation \cite{DDM96} to the Cartan ray. The first values, checked independently by the Weyl dimension formula in \texttt{zabel.py}, are
$$\begin{array}{|c|c|c|c|c|c|c|c|c|}\hline
k&A_1&A_2&G_2&D_4&F_4&E_6&E_7&E_8\\\hline
1&3&8&14&28&52&78&133&248\\\hline
2&5&27&77&300&1053&2430&7371&27000\\\hline
3&7&64&273&1925&12376&43758&238602&1763125\\\hline
\end{array}$$
For $k=1$ the entire signed measure reduces to
$$L_{\mathfrak g}(\theta)=\delta_4+\delta_{5-\mu}+\delta_{6+\mu}-\delta_\mu-\delta_{1+\mu}-\delta_2$$
Its 6 locations are affine functions of $\mu$, and its second moment is always 72. This is the kind of uniform description that motivates a return to the graph question.

Indeed, $E_6,E_7,E_8$ are simply laced, so $h=h^\vee$ and their ternary indices satisfy $l=h/6=1/\mu$. The equality identifies the grid parameters, but does not determine the root weights of Theorem \ref{thm:egrids}.

\begin{quest}\label{q:exceptional}
Is there a uniform root-theoretic description of the circular laws of $E_6,E_7,E_8$, with their specified roots, at
$$\mu\in\{1/2,1/3,1/5\}$$
or at least of the masses of their rooted laws? Can such a description extend to a natural family of measures along the whole exceptional series?
\end{quest}

The desired construction should explain the weights and their dependence on the root, as the signed Cartan construction explains its affine atoms.

\section{Deformed Fourier matrices, quantum groups}\label{sec:hadamard}

Throughout the remaining sections, $X,Y$ are finite abelian groups of orders $M,N\geq2$, written additively. We recall the quantum group construction and the crossed coproduct descriptions of \cite{BBi15,Bic15,BO26}, leading to Theorem \ref{thm:dita-all} and the word count of Theorem \ref{thm:dita-words}. Commuting-square methods also give isolation criteria \cite{Nic06}, relative commutant calculations and intermediate subfactors \cite{Nic10}. Group-type descriptions appear in \cite{Bur15}, while \cite{Mon23} uses colored planar algebras and spectral criteria to study infinite depth.

\subsection{Hadamard matrices and quantum groups}\label{sec:prelim-cqg}

A compact quantum group is described by a unital $C^*$-algebra $C(G)$ with a coassociative comultiplication and the usual cancellation conditions. Its Haar state $h$ satisfies
$$(h\otimes Id)\Delta(a)=(Id\otimes h)\Delta(a)=h(a)1$$
For a unitary representation $U$ on a finite-dimensional space, the character is $\chi(U)=(Tr\otimes Id)(U)$. The characters satisfy $\chi(U\otimes V)=\chi(U)\chi(V)$, and Haar averaging gives the orthogonal projection onto $Fix(U)$. In particular,
$$h(\chi(U)^n)=\dim Fix(U^{\otimes n})$$
We use the representation-theoretic conventions of \cite{NT13}.

\begin{defi}
A complex Hadamard matrix of size $k$ is a matrix $H\in M_k(\mathbb T)$ whose rows are pairwise orthogonal. A magic unitary is a square matrix of projections whose entries in each row and each column sum to one. The universal algebra generated by a $k\times k$ magic unitary is $C(S_k^+)$, with
$$\Delta(v_{ij})=\sum_av_{ia}\otimes v_{aj}$$
For the rows $H_i$ of $H$, put
$$\xi_{ij}=\frac1{\sqrt k}\cdot\frac{H_i}{H_j},\quad
p_{ij}=Proj(\xi_{ij})$$
where the quotients are coordinatewise. The Hopf image $G_H$ is the smallest compact quantum subgroup of $S_k^+$ through which the representation $v_{ij}\to p_{ij}$ factors.
\end{defi}

Associated to $H$ is a subfactor of index $k$, called Hadamard subfactor, which comes from the commuting square formed by the diagonal algebra of $M_k(\mathbb C)$ and by its conjugate by $H$ \cite{Jon99,JS97}. Its box spaces, and so its rooted law, can be read on $G_H$, as follows:

\begin{prop}\label{prop:hadamard-law}
The matrix $(p_{ij})$ is magic. If $u$ is the fundamental representation of $G_H$, then the rooted law of the Hadamard subfactor is the Haar law of the positive main character
$$\chi_H=\sum_i u_{ii},\quad
c_n=h(\chi_H^n)=\dim Fix(u^{\otimes n})$$
for every $n\geq0$.
\end{prop}
\begin{proof}
For fixed $j$, the vectors $\xi_{ij}$ are an orthonormal basis, by row orthogonality of $H$. For fixed $i$, the same follows by conjugating row orthogonality. This proves the magic relations. The identification of the Hadamard box spaces with the invariant tensors of the Hopf image is the construction of \cite{BBS09}. Haar averaging gives their dimensions. Positivity follows because every $u_{ii}$ is a projection, and $0\leq\chi_H\leq k$.
\end{proof}

\begin{remark}
The projection model also gives a direct finite-matrix check of the moments. For multi-indices $i,j\in\{1,\ldots,k\}^n$, put
$$(T_n)_{i,j}=\frac1k Tr(p_{i_1j_1}\ldots p_{i_nj_n})$$
The Ces\`aro averages of $T_n^r$ converge to the Haar projection onto $Fix(u^{\otimes n})$. Hence $c_n$ is the multiplicity of the eigenvalue one of $T_n$. Writing $p_{ij}=\xi_{ij}\xi_{ij}^*$ makes each trace a cyclic product of inner products of the vectors $\xi_{ij}$. The script \texttt{armstrong.py} uses this formula for sample $F_2\otimes_QF_3$ and $F_3\otimes_QF_2$ matrices through degree 3, and for $F_2\otimes_QF_2$ with phase orders $3,4,8$ through degree 5. The eigenvalue comparison uses tolerance $10^{-6}$.
\end{remark}

A representation is inner faithful when its Hopf image is the whole quantum group. In the Fourier setting the crossed coproduct model reduces the computation to a subgroup $N_Q$ of a free abelian lattice.

\begin{defi}
A Gram blowup of a rooted law at index $MN$ is a compact group $K$, with Haar probability measure, and a continuous matrix-valued function $Z:K\to M_{M,N}(\mathbb T)$ such that
$$\mu=\left(1-\frac1N\right)\delta_0+
\frac1N law(ZZ^*)$$
The matrix law is taken with respect to Haar measure and the normalized trace $M^{-1}Tr$.
\end{defi}

The positive moments of this expression are $(MN)^{-1}\int_K Tr((ZZ^*)^n)$. Its first moment is one because $Tr(ZZ^*)=MN$. This viewpoint comes from \cite{BBi15}. The function algebra $C(K)\otimes M_M(\mathbb C)$ carries the matrix variable and the state giving its law. The character can therefore be studied before its matrix spectral decomposition is pushed down to a scalar measure.

\subsection{Di\c t\u a deformations and the generic Gram law}

Choose a decomposition of $X$ into cyclic groups and let $F_X$ be the tensor product of the corresponding Fourier matrices $F_n=(w^{jk})_{jk}$, with $w=e^{2\pi i/n}$ and $j,k\in\{0,\ldots,n-1\}$. By \cite[Proposition 3.1 (2) and (3)]{BBi15}, $F_X$ is the matrix of the Fourier transform of $X$ and $(F_X)_{ij}$ is a bicharacter of $i,j\in X$, which is symmetric and nondegenerate. The matrix $F_Y$ is defined from a decomposition of $Y$ into cyclic groups. For a matrix $Q\in\mathbb T^{X\times Y}$, the Di\c t\u a deformation \cite{Dit04} of $F_X\otimes F_Y$ is the matrix $H=F_X\otimes_Q F_Y$ indexed by $X\times Y$ with entries
$$H_{ia,jb}=Q_{ib}(F_X)_{ij}(F_Y)_{ab},
\quad i,j\in X,\ a,b\in Y$$
For $(i,a),(k,c)\in X\times Y$ we have
$$\sum_{j\in X,b\in Y}H_{ia,jb}\bar H_{kc,jb}
=\sum_{b\in Y}Q_{ib}\bar Q_{kb}(F_Y)_{ab}(\bar F_Y)_{cb}\sum_{j\in X}(F_X)_{ij}(\bar F_X)_{kj}
=MN\delta_{ik}\delta_{ac}$$
by the orthogonality of the rows of $F_X$ and then of $F_Y$. Thus $H$ is a complex Hadamard matrix of size $MN$, and it defines a Hadamard subfactor of index $MN$. If $\Delta$ and $\Delta'$ are diagonal unitary matrices indexed by $X$ and $Y$, replacing $Q$ by $\Delta Q\Delta'$ multiplies each row of $H$ and each column of $H$ by a phase. This does not change the lines spanned by the coordinatewise quotients of the rows, hence it does not change the magic unitary of Section \ref{sec:prelim-cqg}, the quantum group $G_H$ and the rooted law. We may therefore assume that $Q$ is dephased, that is $Q_{0b}=Q_{i0}=1$ for any $i\in X$ and any $b\in Y$. A dephased matrix $Q$ is called generic if its $(M-1)(N-1)$ remaining entries are multiplicatively independent \cite[Definition 4.3]{BBi15}.

\begin{thm}[\cite{BBi15}]\label{thm:gram}
Suppose that $Q$ is generic. Let $Z$ be an $X\times Y$ random matrix with independent entries uniformly distributed on $\mathbb T$, and put $B=ZZ^*$. Then for any integer $n\geq1$ the $n$-th moment of the rooted law $\mu_H$ of the Hadamard subfactor is
\begin{equation}\label{eq:gram}
c_n=\frac1{MN}\mathbb E(Tr(B^n))
\end{equation}
Equivalently, $\mu_H=(1-\frac1N)\delta_0+\frac1N law(B)$, where $law(B)$ is the mean eigenvalue distribution of $B$ for the normalized trace $\frac1M Tr$.
\end{thm}
\begin{proof}
By Section \ref{sec:prelim-cqg}, the rooted law is the law of the main character of $G_H$ for the Haar state. For generic $Q$ this law is the measure of the statement by \cite[Theorem 4.6 and Proposition 5.2]{BBi15}. The moments of order $n\geq1$ of this measure are $\frac1N\cdot\frac1M\mathbb E(Tr(B^n))$, which is \eqref{eq:gram}.
\end{proof}

Theorem \ref{thm:gram} is also the case $N_Q=\{1\}$ of Theorem \ref{thm:gram-all} below. For 2 rows the matrix integral reduces to a planar random walk, see Proposition \ref{prop:two-rows} below.

\subsection{The crossed coproduct model}

We recall the factorization of \cite[Section 3]{BBi15}. The free wreath product $Y \wr_* X$ of \cite[Definition 3.4]{BBi15}, a variant of the construction of \cite{Bic04}, has the following function algebra. It is generated by a magic unitary $v=(v_{ij})_{i,j\in X}$ whose entries depend only on $i-j$, and by magic unitaries $u^{(i)}=(u^{(i)}_{ab})_{a,b\in Y}$, $i\in X$, whose entries depend only on $a-b$, with the relations
$$u^{(i)}_{ab}v_{ij}=v_{ij}u^{(i)}_{ab}$$
By \cite[Lemma 3.5]{BBi15}, the relations on $v$ and on each $u^{(i)}$ are those of the regular magic unitaries of $C(X)$ and $C(Y)$.

Let $Y^{*X}$ be the free product of copies $Y^{(i)}$, $i\in X$, of $Y$, and denote by $c^{(i)}$ the element $c\in Y$ of the $i$-th copy. The group $X$ acts on $Y^{*X}$ by $s\cdot c^{(i)}=c^{(i+s)}$. For an action of $X$ on a discrete group $\Gamma$, the crossed coproduct $C^*(\Gamma)\rtimes C(X)$ is the algebra $C^*(\Gamma)\otimes C(X)$ with the comultiplication \cite[Definition 3.7]{BBi15}
\begin{equation}\label{eq:dita-coproduct}
\Delta(g\otimes\delta_k)=\sum_{l\in X}
(g\otimes\delta_l)\otimes((-l)\cdot g\otimes\delta_{k-l}),
\quad g\in\Gamma,\ k\in X
\end{equation}
By \cite[Lemma 3.8]{BBi15}, the Fourier transform of each copy of $C(Y)$ gives an isomorphism $C(Y \wr_* X)\simeq C^*(Y^{*X})\rtimes C(X)$, under which $v_{ij}$ corresponds to $1\otimes\delta_{i-j}$ and $c^{(i)}\otimes1$ corresponds to $\sum_{a\in Y}(F_Y)_{ac}u^{(i)}_{a0}$.

On the Hilbert space $\mathbb C^{X\times Y}$ with standard basis $(e_j\otimes e_e)_{j\in X,e\in Y}$ we use the orthonormal basis
$$\xi_{ke}=M^{-\frac12}\sum_{j\in X}(F_X)_{jk}e_j\otimes e_e,
\quad k\in X,\ e\in Y$$
and the subspaces $E_k=span\{\xi_{ke}:e\in Y\}$.

\begin{prop}[\cite{BBi15}]\label{prop:dita-block}
Let $u$ be the magic unitary of $H=F_X\otimes_Q F_Y$. The representation $\pi$ of $C(S_{X\times Y}^+)$ defined by $u$ factors through $C(Y \wr_* X)\simeq C^*(Y^{*X})\rtimes C(X)$. The function $f\in C(X)$ acts on $E_k$ by the scalar $f(k)$, and
\begin{equation}\label{eq:dita-shift}
\pi(c^{(i)})\xi_{ke}
=\frac{Q_{i,e-c}Q_{i-k,e}}{Q_{i,e}Q_{i-k,e-c}}\xi_{k,e-c},
\quad i,k\in X,\ c,e\in Y
\end{equation}
In particular the restriction of $\pi$ to $C(X)$ is faithful.
\end{prop}
\begin{proof}
The factorization is \cite[Proposition 3.6]{BBi15}, composed with the isomorphism of \cite[Lemma 3.8]{BBi15}. The factorizing map is surjective, since $v_{ij}=\sum_au_{ia,jb}$ and $u^{(i)}_{ab}=\sum_ju_{ia,jb}$ are images of sums of entries of the fundamental magic unitary. Formula \eqref{eq:dita-shift} is \cite[Lemma 3.9]{BBi15}, where the basis vectors are not normalized. Finally, $v_{ij}=\sum_{a\in Y}u_{ia,jb}$ is the orthogonal projection onto $E_{i-j}$ by the formula for $W_{ia,jb}(\varepsilon_{ke})$ in the proof of \cite[Lemma 3.9]{BBi15}, so that $1\otimes\delta_k$ acts as the orthogonal projection onto $E_k$.
\end{proof}

By \eqref{eq:dita-shift}, a word $c_1^{(i_1)}\ldots c_n^{(i_n)}$ maps $\xi_{ke}$ to a multiple of $\xi_{k,e-d}$, where $d=c_1+\ldots+c_n$ is its degree. The degree defines a homomorphism $\deg:Y^{*X}\to Y$, and the elements of degree zero act diagonally in the basis $(\xi_{ke})$. Following \cite[Theorem 3.10]{BBi15}, let $\Gamma_{X,Y}$ be the quotient of $Y^{*X}$ by the commutators of elements of degree zero. This subgroup of commutators is stable under the action of $X$, so the representation $\pi$ factors through $C^*(\Gamma_{X,Y})\rtimes C(X)$ \cite[Theorem 3.10]{BBi15}.

We write $x^jy^d$ for the basis element $(j,d)$ of the group ring $\mathbb Z[X\times Y]$, and $I_X\subset\mathbb Z[X]$, $I_Y\subset\mathbb Z[Y]$ for the augmentation ideals. Thus $I_X$ is free abelian with basis $x^i-1$, $i\neq0$, and $I_Y$ is free abelian with basis $y^c-1$, $c\neq0$.

\begin{lemma}\label{lem:dita-lattice}
Let $T$ be the kernel of the degree map $\Gamma_{X,Y}\to Y$, and put $t_{ic}=(-c)^{(0)}c^{(i)}$ for $i\in X$ and $c\in Y$. The map $c\to c^{(0)}$ splits the degree map, so that $\Gamma_{X,Y}=T\rtimes Y$. The group $T$ is abelian, and there exists an isomorphism
\begin{equation}\label{eq:dita-lattice}
T\to I_X\otimes I_Y\subset\mathbb Z[X\times Y],
\quad t_{ic}\to(x^i-1)(y^c-1)
\end{equation}
under which conjugation by $(-d)^{(0)}$ is the multiplication by $y^d$ and the action of $s\in X$ is the multiplication by $x^s$. In particular the elements $t_{ic}$ with $i\neq0$ and $c\neq0$ form a basis of $T$.
\end{lemma}

\begin{proof}
The group $T$ is abelian by definition of $\Gamma_{X,Y}$, and $\deg(c^{(0)})=c$. Since $c^{(i)}=c^{(0)}t_{ic}$, the relation $c^{(i)}d^{(i)}=(c+d)^{(i)}$ becomes
\begin{equation}\label{eq:dita-conjugation}
(-d)^{(0)}t_{ic}d^{(0)}=t_{i,c+d}t_{id}^{-1}
\end{equation}
The group $\Gamma_{X,Y}$ is generated by the elements $c^{(0)}$ and $t_{ic}$. By \eqref{eq:dita-conjugation}, the conjugates of the elements $t_{ic}$ by the elements of $Y^{(0)}$ are products of elements $t_{ic}$. We deduce that $T$ is generated by the $t_{ic}$. Note that $t_{0c}=t_{i0}=1$.

Conversely, consider the semidirect product $(I_X\otimes I_Y)\rtimes Y$ in which conjugation by $-d$ is the multiplication by $y^d$, and put $w_{ic}=(x^i-1)(y^c-1)$. The identity
$$y^d w_{ic}+w_{id}=(x^i-1)(y^{c+d}-y^d)+(x^i-1)(y^d-1)=w_{i,c+d}$$
shows that $c\to cw_{ic}$ is a group homomorphism $Y\to(I_X\otimes I_Y)\rtimes Y$ for each $i\in X$. These homomorphisms define a homomorphism on $Y^{*X}$. Its restriction to the elements of degree zero takes values in the abelian group $I_X\otimes I_Y$, hence it factors through $\Gamma_{X,Y}$. It sends $t_{ic}$ to $w_{ic}$. Since the elements $w_{ic}$ with $i,c\neq0$ form a basis of $I_X\otimes I_Y$ and generate the image of $T$, while the $t_{ic}$ generate $T$, the restriction to $T$ is an isomorphism, and \eqref{eq:dita-conjugation} gives the description of the conjugation. Finally,
$$s\cdot t_{ic}=(-c)^{(s)}c^{(s+i)}=(-c)^{(0)}t_{s,-c}c^{(0)}t_{s+i,c}$$
which corresponds to $y^cw_{s,-c}+w_{s+i,c}=(x^s-1)(1-y^c)+(x^{s+i}-1)(y^c-1)=x^sw_{ic}$.
\end{proof}

The group $T$ is free abelian of rank $(M-1)(N-1)$, which is \cite[Proposition 4.2]{BBi15}. The module structure \eqref{eq:dita-lattice} is also used in \cite[Section 4.2]{BO26}. We identify $T$ with $I_X\otimes I_Y$ through \eqref{eq:dita-lattice} and write its group law additively when we compute in $\mathbb Z[X\times Y]$.

\subsection{The quotient determined by the parameter matrix}

An element of $\Gamma_{X,Y}$ can act by a different scalar on each subspace $E_k$. Since an element $f\in C(X)$ acts on $E_k$ by the scalar $f(k)$, such elements give candidate relations $r\otimes1=1\otimes f$ in the Hopf image. The following notion \cite[Definition 2.1]{Bic15} selects the relations which are compatible with the Hopf algebra structure.

\begin{defi}
Let $X$ act on a discrete group $\Gamma$. A quotient datum consists of an $X$-stable normal subgroup $D\subset\Gamma$ and a group homomorphism $\Phi:D\to C(X,\mathbb T)$, $r\to\Phi_r$, such that $\Phi_{grg^{-1}}=\Phi_r$ for any $g\in\Gamma$ and
$$\Phi_r(l+h)=\Phi_r(l)\Phi_{(-l)\cdot r}(h),
\quad l,h\in X,\ r\in D$$
We denote by $\mathbb C[\Gamma/D]\rtimes_\Phi C(X)$ the Hopf algebra $k[\Gamma/D]\rtimes_\Phi k^X$ of \cite[Proposition 2.4]{Bic15}, for $k=\mathbb C$ and $G=H=X$.
\end{defi}

The basic properties of this algebra, that we will need in what follows, are as follows:

\begin{prop}\label{prop:dita-quotient}
Let $(D,\Phi)$ be a quotient datum for $\Gamma$.
\begin{enumerate}
\item $\mathbb C[\Gamma/D]\rtimes_\Phi C(X)$ is the quotient of $\mathbb C[\Gamma]\rtimes C(X)$ by the ideal $J$ generated by the elements $r\otimes1-1\otimes\Phi_r$, with $r\in D$, and it has a basis indexed by $(\Gamma/D)\times X$.

\item $\mathbb C[\Gamma/D]\rtimes_\Phi C(X)$ is a compact Hopf $*$-algebra, with involution $(g\otimes f)^*=g^{-1}\otimes\bar f$.

\item $\mathbb C[\Gamma/D]\rtimes_\Phi C(X)$ embeds as a dense $*$-subalgebra of its enveloping $C^*$-algebra, denoted $C^*(\Gamma/D)\rtimes_\Phi C(X)$, which is the quotient of $C^*(\Gamma)\rtimes C(X)$ by the closed ideal generated by the elements $r\otimes1-1\otimes\Phi_r$, with $r\in D$.
\end{enumerate}
\end{prop}

\begin{proof}
(1) By \cite[Proposition 2.4]{Bic15}, $\mathbb C[\Gamma/D]\rtimes_\Phi C(X)$ has a basis indexed by $(\Gamma/D)\times X$, and there is a surjective Hopf algebra map $q$ from $\mathbb C[\Gamma]\rtimes C(X)$ onto it, with
$$q(r\otimes1)=q(1\otimes\Phi_r)$$
for any $r\in D$. Thus $J\subset\ker q$. Conversely, if $j$ is a section of $\Gamma\to\Gamma/D$, then
$$g\otimes\delta_h=(gj(\bar g)^{-1}\otimes1)(j(\bar g)\otimes\delta_h)$$
is congruent modulo $J$ to a multiple of $j(\bar g)\otimes\delta_h$. Thus the elements $j(c)\otimes\delta_h$, with $c\in\Gamma/D$ and $h\in X$, span the quotient by $J$, whereas their images under $q$ are linearly independent, and we conclude that we have $\ker q=J$.

(2) The involution $(g\otimes f)^*=g^{-1}\otimes\bar f$ of $\mathbb C[\Gamma]\rtimes C(X)$ preserves the relations $r\otimes1=1\otimes\Phi_r$, since $\bar\Phi_r=\Phi_{r^{-1}}$. Hence $\mathbb C[\Gamma/D]\rtimes_\Phi C(X)$ is a Hopf $*$-algebra, and $q$ is a map of Hopf $*$-algebras. The quotient is spanned by the images of coefficients of unitary representations, which remain unitary. It is therefore compact, giving the result.

(3) The Haar state of a compact Hopf algebra is faithful \cite[Theorems 3.7 and 4.4]{DK94}, so the GNS representation is injective, and this gives the embedding. As for the description of the enveloping $C^*$-algebra, this follows from (1).
\end{proof}

Since the antipode of $C(S^+_{X\times Y})$ satisfies $S^2=Id$ and $\Phi$ takes values in $\mathbb T$, the Hopf images below can be computed on the Hopf $*$-algebras $\mathbb C[\Gamma]\rtimes C(X)$ and $\mathbb C[\Gamma/D]\rtimes_\Phi C(X)$. The criterion that we will use is as follows:

\begin{prop}[\cite{Bic15}]\label{prop:hopf-image}
Let $\rho$ be a representation of $C^*(\Gamma)\rtimes C(X)$ whose restriction to $C(X)$ is faithful, and let $D$ be the set of the elements $r\in\Gamma$ such that, for some $f\in C(X)$,
$$\rho(r\otimes1)=\rho(1\otimes f)$$
Assume that $D$ is $X$-stable, and carries a homomorphism $\Phi$ with $\rho(r\otimes1)=\rho(1\otimes\Phi_r)$, which makes $(D,\Phi)$ a quotient datum. Then the Hopf image of $\rho$ is $C^*(\Gamma/D)\rtimes_\Phi C(X)$.
\end{prop}

\begin{proof}
This is \cite[Corollary 4.3]{Bic15}. Indeed, the subgroup $N$ there is defined by the condition $\rho(h\cdot r\otimes1)=\rho(1\otimes f)$ for every $h\in X$, with $f$ invertible, and it is equal to $D$. The condition for every $h$ reduces to $r\in D$, since $D$ is $X$-stable. As for $f$, this is automatically invertible, since $\rho(r\otimes1)$ is invertible, and $\rho(1\otimes\delta_k)\neq0$ for every $k\in X$, by faithfulness.
\end{proof}

We apply this criterion to the representation $\pi$ of $C^*(\Gamma_{X,Y})\rtimes C(X)$ of Proposition \ref{prop:dita-block}. The parameters are encoded by $\theta_{jd}\in\mathbb R/\mathbb Z$ with $e^{2\pi i\theta_{jd}}=Q_{-j,d}$. For $r=\sum_{j,d}r_{jd}x^jy^d$ in $\mathbb Z[X\times Y]$, define $F_r=r\theta$ by convolution,
\begin{equation}\label{eq:dita-convolution}
F_r(k,e)=\sum_{j\in X,d\in Y}r_{jd}\theta_{k-j,e-d}
\quad{\rm in}\ \mathbb R/\mathbb Z,
\quad k\in X,\ e\in Y
\end{equation}
Only integer multiples of the $\theta_{jd}$ occur, so this definition requires no choice of logarithm. For any abelian group $A$, let $S_A\subset A^{X\times Y}$ be the subgroup of the functions of the form $F(k,e)=\lambda(k)+\lambda'(e)$. A function $F$ belongs to $S_A$ if and only if
\begin{equation}\label{eq:dita-difference}
F(k,e)-F(0,e)-F(k,0)+F(0,0)=0,
\quad k\in X,\ e\in Y
\end{equation}
Indeed, if \eqref{eq:dita-difference} holds, one can take $\lambda(k)=F(k,0)-F(0,0)$ and $\lambda'(e)=F(0,e)$, and the converse is clear. If $X,Y$ are cyclic with generators corresponding to $x,y$, condition \eqref{eq:dita-difference} is equivalent to $(x-1)(y-1)F=0$, since summing the latter identity along the 2 cyclic factors gives \eqref{eq:dita-difference}. In terms of these functions, the relations that we are looking for are as follows:

\begin{defi}
The relation subgroup of a dephased matrix $Q\in\mathbb T^{X\times Y}$ is
$$N_Q=\{r\in I_X\otimes I_Y:F_r\in S_{\mathbb R/\mathbb Z}\}$$
and for $r\in N_Q$ we define a function $\Phi_r:X\to\mathbb T$ by $\Phi_r(k)=e^{2\pi i(F_r(0,0)-F_r(k,0))}$.
\end{defi}

With these notions in hand, the quantum group can be computed, as follows:

\begin{thm}[\cite{BO26}]\label{thm:dita-all}
For a dephased $Q\in\mathbb T^{X\times Y}$ and $H=F_X\otimes_QF_Y$, we have
$$C(G_H)=C^*(\Gamma_{X,Y}/N_Q)\rtimes_\Phi C(X)$$
where $(N_Q,\Phi)$ is a quotient datum for $\Gamma_{X,Y}$.
\end{thm}

\begin{proof}
This is \cite[Theorem 4.4]{BO26}. We include here the proof, with all the details on the $C^*$-algebra side, the subsequent blowup construction being based on them.

We first compute the action of $T$. Let $\varphi_{jd}=\theta_{-j,d}$, so that $Q_{jd}=e^{2\pi i\varphi_{jd}}$ and $\varphi_{0d}=0$. Applying \eqref{eq:dita-shift} first to $c^{(i)}$ and then to $(-c)^{(0)}$, we see that $t_{ic}$ acts on $\xi_{ke}$ by the multiplication by $e^{2\pi i V}$, where
$$V=\varphi_{i,e-c}-\varphi_{i,e}+\varphi_{i-k,e}-\varphi_{i-k,e-c}+\varphi_{-k,e-c}-\varphi_{-k,e}$$
On the other hand, for $r=(x^i-1)(y^c-1)=x^iy^c-x^i-y^c+1$, the definition \eqref{eq:dita-convolution} gives
$$F_r(k,e)=\varphi_{i-k,e-c}-\varphi_{i-k,e}-\varphi_{-k,e-c}+\varphi_{-k,e}
\quad{\rm and}\quad
F_r(0,e)=\varphi_{i,e-c}-\varphi_{i,e}$$
Thus $V=F_r(0,e)-F_r(k,e)$. Both sides are additive in $r$, and the group law of $T$ corresponds to the product of the diagonal operators. By Lemma \ref{lem:dita-lattice}, every $r\in T$ therefore satisfies
\begin{equation}\label{eq:dita-diagonal}
\pi(r)\xi_{ke}=e^{2\pi i(F_r(0,e)-F_r(k,e))}\xi_{ke},
\quad k\in X,\ e\in Y
\end{equation}

Now, we identify the set $D$ of Proposition \ref{prop:hopf-image}. By Proposition \ref{prop:dita-block}, the element $1\otimes f$ acts on $E_k$ by $f(k)$. An element of nonzero degree $d$ maps $\xi_{ke}$ to a multiple of $\xi_{k,e-d}\neq\xi_{ke}$, so it does not belong to $D$. By \eqref{eq:dita-diagonal}, an element $r\in T$ acts by a scalar on every $E_k$ if and only if $F_r(0,e)-F_r(k,e)$ does not depend on $e$ for any $k$, that is, if and only if \eqref{eq:dita-difference} holds for $F_r$. Hence $D=N_Q$, and putting $e=0$ shows that $r\in N_Q$ acts on $E_k$ by $\Phi_r(k)$.

We check that $(N_Q,\Phi)$ is a quotient datum. For $s\in X$ and $d\in Y$, the definition \eqref{eq:dita-convolution} gives $F_{x^sy^dr}(k,e)=F_r(k-s,e-d)$, and condition \eqref{eq:dita-difference} is invariant under these translations. By Lemma \ref{lem:dita-lattice}, the multiplications by $x^s$ and by $y^d$ implement the action of $X$ and the conjugation by the elements of $Y^{(0)}$, while the conjugation by the elements of the abelian group $T$ is trivial. Hence $N_Q$ is normal and $X$-stable.

The map $r\to\Phi_r$ is a homomorphism since $r\to F_r$ is additive. Every element of $\Gamma_{X,Y}$ maps each subspace $E_k$ onto itself, so that $grg^{-1}$ acts on $E_k$ by the same scalar as $r$, and $\Phi_{grg^{-1}}=\Phi_r$. Finally, since $(-l)\cdot r=x^{-l}r$ and $F_{x^{-l}r}(h,0)=F_r(h+l,0)$, we have
$$\Phi_r(l)\Phi_{(-l)\cdot r}(h)
=e^{2\pi i(F_r(0,0)-F_r(l,0))}e^{2\pi i(F_r(l,0)-F_r(l+h,0))}
=\Phi_r(l+h)$$
Thus Proposition \ref{prop:hopf-image} applies, since the restriction of $\pi$ to $C(X)$ is faithful by Proposition \ref{prop:dita-block}, and gives the result.
\end{proof}

Regarding now the computation of $N_Q$, write $Q_{jd}=\eta_{jd}\prod_{s=1}^m\gamma_s^{\beta_s(j,d)}$, where the $\eta_{jd}$ are roots of unity, $\beta_s(j,d)\in\mathbb Z$, and the images of $\gamma_1,\ldots,\gamma_m$ form a basis of the quotient of the subgroup of $\mathbb T$ generated by the entries of $Q$ by its torsion subgroup. Such a decomposition exists since this subgroup is finitely generated, so that the quotient is free abelian. With this convention, we have:

\begin{prop}[\cite{BO26}]\label{prop:dita-criterion}
Let $Q\in\mathbb T^{X\times Y}$ be dephased and let $H=F_X\otimes_QF_Y$.
\begin{enumerate}
\item We have $N_Q=\{1\}$ if and only if, for every pair of nontrivial characters $\chi$ of $X$ and $\psi$ of $Y$, there exists $s$ such that
$$\sum_{j\in X,d\in Y}\beta_s(j,d)\bar\chi(j)\psi(d)
\neq0$$
\item If $Q_{jd}=e^{2\pi ia_{jd}/L}$ with $a_{jd}\in\mathbb Z_L$ for any $j,d$, then $G_H$ is finite of order $MN[T:N_Q]$, and the index $[T:N_Q]$ is the order of the image of the homomorphism
$$T\to\mathbb Z_L^{X\times Y}/S_{\mathbb Z_L},
\quad r\to r\widehat\theta\bmod S_{\mathbb Z_L}$$
where $\widehat\theta_{jd}=a_{-j,d}$ and the product is the convolution.
\end{enumerate}
\end{prop}

\begin{proof}
This is part of \cite[Theorem 4.4]{BO26}, and we refer to that paper for the proof, which is based on the above description of $N_Q$.
\end{proof}

Finally, here is the formula for the moments, which is our main tool in what follows:

\begin{thm}\label{thm:dita-words}
Let $Q\in\mathbb T^{X\times Y}$ be dephased and let $H=F_X\otimes_QF_Y$. For any $n\geq1$, the $n$-th moment of $\mu_H$ is
\begin{equation}\label{eq:dita-words}
c_n=\frac1M\#\Big\{((i_1,b_1),\ldots,(i_n,b_n))\in(X\times Y)^n:
b_1^{(i_1)}\ldots b_n^{(i_n)}\in N_Q\Big\}
\end{equation}
In particular, for fixed $X$ and $Y$ the rooted law depends only on $N_Q$.
\end{thm}

\begin{proof}
In the $i$-th copy of $C(Y)$, the diagonal entry $u^{(i)}_{aa}$ is the indicator function of the identity of $Y$, which corresponds under the Fourier transform to $N^{-1}\sum_{b\in Y}b^{(i)}$. Since $v_{ii}$ corresponds to $1\otimes\delta_0$, summing the $MN$ diagonal entries $u^{(i)}_{aa}v_{ii}$ of the fundamental magic unitary gives the main character
$$\chi_H=\Big(\sum_{i\in X,b\in Y}b^{(i)}\Big)\otimes\delta_0$$
Let $h$ be the Haar state of $C(G_H)$. Compose evaluation at $0\in X$ with the group quotient to obtain the Hopf algebra map $g\otimes f\to f(0)\bar g$, where $\bar g$ denotes the class of $g$ in $\Gamma_{X,Y}/N_Q$. For $r\in N_Q$ we have $\bar r=1$ and $\Phi_r(0)=1$, so this map respects the relations of Proposition \ref{prop:dita-quotient}. Denote the induced map from $C(G_H)$ to $C^*(\Gamma_{X,Y}/N_Q)$ by $ev_0$. Applying $ev_0\otimes Id$ to the coproduct of $g\otimes\delta_0$ in \eqref{eq:dita-coproduct} gives $\bar g\otimes(g\otimes\delta_0)$, and the invariance $(ev_0\otimes h)\Delta=h(\cdot)1$ then gives
$$h(g\otimes\delta_0)\bar g=h(g\otimes\delta_0)1$$
so that $h(g\otimes\delta_0)=0$ if $\bar g\neq1$. If $\bar g=1$, then $g\in N_Q$ and $g\otimes\delta_0=(1\otimes\Phi_g)(1\otimes\delta_0)=1\otimes\delta_0$ in $C(G_H)$.

The restriction of $h$ to the image of $C(X)$, which is isomorphic to $C(X)$ by the faithfulness of $\pi$ on $C(X)$, is the Haar state of $C(X)$. Hence $h(1\otimes\delta_0)=\frac1M$, and by expanding $h(\chi_H^n)$ we obtain the formula in the statement.
\end{proof}

\section{Gram blowup, rigidity and transposition}\label{sec:fourier-blowups}

We keep the notation of Section \ref{sec:hadamard}. In the spirit of \cite{BBi15}, we call blowup of the rooted law a formula which expresses it through the law of a positive matrix valued function on a compact group, as in Theorem \ref{thm:gram}. We prove such a formula for every parameter matrix in Theorem \ref{thm:gram-all}, and we compute the compact group explicitly from the entries of $Q$ in Proposition \ref{prop:gram-generators}. The examples at indices 6, 8 and 9, discussed in Sections \ref{sec:lattices} and \ref{sec:growth}, are then direct computations of this group.

For $z\in\mathbb T$ we denote by $\langle z\rangle$ the closed subgroup of $\mathbb T$ generated by $z$, which is finite cyclic if $z$ is a root of unity, and equal to $\mathbb T$ otherwise. Also, we use the letter $\varepsilon$ for a random sign, with $\mathbb P(\varepsilon=1)=\mathbb P(\varepsilon=-1)=\frac12$, independent of the other random variables. For $M=2$ or $N=2$, the subgroup $N_Q$ is also computed in \cite[Theorems 4.5 and 4.6]{BO26}, in agreement with the results below.

\subsection{The Gram model for arbitrary parameters}

Let $Q$ be dephased. Put $\Gamma_Q=\Gamma_{X,Y}/N_Q$ and $T_Q=T/N_Q$. For $i\in X$ and $j\in Y$, we have $t_{i,-j}=j^{(0)}(-j)^{(i)}$, and the elements $t_{i,-j}$ with $i,j\neq0$ form a basis of $T$ by Lemma \ref{lem:dita-lattice}. Let $\widehat T_Q$ be the Pontryagin dual of the discrete abelian group $T_Q$, with its Haar probability measure. The map
$$\widehat T_Q\to\mathbb T^{X\times Y},
\quad \chi\to Z(\chi)=(\chi(t_{i,-j}))_{i\in X,j\in Y}$$
where $\chi(t_{i,-j})$ means the value of $\chi$ on the image of $t_{i,-j}$ in $T_Q$, is an injective continuous homomorphism, since these images generate $T_Q$. A character of $T$ is determined by its values on the basis $(t_{i,-j})_{i,j\neq0}$, which can be chosen arbitrarily in $\mathbb T$, and it factors through $T_Q$ if and only if it is trivial on $N_Q$. Hence the image of this map is the closed subgroup
\begin{equation}\label{eq:gram-group}
K_Q=\Big\{Z\in\mathbb T^{X\times Y}:Z_{0j}=Z_{i0}=1\ {\rm and}\ \prod_{i,j\neq0}Z_{ij}^{m_{ij}}=1\ {\rm whenever}\ \sum_{i,j\neq0}m_{ij}t_{i,-j}\in N_Q\Big\}
\end{equation}
and it carries the image of the Haar measure of $\widehat T_Q$, which is its Haar measure. Deleting the first row and the first column identifies $K_Q$ with a closed subgroup of $\mathbb T^{(M-1)(N-1)}$.

\begin{thm}[The Gram blowup for all parameters]\label{thm:gram-all}
Let $Q\in\mathbb T^{X\times Y}$ be dephased and $H=F_X\otimes_Q F_Y$. Then for any integer $n\geq1$ we have
\begin{equation}\label{eq:gram-all}
c_n=\frac1{MN}\int_{K_Q}Tr((ZZ^*)^n)dZ
\end{equation}
Equivalently, the rooted law is
$$\mu_H=\Big(1-\frac1N\Big)\delta_0+\frac1N law(ZZ^*)$$
where $Z$ is Haar distributed on $K_Q$ and $law(ZZ^*)$ is the mean eigenvalue distribution of $ZZ^*$ for the normalized trace $\frac1M Tr$.
\end{thm}

\begin{proof}
Let $Y_i$ be the image of the $i$-th copy of $Y$ in $\Gamma_Q$. Since $N_Q\subset T$, the degree map induces a homomorphism $\Gamma_Q\to Y$ with kernel $T_Q$, which is an isomorphism on each $Y_i$. Consequently, each left coset of $Y_i$ contains exactly one element of $T_Q$. We keep the coset spaces for the different $i$ as separate copies, also when some of the subgroups $Y_i$ coincide. Consider the operator
$$J:l^2(\Gamma_Q)\to\bigoplus_{i\in X}l^2(\Gamma_Q/Y_i),
\quad J\delta_\gamma=\bigoplus_{i\in X}\delta_{\gamma Y_i}$$
The adjoint $J^*$ sums over the elements of a coset. Hence $J^*J\delta_\gamma=\sum_{i\in X,b\in Y}\delta_{\gamma b^{(i)}}$, where the letters $b^{(i)}$ are those of \eqref{eq:dita-words}. The coefficient $\langle(J^*J)^n\delta_1,\delta_1\rangle$ is the number of words of length $n$ in these letters whose product is equal to $1$ in $\Gamma_Q$, that is, whose product in $\Gamma_{X,Y}$ belongs to $N_Q$. By \eqref{eq:dita-words}, this number is $Mc_n$. Since $J^*J$ commutes with the left translations of $\Gamma_Q$, every diagonal coefficient of $(J^*J)^n$ is equal to $Mc_n$.

Since $\Gamma_Q=T_Q\rtimes Y$, every $\gamma\in\Gamma_Q$ has a unique expression $\gamma=uj^{(0)}$ with $u\in T_Q$ and $j\in Y$. This gives an identification $l^2(\Gamma_Q)=l^2(T_Q)\otimes l^2(Y)$, and we identify $l^2(\Gamma_Q/Y_i)$ with $l^2(T_Q)$ by means of the representatives of degree zero. The coset of $uj^{(0)}$ in $\Gamma_Q/Y_i$ has the representative of degree zero $uj^{(0)}(-j)^{(i)}=ut_{i,-j}$. Thus the $(i,j)$ entry of $J$ is the translation by $t_{i,-j}$ on $l^2(T_Q)$.

The Fourier transform $l^2(T_Q)\to L^2(\widehat T_Q)$, $\delta_u\to(\chi\to\chi(u))$, is unitary and transforms the translation by $t$ into the multiplication by the function $\chi\to\chi(t)$. Hence it transforms $J$ into the multiplication by the matrix valued function $\chi\to Z(\chi)$, and $\delta_{j^{(0)}}=\delta_1\otimes\delta_j$ into the constant function $1$ times $\delta_j$. We obtain
$$Mc_n=\langle(J^*J)^n\delta_{j^{(0)}},\delta_{j^{(0)}}\rangle
=\int_{\widehat T_Q}((Z(\chi)^*Z(\chi))^n)_{jj}d\chi,
\quad j\in Y$$
Summing over $j\in Y$ gives
$$MNc_n=\int_{\widehat T_Q}Tr((Z^*Z)^n)d\chi
=\int_{K_Q}Tr((ZZ^*)^n)dZ$$
since $Z^*Z$ and $ZZ^*$ have the same nonzero eigenvalues with multiplicities. This proves \eqref{eq:gram-all}. Finally, $ZZ^*\geq0$ and $Tr(ZZ^*)=MN$, so that the measure of the statement is a probability measure supported in $[0,MN]$. Its moments of order $n\geq1$ are the right-hand side of \eqref{eq:gram-all}. Recall that 2 probability measures with compact supports and the same moments coincide.
\end{proof}

The next result computes $K_Q$ directly from the entries of $Q$.

\begin{prop}[Generators from the parameter matrix]\label{prop:gram-generators}
For $i,k\in X$ and $e,j\in Y$, put
\begin{equation}\label{eq:gram-generators}
R_k(i,e)=\frac{Q_{i,e}Q_{-k,e}}{Q_{i-k,e}Q_{0,e}},\quad
g_{k,e}(i,j)=\frac{R_k(i,e+j)R_k(i,0)}{R_k(i,e)R_k(i,j)}
\end{equation}
Then $K_Q$ is the closed subgroup of $\mathbb T^{X\times Y}$ generated by the matrices $g_{k,e}$, $k\in X$, $e\in Y$. The same formula applies to a parameter matrix which is not dephased.
\end{prop}

\begin{proof}
If $Q$ is replaced by $\Delta Q\Delta'$ with diagonal unitary matrices $\Delta,\Delta'$, each $R_k(i,e)$ is multiplied by a factor which does not depend on $e$, and this factor cancels in $g_{k,e}$. We may therefore assume that $Q$ is dephased. Applying \eqref{eq:dita-shift} first to $(-j)^{(i)}$ and then to $j^{(0)}$, we see that $t_{i,-j}=j^{(0)}(-j)^{(i)}$ maps $\xi_{ke}$ to
$$\frac{Q_{i,e+j}Q_{i-k,e}}{Q_{i,e}Q_{i-k,e+j}}\cdot\frac{Q_{0,e}Q_{-k,e+j}}{Q_{0,e+j}Q_{-k,e}}\xi_{ke}
=\frac{R_k(i,e+j)}{R_k(i,e)}\xi_{ke}$$
Let $(m_{ij})$ be an integer matrix. By the proof of Theorem \ref{thm:dita-all}, the element $\sum_{i,j}m_{ij}t_{i,-j}$ of $T$ belongs to $N_Q$ if and only if it acts by a scalar on each subspace $E_k$. Dividing by the value at $e=0$, this means that
$$\prod_{i\in X,j\in Y}g_{k,e}(i,j)^{m_{ij}}=1,
\quad k\in X,\ e\in Y$$
Note that $g_{k,e}(i,j)=1$ if $i=0$ or $j=0$. By \eqref{eq:gram-group}, $K_Q$ is the subgroup of the matrices $Z$ with first row and first column equal to one which are annihilated by the characters $Z\to\prod Z_{ij}^{m_{ij}}$ with $\sum m_{ij}t_{i,-j}\in N_Q$. We have just shown that these characters are exactly the characters which are trivial on the matrices $g_{k,e}$. By Pontryagin duality, a closed subgroup of a compact abelian group is the annihilator of its annihilator. Hence the closed subgroup generated by the $g_{k,e}$ is $K_Q$.
\end{proof}

In terms of the second differences $W_{ij}=Q_{i+k,j+e}Q_{ij}/(Q_{i+k,j}Q_{i,j+e})$ of the matrix $Q$, the above formula shows that $g_{-k,e}^{-1}$ is the dephased form $(W_{ij}W_{00}/W_{i0}W_{0j})_{ij}$ of the matrix $W$, a fact which is checked by the program \texttt{pantani.py} too. This gives the description of $K_Q$ in Theorem A.

If $N_Q=\{1\}$, then $K_Q$ is the subgroup of all matrices with first row and first column equal to one, and Theorem \ref{thm:gram-all} gives back Theorem \ref{thm:gram}, since the eigenvalues of $ZZ^*$ do not change when the rows and columns of $Z$ are multiplied by phases. If $Q=1$, then $K_Q=\{1\}$, the matrix $Z$ is the all-ones matrix and $\mu_H=(1-\frac1{MN})\delta_0+\frac1{MN}\delta_{MN}$. If the entries of $Q$ are roots of unity, then $K_Q$ is finite and the integral in \eqref{eq:gram-all} is a uniform average over $K_Q$.

\subsection{Degenerations}

We denote by $c_n^{gen}$ the moments of $\mu_H$ for a generic parameter matrix. By Theorem \ref{thm:gram} they depend only on $M$ and $N$. The next result shows that the rooted law detects whether $G_H$ is a proper quantum subgroup of the generic quantum group, that is, whether $C(G_H)$ is a proper quotient of $C^*(\Gamma_{X,Y})\rtimes C(X)$.

\begin{thm}\label{thm:rigidity}
Let $Q\in\mathbb T^{X\times Y}$ be dephased and let $H=F_X\otimes_QF_Y$. For any integer $n\geq1$ we have
\begin{equation}\label{eq:rigidity}
c_n-c_n^{gen}=\frac1M\#\Big\{((i_1,b_1),\ldots,(i_n,b_n))\in(X\times Y)^n:b_1^{(i_1)}\ldots b_n^{(i_n)}\in N_Q-\{1\}\Big\}\geq0
\end{equation}
If $N_Q\neq\{1\}$, the smallest integer $n$ such that $c_n>c_n^{gen}$ is the smallest length of a nontrivial element of $N_Q$ with respect to the generators $b^{(i)}$, $i\in X$, $b\in Y-\{0\}$, of $\Gamma_{X,Y}$.
\end{thm}

\begin{proof}
A dephased matrix whose $(M-1)(N-1)$ free entries are multiplicatively independent exists, and for such a matrix $N_Q=\{1\}$ by Proposition \ref{prop:dita-criterion} (1), applied with the $\gamma_s$ equal to these entries. Applying \eqref{eq:dita-words} to this matrix and to $Q$, and subtracting, gives \eqref{eq:rigidity}.

For $\gamma\in\Gamma_{X,Y}$, let $W_n(\gamma)$ be the number of words of length $n$ in the letters $b^{(i)}$, $(i,b)\in X\times Y$, with product $\gamma$. Since the letter $0^{(i)}$ is the identity, $W_n(\gamma)>0$ if and only if the length $|\gamma|$ of $\gamma$ with respect to the generators $b^{(i)}$, $b\neq0$, satisfies $|\gamma|\leq n$. By \eqref{eq:rigidity}, we have
$$c_n-c_n^{gen}=\frac1M\sum_{\gamma\in N_Q-\{1\}}W_n(\gamma)$$
which is positive if and only if $N_Q-\{1\}$ contains an element of length at most $n$. This proves the last assertion.
\end{proof}

As a consequence, the rooted law detects the degenerations of the quantum group:

\begin{thm}\label{thm:degenerate}
Let $Q\in\mathbb T^{X\times Y}$ be dephased and let $H=F_X\otimes_QF_Y$. The following are equivalent:
\begin{enumerate}
\item $N_Q=\{1\}$.
\item The quotient map $C^*(\Gamma_{X,Y})\rtimes C(X)\to C(G_H)$ is an isomorphism.
\item The rooted law $\mu_H$ is the generic law of Theorem \ref{thm:gram}.
\item The box spaces of the Hadamard subfactor of $H$ satisfy $\dim P_n=c_n^{gen}$ for any $n\geq0$.
\end{enumerate}
\end{thm}

\begin{proof}
The equivalence of (1) and (3) follows from Theorem \ref{thm:rigidity}, since a probability measure with compact support is determined by its moments. Also, (3) and (4) are equivalent since $\dim P_n=c_n$, by Proposition \ref{prop:hadamard-law}.

If $N_Q\neq\{1\}$ and $r\in N_Q-\{1\}$, the element $r\otimes1-1\otimes\Phi_r$ of the dense subalgebra $\mathbb C[\Gamma_{X,Y}]\rtimes C(X)$ of $C^*(\Gamma_{X,Y})\rtimes C(X)$ is nonzero, since $r\neq1$. By Proposition \ref{prop:dita-quotient}, it lies in the kernel of the quotient map onto $C^*(\Gamma_{X,Y}/N_Q)\rtimes_\Phi C(X)$, which is $C(G_H)$ by Theorem \ref{thm:dita-all}. Conversely, if $N_Q=\{1\}$, the quotient map is the identity. This proves the equivalence of (1) and (2).
\end{proof}

For $X=\mathbb Z_2$ and $Y$ of prime order, \cite[Theorem 5.1]{Bic15} shows that the representation of $C^*(\Gamma_{X,Y})\rtimes C(X)$ is inner faithful as soon as one entry of $Q$ is not a root of unity, that is $N_Q=\{1\}$, by Theorem \ref{thm:degenerate}. Thus, in this case, the rooted law is either the generic one, or that of a finite quantum group. This fails for $Y=\mathbb Z_4$, where intermediate cases appear, with phase groups which are infinite, but smaller than the generic one, see Proposition \ref{prop:eight-cyclic}. In Example \ref{ex:six-twelfth} below the first difference occurs at $n=8$.

\begin{remark}
The results of this section can be checked on examples, as follows. For a root-of-unity matrix $Q_{ij}=e^{2\pi i a_{ij}/L}$, the exact verification uses 2 descriptions. The block action of Proposition \ref{prop:dita-block} counts words modulo blockwise scalar phases. The generators of Proposition \ref{prop:gram-generators} give an integer matrix of phase exponents modulo $L$. If $s_1,\ldots,s_r$ are its nonzero Smith factors, its image has order
$$|K_Q|=[T:N_Q]=\prod_{a=1}^r\frac L{\gcd(L,s_a)}$$
The equality follows by applying the invertible integer row and column operations to the map modulo $L$. Expanding a Gram trace then amounts to retaining the monomials whose character is trivial on all generators. The program \texttt{wordgram.py} compares these counts for the following inputs and their transposes, in agreement with Theorem \ref{thm:transpose-all} below. Here $A=(a_{ij})_{i,j\neq0}$ is the free part of the exponent matrix, and $r$ is the highest checked moment. The first 4 positive moments are shown.
$$\begin{array}{|c|c|c|c|c|c|c|}\hline
X&Y&L&A&|K_Q|&r&(c_1,c_2,c_3,c_4)\\\hline
\mathbb Z_2&\mathbb Z_2&4&(1)&1&5&(1,4,16,64)\\\hline
\mathbb Z_2&\mathbb Z_2&8&(1)&2&5&(1,3,10,36)\\\hline
\mathbb Z_2&\mathbb Z_3&12&(0\ \ 1)&12&8&(1,4,18,86)\\\hline
\mathbb Z_2&\mathbb Z_3&7&(1\ \ 3)&7&5&(1,4,18,86)\\\hline
\mathbb Z_2&\mathbb Z_2^2&3&(1\ \ 2\ \ 2)&9&4&(1,5,28,169)\\\hline
\mathbb Z_3&\mathbb Z_3&6&(^1_0{\ }^{\ 0}_{-1})&4&5&(1,5,31,225)\\\hline
\mathbb Z_2^2&\mathbb Z_2&4&(1\ \ 2\ \ 3)^t&1&4&(1,8,64,512)\\\hline
\end{array}$$
For $\mathbb Z_2^2$ the elements are ordered as $00,01,10,11$. The same program checks the 3 rational atomic laws at index 6, the first 7 moments of the infinite-order index-9 slice, and the first exceptional moment of Proposition \ref{prop:eight-first-moment} for $m=1,2,3,4$.
\end{remark}

\subsection{Transposition}

The matrix $H^t$ is obtained from $F_Y\otimes_{Q^t}F_X$ by the same permutation $(i,a)\to(a,i)$ of the rows and of the columns, since $(H^t)_{ia,jb}=H_{jb,ia}=(Q^t)_{aj}(F_Y)_{ab}(F_X)_{ij}$. By \cite[Theorem 2.3]{BO26}, for any complex Hadamard matrix $H$, the matrices $H$ and $H^t$ have the same rooted law. In the present setting this equality can be seen on the phase groups, which are transposes of each other. On the other hand the quantum groups themselves are different, as observed in \cite[Section 4.3]{BO26} for $X=\mathbb Z_2$ and $Y=\mathbb Z_4$, where the numbers of characters are $32$ and $64$. Here is the general statement, for arbitrary finite abelian groups $X,Y$.

\begin{thm}\label{thm:transpose-all}
Let $Q\in\mathbb T^{X\times Y}$ be dephased.
\begin{enumerate}
\item We have $K_{Q^t}=\{Z^t:Z\in K_Q\}$. In particular the Gram formula of Theorem \ref{thm:gram-all} gives the same rooted law for $F_X\otimes_QF_Y$ and for $F_Y\otimes_{Q^t}F_X$.
\item If $Q$ is generic, the set of characters of the $*$-algebra $C(G_H)$ of $H=F_X\otimes_QF_Y$ has $MN^M$ elements. In particular, if $M\neq N$, the quantum groups of $F_X\otimes_QF_Y$ and of $F_Y\otimes_{Q^t}F_X$ are not isomorphic.
\end{enumerate}
\end{thm}

\begin{proof}
1. For $a,b\in X$ and $c,d\in Y$ consider the multiplicative double difference
$$\partial Q(a,b;c,d)=\frac{Q_{a+b,c+d}Q_{a,c}Q_{b,c}Q_{a,d}Q_{b,d}}{Q_{a,c+d}Q_{b,c+d}Q_{a+b,c}Q_{a+b,d}}$$
This expression is symmetric in $(a,b)$ and in $(c,d)$, and satisfies $\partial Q^t(c,d;a,b)=\partial Q(a,b;c,d)$. Since $Q$ is dephased, we have $R_k(i,0)=1$ and $R_k(i,e)=Q_{i,e}Q_{-k,e}Q_{i-k,e}^{-1}$ in \eqref{eq:gram-generators}, and a direct expansion, which is also checked by the program \texttt{pantani.py}, gives
$$g_{k,e}(i,j)=\partial Q(i,-k;e,j)^{-1},
\quad i,k\in X,\ e,j\in Y$$
Let $g'_{e,k}$, $e\in Y$, $k\in X$, be the matrices \eqref{eq:gram-generators} for the parameter matrix $Q^t\in\mathbb T^{Y\times X}$. For $j\in Y$ and $i\in X$ we obtain
$$g'_{e,k}(j,i)
=\partial Q^t(j,-e;k,i)^{-1}
=\partial Q(k,i;j,-e)^{-1}
=\partial Q(i,k;-e,j)^{-1}
=g_{-k,-e}(i,j)$$
Thus the generators of Proposition \ref{prop:gram-generators} for $Q^t$ are the transposes of the generators for $Q$, and $K_{Q^t}=\{Z^t:Z\in K_Q\}$. By Theorem \ref{thm:gram-all}, the moments of order $n\geq1$ of the rooted law of $F_Y\otimes_{Q^t}F_X$ are
$$\frac1{MN}\int_{K_Q}Tr((Z^t\bar Z)^n)dZ
=\frac1{MN}\int_{K_Q}Tr(((Z^*Z)^n)^t)dZ
=\frac1{MN}\int_{K_Q}Tr((ZZ^*)^n)dZ$$
which are the moments of $\mu_H$. Both laws are supported in $[0,MN]$, hence they coincide.

2. If $Q$ is generic, then $N_Q=\{1\}$ by Proposition \ref{prop:dita-criterion}, and $C(G_H)=C^*(\Gamma_{X,Y})\rtimes C(X)$ contains the Hopf $*$-algebra $\mathbb C[\Gamma_{X,Y}]\rtimes C(X)$ as a dense subalgebra, the product being that of the tensor product. A character of a tensor product of unital algebras is determined by its restrictions to the 2 factors, and any pair of characters of the factors defines a character, since the factors commute.

The $*$-characters of $C(X)$ are the $M$ evaluations at the points of $X$. As for the $*$-characters of $\mathbb C[\Gamma_{X,Y}]$, these are the homomorphisms $\Gamma_{X,Y}\to\mathbb T$, that is, the characters of the abelianization of $\Gamma_{X,Y}$. Now $\Gamma_{X,Y}$ is the quotient of the free product $Y^{*X}$ by a subgroup of its commutator subgroup, so its abelianization is that of $Y^{*X}$, namely the group $Y^X$, of order $N^M$. Hence we have $MN^M$ algebraic $*$-characters. Each of them extends to $C(G_H)$, by the universal property of the enveloping $C^*$-algebra, from Proposition \ref{prop:dita-quotient}, and density gives uniqueness.

For $F_Y\otimes_{Q^t}F_X$ we obtain $NM^N$ characters. The equality $MN^M=NM^N$ means
$$\frac{\log N}{N-1}=\frac{\log M}{M-1}$$
and the function $t\to\frac{\log t}{t-1}$ is strictly decreasing on $(1,\infty)$, since $\log t>1-\frac1t$ there. Hence $M=N$. A compact quantum group isomorphism induces a bijection between the characters of the $C^*$-algebras, which proves the last assertion.
\end{proof}

The above characters, which are the points of the maximal classical subgroup of $G_H$, should not be confused with the one-dimensional representations of $G_H$, which correspond to the group-like elements of $C(G_H)$, such as the elements $1\otimes\chi$, with $\chi$ being a character of $X$. Summarizing, the rooted law detects the degenerations of the quantum group, by Theorem \ref{thm:degenerate}, but does not distinguish $H$ from $H^t$.

\section{Coverings, honeycomb lattice, index 6}\label{sec:lattices}

We explain here the geometric meaning of the Gram formula, in terms of the abelian coverings of the complete bipartite graph, and then we discuss the case of the deformations with 2 rows, where the lattices which appear are well-known ones. At index 6 this leads to a complete description of the quantum groups, of their coset graphs, and of the rooted laws.

\subsection{The covering graph}

Let $Q\in\mathbb T^{X\times Y}$ be dephased, and consider the operator $J$ from the proof of Theorem \ref{thm:gram-all}. This is the incidence operator of the following bipartite graph $B_Q$. The white vertices are the elements of $\Gamma_Q$. The black vertices are the cosets $\gamma Y_i$, with $\gamma\in\Gamma_Q$ and $i\in X$, the cosets coming from different indices $i$ being kept as separate vertices, and we say that $\gamma Y_i$ is of type $i$. The edges join each $\gamma$ to the $M$ cosets $\gamma Y_i$. Thus each white vertex has valence $M$, and each black vertex has valence $N$. Observe that the graph $B_Q$ is finite precisely when the group $\Gamma_Q$ is finite, which is the case when all the entries of the parameter matrix $Q$ are roots of unity.

Let also $K_{X,Y}$ be the complete bipartite graph having $Y$ as set of white vertices, and $X$ as set of black vertices. We orient its edges from white to black, so that its first homology group is the group of integer matrices $m\in\mathbb Z^{X\times Y}$ whose sums on the rows and on the columns vanish. Recall that the first homology group of a connected graph is free abelian, of rank equal to the number of edges which are not in a spanning tree. Here this rank is $MN-(M+N-1)=(M-1)(N-1)$, which is the rank of the lattice $T$ from Lemma \ref{lem:dita-lattice}. As we will see below, this is not a coincidence, and these 2 lattices can be in fact identified, in a natural way. Our main result here is as follows:

\begin{thm}[The covering graph]\label{thm:covering}
Let $Q\in\mathbb T^{X\times Y}$ be dephased.
\begin{enumerate}
\item The group $T_Q$ acts freely on $B_Q$, by left multiplication, and $B_Q/T_Q=K_{X,Y}$. Thus $B_Q$ is a connected regular covering of $K_{X,Y}$, with group $T_Q$.

\item The map $m\to\sum_{i,j}m_{ij}x^iy^{-j}$ is an isomorphism $H_1(K_{X,Y},\mathbb Z)\simeq I_X\otimes I_Y$, and under this isomorphism the covering $B_Q$ corresponds to the subgroup $N_Q$. In particular, when $N_Q=\{1\}$, and so for generic $Q$, the graph $B_Q$ is the maximal abelian covering of $K_{X,Y}$.
\end{enumerate}
\end{thm}

\begin{proof}
1. The action on the white vertices is the left regular one, which is free. If $t\in T_Q$ fixes a coset $\gamma Y_i$, then $\gamma^{-1}t\gamma$ belongs to $T_Q\cap Y_i$, which is trivial, since the degree map is injective on $Y_i$ and vanishes on $T_Q$. Thus the action is free. Since $\Gamma_Q=T_Q\rtimes Y$, the orbits of the white vertices are labeled by their degrees $j\in Y$, and since each coset of $Y_i$ contains exactly one element of $T_Q$, there is one orbit of black vertices for each type $i\in X$. A white vertex of degree $j$ is joined to exactly one black vertex of each type, so the quotient is $K_{X,Y}$. Finally, $B_Q$ is connected, since $\gamma$ and $\gamma b^{(i)}$ are joined to $\gamma Y_i$, and the elements $b^{(i)}$ generate $\Gamma_Q$.

2. The elements of $I_X\otimes I_Y$ are the elements of $\mathbb Z[X\times Y]$ whose coefficients have vanishing sums on the rows and on the columns, which gives the isomorphism. We identify the fiber of $B_Q$ over a white vertex $j$ with $T_Q$, via $u\to uj^{(0)}$, and the fiber over a black vertex $i$ with $T_Q$ too, via $u\to uY_i$. As explained in the proof of Theorem \ref{thm:gram-all}, the white vertex $uj^{(0)}$ is joined to the black vertex $ut_{i,-j}Y_i$. Thus the lift of the closed path $j\to i\to j'\to i'\to j$ of $K_{X,Y}$ multiplies the coordinate in $T_Q$ by the image of the element
$$t_{i,-j}t_{i,-j'}^{-1}t_{i',-j'}t_{i',-j}^{-1}\to(x^i-x^{i'})(y^{-j}-y^{-j'})$$
where at right we have used the isomorphism of Lemma \ref{lem:dita-lattice}. This is precisely the image of the 4-cycle under the isomorphism in the statement, and since the 4-cycles generate the homology, the monodromy map $H_1(K_{X,Y},\mathbb Z)\to T_Q$ of the covering is the quotient map $I_X\otimes I_Y\to(I_X\otimes I_Y)/N_Q$, whose kernel is $N_Q$. The coverings with abelian group correspond to the subgroups of the homology, the maximal one corresponding to the trivial subgroup, see \cite{KS01,Sun13}.
\end{proof}

In terms of walks on these graphs, we have the following result:

\begin{prop}\label{prop:covering-walks}
The operator $J^*J$ is $M$ plus the adjacency operator of the Cayley graph of $\Gamma_Q$ with respect to the $M(N-1)$ generators $b^{(i)}$, with $i\in X$ and $b\in Y-\{0\}$. Also, if $w_n$ and $w_n^{(i)}$ denote the numbers of closed walks of length $2n$ in $B_Q$ based at a white vertex, and at a black vertex of type $i$, then
$$c_n=\frac{w_n}M=\frac1{MN}\sum_{i\in X}w_n^{(i)}$$
for any $n\geq1$.
\end{prop}

\begin{proof}
We have $J^*J\delta_\gamma=\sum_{i\in X,b\in Y}\delta_{\gamma b^{(i)}}$, and the $M$ terms with $b=0$ give $M\delta_\gamma$. The formula $c_n=\frac{w_n}M$ is the one in Theorem \ref{thm:dita-words}, and the last formula follows from the proof of Theorem \ref{thm:gram-all}, the diagonal entry of $(JJ^*)^n$ at a black vertex of type $i$ being $\int_{K_Q}((ZZ^*)^n)_{ii}dZ$.
\end{proof}

In terms of spectral measures, the first formula reads $\mu_H=(1-\frac1M)\delta_0+\frac1M law(J^*J)$, with the law being taken at a white vertex, while Theorem \ref{thm:gram-all} involves $JJ^*$, at the black vertices. Also, the graph $B_Q$ is amenable, $\Gamma_Q$ being an extension of an abelian group by a finite group, so the norm of $J^*J$ is $MN$, which is the index, in agreement with the coamenability observation in \cite{BO26}.

The phase group has a natural interpretation too. A matrix $Z\in\mathbb T^{X\times Y}$ can be viewed as a connection on $K_{X,Y}$, with the group $\mathbb T$ as structure group, and the multiplications of the rows and columns by phases are the gauge transformations. The dephased matrices are then the gauge classes, which form the torus $H^1(K_{X,Y},\mathbb T)=Hom(I_X\otimes I_Y,\mathbb T)$, the character associated to $Z$ being the holonomy, given on the above 4-cycles by
$$\frac{Z_{ij}Z_{i'j'}}{Z_{ij'}Z_{i'j}}$$
In this picture $K_Q$ is the group of the gauge classes having trivial holonomy along the cycles which belong to $N_Q$, the matrix $ZZ^*$ is the operator $JJ^*$ of the complete bipartite graph, twisted by the connection $Z$, and Theorem \ref{thm:gram-all} is the decomposition of $JJ^*$ over the characters of the covering group, in the sense of the Floquet--Bloch theory \cite{KS01,Sun13}. Observe also that the generators of Proposition \ref{prop:gram-generators} are second differences of the gauge class of $Q$ itself, with respect to the translations of $X$ and of $Y$: with $\partial_k$ and $\partial_e$ denoting the corresponding multiplicative differences, $g_{k,e}$ is the gauge class of $\partial_{-k}\partial_eQ$, up to inversion.

\subsection{Deformations with 2 rows}

For $X=\mathbb Z_2$ write the second row of the dephased matrix $Q$ as $(q_b)_{b\in Y}$, with $q_0=1$. For $Z\in K_Q$ we write $Z_b=Z_{1b}$, so that $Z_0=1$.

\begin{cor}[The blowup for 2 rows]\label{cor:two-rows-abelian}
Let $H=F_2\otimes_QF_Y$. The group $K_Q\subset\mathbb T^Y$ is the closed subgroup generated by the vectors
$$g_s(b)=\Big(\frac{q_0q_{s+b}}{q_sq_b}\Big)^2,\quad s,b\in Y$$
If $(Z_b)_{b\in Y}$ is Haar distributed on $K_Q$, then
$$\mu_H=\Big(1-\frac1N\Big)\delta_0+\frac1N law\Big(N+\varepsilon\Big|\sum_{b\in Y}Z_b\Big|\Big)$$
where $\varepsilon$ is an independent uniform sign.
\end{cor}

\begin{proof}
Since $-1=1$ in $\mathbb Z_2$, the elements of \eqref{eq:gram-generators} are $R_1(1,e)=q_e^2$ and $R_k(i,e)=1$ if $i=0$ or $k=0$. Hence the only nontrivial generators are $g_{1,s}(1,b)=g_s(b)$, since $q_0=1$. The rows of $Z$ are $(1)_{b\in Y}$ and $(Z_b)_{b\in Y}$, so that
$$ZZ^*=\begin{pmatrix}N&\bar S\\S&N\end{pmatrix},
\quad S=\sum_{b\in Y}Z_b$$
Its eigenvalues are $N\pm|S|$, each with weight $\frac12$ for the normalized trace, and Theorem \ref{thm:gram-all} gives the result.
\end{proof}

For a row which is not dephased, Proposition \ref{prop:gram-generators} shows that one may take $q_b=Q_{1b}/Q_{0b}$ in the formula for $g_s(b)$, which does not change if all these ratios are multiplied by the same scalar.

The white vertices of $B_Q$ have valence 2. Suppressing them gives the coset graph $C_Q$, with vertex sets $\Gamma_Q/Y_0$ and $\Gamma_Q/Y_1$; its edge $\gamma$ joins $\gamma Y_0$ to $\gamma Y_1$. Using the representatives of degree zero identifies both vertex sets with $T_Q$, and the left vertex $u$ is joined to $ut_{1,-b}$ on the right, for $b\in Y$. Thus $C_Q$ is $N$-regular and, by Theorem \ref{thm:covering}, is an abelian covering of the graph with 2 vertices and $N$ edges. In the generic case we have:

\begin{prop}\label{prop:two-rows}
Let $X=\mathbb Z_2$ and let $Q$ be generic. Let $\zeta_b$, $b\in Y$, be independent random variables uniformly distributed on $\mathbb T$, and let $s=\sum_{b\in Y}\zeta_b$. Then $\mu_H=(1-\frac1N)\delta_0+\frac1N law(N+\varepsilon|s|)$, and
\begin{equation}\label{eq:two-rows-moments}
c_n=\frac1N\sum_{j=0}^{[n/2]}\binom n{2j}N^{n-2j}W_N(2j),
\quad
W_N(2j)=\mathbb E|s|^{2j}=\sum_{(j_b)\in\mathbb N^Y,\sum_b j_b=j}\binom j{(j_b)_{b\in Y}}^2
\end{equation}
for any integer $n\geq1$.
\end{prop}
\begin{proof}
By Theorem \ref{thm:gram}, and since multiplying each column of $Z$ by the conjugate of its first entry does not change $ZZ^*$, we can assume that the first row of $Z$ is $(1,\ldots,1)$ and that the second row consists of independent uniform phases $\zeta_b$. Hence
$$ZZ^*=\begin{pmatrix}N&\bar s\\s&N\end{pmatrix}$$
whose eigenvalues are $N \pm |s|$. The normalized trace amounts to choosing the sign with equal probabilities, which gives the law. Expanding $(N+\varepsilon|s|)^n$, the odd powers of $\varepsilon$ vanish after averaging over the sign, which gives the first formula in \eqref{eq:two-rows-moments}.

Finally, $|s|^{2j}=\sum \zeta_{a_1}\ldots\zeta_{a_j}\bar\zeta_{b_1}\ldots\bar\zeta_{b_j}$, and the expectation of a term is one if the multisets $\{a_1,\ldots,a_j\}$ and $\{b_1,\ldots,b_j\}$ coincide and zero otherwise. Counting the pairs of words with the same content gives the last formula.
\end{proof}

The numbers $W_N(2j)$ are the even moments of the length of the $N$-step planar random walk with unit steps \cite{BSWZ12}. For $N=3$ a multinomial coefficient with entries $(j-k,a,k-a)$ is $\binom jk\binom ka$, and $\sum_a\binom ka^2=\binom{2k}k$. For $N=4$ we can proceed similarly, by splitting the 4 entries into 2 pairs, and we obtain
$$W_3(2j)=\sum_{k=0}^j\binom jk^2\binom{2k}k,\quad
W_4(2j)=\sum_{k=0}^j\binom jk^2\binom{2k}k\binom{2j-2k}{j-k}$$
Thus $W_3(2j)=1,3,15,93,639,\ldots$ and for $N=4$ we obtain the Domb numbers $1,4,28,256,2716,\ldots$ These numbers count walks on some well-known lattices, as follows:

\begin{prop}\label{prop:two-rows-graph}
Let $X=\mathbb Z_2$ and let $Q$ be generic. Then the coset graph $C_Q$ is the maximal abelian covering of the graph with $2$ vertices and $N$ edges, the number of closed walks of length $2j$ at a vertex of $C_Q$ is $W_N(2j)$, and for $N=2,3,4$ this graph is the line, the honeycomb lattice, and the diamond lattice.
\end{prop}
\begin{proof}
The first assertion follows from Theorem \ref{thm:covering} (2). The adjacency operator of $C_Q$, from the left vertices to the right ones, is the sum of the translations by the elements $t_{1,-b}$, which becomes $\sum_bZ_b$ after Fourier transform, so the number of closed walks of length $2j$ is
$$\int_{K_Q}\Big|\sum_bZ_b\Big|^{2j}dZ=W_N(2j)$$
Finally, with $T=\mathbb Z^{N-1}$ and with $e_b=t_{1,-b}$ as basis, the left vertex $u$ is joined to the right vertices $u,u+e_1,\ldots,u+e_{N-1}$, and for $N=2,3,4$ these are the standard descriptions of the line, of the honeycomb lattice, and of the diamond lattice, see \cite{Sun13}.
\end{proof}

For $N=3$, so that the index is $6$, the formula \eqref{eq:two-rows-moments} gives $c_2=4$, $c_3=18$ and $c_4=\frac13(81+6\cdot9\cdot3+15)=86$. For the trivial deformation $Q=1$ one obtains instead the moments $6^{n-1}$ of the Fourier matrix of $\mathbb Z_6$. The program \texttt{jalabert.py} compares the moments \eqref{eq:two-rows-moments} with the word counts \eqref{eq:dita-words}, and with the numbers of closed walks on the honeycomb, kagome and diamond lattices.

At index 4 the graph with 2 vertices and 2 edges is a cycle, whose coverings are cycles too, and we recover a known formula. Take $X=Y=\mathbb Z_2$ and let $q$ be the only free entry of $Q$. Then $T=\mathbb Z w$ with $w=(x-1)(y-1)$, and Corollary \ref{cor:two-rows-abelian} gives $K_Q=\langle q^4\rangle$. Thus $N_Q=\{1\}$ if $q$ has infinite order, and if $q^4$ has finite order $l$, then $N_Q=l\mathbb Z w$ and $\dim C(G_H)=4l$.

\begin{prop}
Let $M=N=2$ and let $q$ be the free entry of the dephased matrix $Q$. If $q^4$ has finite order $l$, the rooted law is
\begin{equation}\label{eq:four-points}
 \mu_H=\frac12\delta_0+\frac12\sigma_*d_{2l},
 \quad \sigma(u)=(u+u^{-1})^2
\end{equation}
If $q$ has infinite order, the same formula holds with the uniform measure on $\mathbb T$ in place of $d_{2l}$.
\end{prop}
\begin{proof}
By Corollary \ref{cor:two-rows-abelian} we have $\mu_H=\frac12\delta_0+\frac12 law(2+\varepsilon|1+Z|)$, with $Z$ uniform on the $l$-th roots of unity. With $Z=e^{2\pi i k/l}$ the 2 values of $2+\varepsilon|1+Z|$ are $2\pm2\cos(\pi k/l)$, and these are the values of $\sigma(u)=2+2\cos(\pi m/l)$ at $u=e^{2\pi i m/(4l)}$, for $m=k$ and $m=k+l$ modulo $2l$. Since $d_{2l}$ is the uniform measure on the $4l$-th roots of unity, this gives the result. For $q$ of infinite order the proof is similar. Here $C_Q$ is a cycle of length $2l$, with $2$ parallel edges if $l=1$, or a line.
\end{proof}

For finite $l$, the measure \eqref{eq:four-points} is the squared law of $\widetilde D_{2l+2}$ at an extending vertex by the $A,D$ formulas of Section \ref{sec:circle}. In particular, phases $q$ of orders $3$ and $12$ give the same rooted law.

Let us go back now to the general case of 2 rows. In terms of quantum groups, the group $\Gamma_{X,Y}=I_Y\rtimes Y$ appearing in the generic case is a crystallographic group, with $Y$ acting on $I_Y\simeq\mathbb Z^{N-1}$ by multiplication, and
$$C(G_H)=C^*(\Gamma_{X,Y})\rtimes C(\mathbb Z_2)$$
For $Y=\mathbb Z_3$ this group is the wallpaper group $p3$, generated by 2 rotations of order 3, namely $a=1^{(0)}$ and $b=1^{(1)}$, with presentation $a^3=b^3=(ab)^3=1$, and its Cayley graph with respect to $a^{\pm1},b^{\pm1}$ is the kagome lattice, as explained in \cite{ABF26}. In terms of the rooted law, we have:

\begin{thm}\label{thm:kagome}
Let $H=F_2\otimes_QF_3$, with $Q$ generic, or more generally with $N_Q=\{1\}$. Then
$$\mu_H=\frac12\delta_0+\frac12 law(2+A)$$
where $A$ is the adjacency operator of the kagome lattice, that is, of the Cayley graph of $\Gamma_{X,Y}$ with respect to $a^{\pm1},b^{\pm1}$, and the law is taken at any vertex.
\end{thm}
\begin{proof}
The elements of $\Gamma_{X,Y}$ are the edges of the coset graph $C_Q$, which is the honeycomb lattice, by Proposition \ref{prop:two-rows-graph}. Now 2 edges $\gamma\neq\gamma'$ have a common vertex precisely when $\gamma Y_i=\gamma'Y_i$ for some $i$, that is, when $\gamma'=\gamma b^{(i)}$ with $b\neq0$. Thus the Cayley graph in the statement is the line graph of the honeycomb lattice, which is the kagome lattice, in agreement with \cite{ABF26}. The formula for $\mu_H$ follows from Proposition \ref{prop:covering-walks}.
\end{proof}

The honeycomb and kagome lattices are shown in Figure \ref{fig:lattices}. The vertices of the kagome lattice are the midpoints of the edges of the honeycomb lattice, and 2 such midpoints are joined when the corresponding edges have a common vertex.

\begin{figure}[htbp]
\centering
\setlength{\unitlength}{7mm}
\begin{picture}(8.7,5)(-0.9,-1)
\linethickness{0.7pt}
\multiput(0,0)(0.866,1.5){3}{\multiput(0,0)(1.732,0){4}{\polygon(0,1)(0.866,0.5)(0.866,-0.5)(0,-1)(-0.866,-0.5)(-0.866,0.5)\multiput(0,-1)(0,2){2}{\circle*{0.2}}\multiput(-0.866,-0.5)(1.732,0){2}{\circle*{0.2}}\multiput(-0.866,0.5)(1.732,0){2}{\circle*{0.2}}}}
\end{picture}\quad\quad
\begin{picture}(8.7,5)(-0.9,-1)
{\color{magenta}\multiput(0,0)(0.866,1.5){3}{\multiput(0,0)(1.732,0){4}{\polygon(0,1)(0.866,0.5)(0.866,-0.5)(0,-1)(-0.866,-0.5)(-0.866,0.5)}}}
\linethickness{0.7pt}
\multiput(0,0)(0.866,1.5){3}{\multiput(0,0)(1.732,0){4}{\polygon(0.866,0)(0.433,0.75)(-0.433,0.75)(-0.866,0)(-0.433,-0.75)(0.433,-0.75)\multiput(-0.866,0)(1.732,0){2}{\circle*{0.2}}\multiput(-0.433,-0.75)(0.866,0){2}{\circle*{0.2}}\multiput(-0.433,0.75)(0.866,0){2}{\circle*{0.2}}}}
\end{picture}
\caption{The honeycomb lattice, and its line graph, the kagome lattice.}
\label{fig:lattices}
\end{figure}

The kagome lattice is known to have a flat band, namely the eigenvalue $-2$ of $A$, of weight $\frac13$. This comes from the identity $2+A=J^*J$, the honeycomb lattice having 3 edges and 2 vertices in each cell, so that $J^*J$ has a kernel of weight $\frac13$. Via $\frac12+\frac12\cdot\frac13=\frac23$, this flat band is the atom $\frac23\delta_0$ of the formula in Proposition \ref{prop:two-rows}.

\subsection{The index 6}

At index 6 we take $Y=\mathbb Z_3$ and write the dephased second row as $(1,p,q)$. For $Z\in K_Q$ we write $(Z_1,Z_2)$ for its free coordinates.

\begin{thm}[All parameters at index 6]\label{thm:six-blowup}
Put $\alpha=p^2$ and $\beta=q^2$. The group $K_Q$ is the closed subgroup of $\mathbb T^2$ generated by
$$(\beta\alpha^{-2},(\alpha\beta)^{-1})
\quad{\rm and}\quad
((\alpha\beta)^{-1},\alpha\beta^{-2})$$
and the rooted law of $F_2\otimes_QF_3$ is
$$\mu_H
=\frac23\delta_0+\frac13 law(3+\varepsilon|1+Z_1+Z_2|)$$
where $(Z_1,Z_2)$ is Haar distributed on $K_Q$. If $p$ or $q$ has infinite order, then $K_Q=\mathbb T^2$. If $p$ and $q$ are roots of unity, $K_Q$ is the finite group generated by the 2 displayed elements.
\end{thm}

\begin{proof}
The generators of Corollary \ref{cor:two-rows-abelian} are $g_0=1$, $g_1=(q^2p^{-4},p^{-2}q^{-2})$ and $g_2=(p^{-2}q^{-2},p^2q^{-4})$, which are the displayed elements. The formula for $\mu_H$ is Corollary \ref{cor:two-rows-abelian}. A character $(z_1,z_2)\to z_1^hz_2^k$ is trivial on $K_Q$ if and only if
$$\alpha^{-2h-k}\beta^{h-k}=1
\quad{\rm and}\quad
\alpha^{k-h}\beta^{-h-2k}=1$$
The matrix of exponents has determinant $(2h+k)(h+2k)+(h-k)^2=3(h^2+hk+k^2)$, which is positive for $(h,k)\neq(0,0)$. Multiplying the 2 relations by the adjugate matrix, we see that a nonzero relation gives $\alpha^{3(h^2+hk+k^2)}=\beta^{3(h^2+hk+k^2)}=1$, so that $p$ and $q$ are roots of unity. Hence if $p$ or $q$ has infinite order, no nontrivial character is trivial on $K_Q$, and $K_Q=\mathbb T^2$. If $p$ and $q$ are roots of unity, the 2 generators have finite order and generate a finite group.
\end{proof}

The first assertion is the dual form of \cite[Theorem 4.6]{BO26}, where the group $N_Q$ is computed, and where the ring of Eisenstein integers appears, in the computation of the order of $G_H$. If $p$ has infinite order, the law is the generic law of Proposition \ref{prop:two-rows} for every $q$, including dependent choices such as $q=p^3$. This agrees with the case $N=3$ of \cite[Theorem 5.1]{Bic15}.

\begin{cor}[The row $(1,1,q)$]\label{cor:six-one-phase}
For $p=1$, the random pair $(Z_1,Z_2)$ of Theorem \ref{thm:six-blowup} has the law of $(U,\bar UV)$, where $U$ and $V$ are independent and Haar distributed on $\langle q^2\rangle$ and on $\langle q^6\rangle$. If $q^2$ has finite order $l$, the group $K_Q$ has $l^2/\gcd(l,3)$ elements.
\end{cor}

\begin{proof}
With $\alpha=1$ the generators of Theorem \ref{thm:six-blowup} are $(\beta,\beta^{-1})$ and $(\beta^{-1},\beta^{-2})$. The subgroup generated by them consists of the pairs $(\beta^{a-b},\beta^{-a-2b})=(U,\bar UV)$ with $U=\beta^{a-b}$ and $V=\beta^{-3b}$, $a,b\in\mathbb Z$. As $(a,b)$ runs over $\mathbb Z^2$, the pair $(U,V)$ runs over $\{\beta^u:u\in\mathbb Z\}\times\{\beta^{3v}:v\in\mathbb Z\}$. The map $(u,v)\to(u,\bar uv)$ is an automorphism of the compact group $\mathbb T^2$, so that $K_Q$ is the image of $\langle\beta\rangle\times\langle\beta^3\rangle$ and the Haar measures correspond. In the finite case, $\langle\beta\rangle$ has $l$ elements and $\langle\beta^3\rangle$ has $l/\gcd(l,3)$ elements.
\end{proof}

We can now describe all the quantum groups at index 6. Let $\omega=e^{2\pi i/3}$ and $\zeta=1+\omega=e^{\pi i/3}$, and let $\mathbb Z[\omega]=\mathbb Z\oplus\mathbb Z\zeta$ be the ring of Eisenstein integers, whose norm is given by $|k+h\zeta|^2=k^2+kh+h^2$. The honeycomb lattice of Proposition \ref{prop:two-rows-graph} has $T$ as group of translations, and for a nonzero ideal $I=(b+c\zeta)$ of $\mathbb Z[\omega]$ the quotient of the hexagonal tessellation by the translations in $I$ is the map $\{6,3\}_{b,c}$ on the torus of Coxeter and Moser \cite{CM80}, with $b^2+bc+c^2$ hexagonal faces. The units being the powers of $\zeta$, any nonzero ideal has a unique generator $b+c\zeta$ with $b\geq1$ and $c\geq0$.

\begin{thm}\label{thm:eisenstein}
Let $H=F_2\otimes_QF_3$, with the dephased second row of $Q$ being $(1,p,q)$, and put $\alpha=p^2$, $\beta=q^2$. Identify $T$ with $\mathbb Z[\omega]$, via $t_{11}\to1$ and $t_{12}\to\zeta$.
\begin{enumerate}
\item $N_Q$ is an ideal of $\mathbb Z[\omega]$, which is zero precisely when $p$ or $q$ has infinite order.

\item Every ideal occurs. More precisely, for $k+h\zeta\neq0$, with norm $n=k^2+kh+h^2$, the parameters
$$\alpha=e^{-2\pi i(h+2k)/(3n)},\quad\beta=e^{2\pi i(h-k)/(3n)}$$
give the ideal $N_Q=(k+h\zeta)$.
\end{enumerate}
\end{thm}

\begin{proof}
1. By Lemma \ref{lem:dita-lattice} the conjugation by $(-1)^{(0)}$ is the multiplication by $y$ on $T=(x-1)I_Y$. Since $1+y+y^2$ acts by zero on $I_Y$, the group $T$ is a module over $\mathbb Z[y]/(1+y+y^2)=\mathbb Z[\omega]$, which is free of rank 1, generated by $t_{11}=(x-1)(y-1)$, and we have $t_{12}=(1+y)t_{11}$, which corresponds to $1+\omega=\zeta$. Now $N_Q$ is a normal subgroup of $\Gamma_{X,Y}$, so it is stable under the multiplication by $y$, and so it is an ideal. The last assertion follows from Theorem \ref{thm:six-blowup}.

2. The free coordinates of $Z\in K_Q$ are $Z_1=\chi(t_{1,-1})=\chi(t_{12})$ and $Z_2=\chi(t_{1,-2})=\chi(t_{11})$, so by \eqref{eq:gram-group} the element $k+h\zeta$ belongs to $N_Q$ precisely when the character $Z_1^hZ_2^k$ is trivial on $K_Q$. By the proof of Theorem \ref{thm:six-blowup}, with $\alpha=e^{2\pi i s}$ and $\beta=e^{2\pi i s'}$, this happens when
$$A(k+h\zeta)\begin{pmatrix}s\\ s'\end{pmatrix}\in\mathbb Z^2,\quad
A(k+h\zeta)=\begin{pmatrix}-(2h+k)&h-k\\ k-h&-(h+2k)\end{pmatrix}$$
The matrix $A(z)$ is additive in $z\in\mathbb Z[\omega]$, and satisfies $A(\zeta z)=PA(z)$, with $P=(^1_1{\ }^{-1}_{\ 0})$, as one checks on the formula. Since $P^2-P+1=0$, we have $A(z)=\iota(z)A(1)$, where $\iota$ is the embedding of the field $\mathbb Q(\omega)$ into $M_2(\mathbb Q)$ given by $\iota(\zeta)=P$.

Now let $z_0=k+h\zeta\neq0$. The matrix $A(z_0)$ has determinant $3n$, and the parameters in the statement are those given by $(s,s')^t=A(z_0)^{-1}(1,0)^t$. For this choice, an element $z$ belongs to $N_Q$ precisely when the vector $\iota(z/z_0)(1,0)^t$ has integer entries, and since $\iota(k'+h'\zeta)(1,0)^t=(k'+h',h')^t$, this happens precisely when $z/z_0\in\mathbb Z[\omega]$.
\end{proof}

When the ideal is nonzero, the quantum group and its law can be described as follows:

\begin{thm}\label{thm:eisenstein-finite}
In the above setting, assume $N_Q=(b+c\zeta)$, with $b\geq1$ and $c\geq0$, and let $n=b^2+bc+c^2$. Then $G_H$ is finite, of order $6n$, the coset graph $C_Q$ is the graph of the map $\{6,3\}_{b,c}$, with $2n$ vertices, and
$$\mu_H=\frac23\delta_0+\frac1{6n}\sum_\chi\left(\delta_{3+|1+\chi(1)+\chi(\zeta)|}+\delta_{3-|1+\chi(1)+\chi(\zeta)|}\right)$$
with the sum being over the $n$ characters $\chi$ of the group $\mathbb Z[\omega]/N_Q$. Also, conjugate ideals give the same rooted law.
\end{thm}

\begin{proof}
By Theorem \ref{thm:eisenstein} the parameters $p,q$ are roots of unity, so $G_H$ is finite, of order $6[T:N_Q]$, by Proposition \ref{prop:dita-criterion}, and the index of an ideal is the norm of its generator. The coset graph is the quotient of the honeycomb lattice by the group of translations $N_Q$, as explained above, and the formula for $\mu_H$ is the one in Theorem \ref{thm:six-blowup}, with $(Z_1,Z_2)=(\chi(\zeta),\chi(1))$.

Regarding the last assertion, the conjugate of $k+h\zeta$ is $(k+h)-h\zeta=\zeta^{-1}(h+k\zeta)$, so passing to the conjugate ideal amounts to exchanging $Z_1$ and $Z_2$, and the number $|1+Z_1+Z_2|$ is symmetric in $Z_1,Z_2$.
\end{proof}

The orders of the phase groups are therefore the numbers of the form $b^2+bc+c^2$, namely $1,3,4,7,9,12,13,16,19,21,\ldots$ and for the first values the graphs obtained are well-known ones. Here are the first 5 cases, with the graphs being identified by an isomorphism test, done by the program \texttt{jalabert.py}, which also checks the above 2 theorems, for the ideals of norm at most 60.

\begin{example}[The complete bipartite graph and the cube]\label{ex:six-cube}
Take $p=1$. At $q=e^{2\pi i/3}$, the pair of Corollary \ref{cor:six-one-phase} is $(U,\bar U)$ with $U^3=1$. Here $T_Q\simeq\mathbb Z_3$, the ideal is $(1+\zeta)$, and the 3 right neighbors of a left vertex exhaust the right vertex set, which gives $K_{3,3}$. At $q=i$, the 2 coordinates are independent uniform signs. Here $T_Q\simeq\mathbb Z_2^2$, the ideal is $(2)$, and the 3 translations omit one perfect matching from $K_{4,4}$, which gives the cube. In these 2 cases the rooted laws are
$$\frac{13\delta_0+4\delta_3+\delta_6}{18}
\quad{\rm and}\quad
\frac{17\delta_0+3\delta_2+3\delta_4+\delta_6}{24}$$
Indeed, $|1+Z_1+Z_2|$ takes the values $3,0,0$ in the first case and $3,1,1,1$ in the second. At $q=\pm1$ the group $K_Q$ is trivial and the law is $\frac{5\delta_0+\delta_6}6$.
\end{example}

The next ideals are those of norm 7, namely $(2+\zeta)$ and its conjugate, which give the same graph and the same law, as follows:

\begin{example}[The Heawood graph]
Let $\eta=e^{2\pi i/7}$ and take $(p,q)=(\eta,\eta^3)$. The generators of Theorem \ref{thm:six-blowup} are $(\eta^2,\eta^6)$ and $(\eta^6,\eta^4)$, which is the cube of the first one. Thus the random row is $(1,u,u^3)$ with $u$ uniform on the seventh roots of unity. The 6 differences of 2 distinct elements of $\{0,1,3\}$ are the nonzero residues modulo $7$. Hence, for $u\neq1$, we have
$$|1+u+u^3|^2=3+\sum_{j=1}^6u^j=2$$
At $u=1$ the length is 3, and Theorem \ref{thm:six-blowup} gives
$$\mu_H=\frac{29}{42}\delta_0+\frac17(\delta_{3-\sqrt2}+\delta_{3+\sqrt2})+\frac1{42}\delta_6$$
Here $T_Q\simeq\mathbb Z_7$, the ideal has norm 7, and the right neighbors of $u$ are $u$, $u+1$ and $u+3$. Since $\{0,1,3\}$ is a perfect difference set modulo $7$, any 2 distinct points lie in exactly one translate of $\{0,1,3\}$. The 7 translates are the lines of the Fano plane, and the coset graph is its incidence graph, which is the Heawood graph, coming from the map $\{6,3\}_{2,1}$.
\end{example}

At norms 9 and 12 we have the ideals $(3)$ and $(2+2\zeta)$, where the rooted law gets closer to the generic one:

\begin{example}[The Pappus and Nauru graphs]\label{ex:six-twelfth}
For the ideal $(3)$ the phase group consists of the pairs of cubic roots of unity, the length $|1+Z_1+Z_2|$ takes the values $3,0,\sqrt3$ with multiplicities $1,2,6$, and we obtain the Pappus graph, coming from the map $\{6,3\}_{3,0}$, and the law
$$\mu_H=\frac{37\delta_0+4\delta_3+\delta_6}{54}+\frac19\left(\delta_{3-\sqrt3}+\delta_{3+\sqrt3}\right)$$
Next, take $p=1$ and $q=e^{2\pi i/12}$. In Corollary \ref{cor:six-one-phase}, $U$ is uniform on the sixth roots of unity and $V$ is a uniform sign. The length $|1+U+\bar UV|$ takes the values $0,1,2,3$ with multiplicities $2,3,6,1$ among the 12 pairs. We obtain
$$\mu_H
=\frac{49\delta_0+6\delta_1+3\delta_2+4\delta_3+3\delta_4+6\delta_5+\delta_6}{72}$$
Here the ideal is $(2+2\zeta)$, and the coset graph is the Nauru graph, coming from the map $\{6,3\}_{2,2}$. The first 7 moments agree with the generic ones given by \eqref{eq:two-rows-moments}. For $n=8$ this measure gives $c_8=58986$, whereas $c_8^{gen}=58974$. By Theorem \ref{thm:rigidity}, the shortest nontrivial element of $N_Q$ has length $8$. The moments of the 3 rational laws above, with denominators 18, 24 and 72, are computed by the program \texttt{wordgram.py} too, including the difference $c_8-c_8^{gen}=12$.
\end{example}

By Theorem \ref{thm:eisenstein-finite} the rooted law does not distinguish an ideal from its conjugate. We do not know if this is the only loss of information. The program \texttt{jalabert.py} shows that the rooted laws of the 243 ideals of norm at most 400 are distinct, up to conjugation.

The above results explain the numerical data in \cite{ABF26}, where the matrices $F_2\otimes_QF_3$ appear as limits of the Karlsson matrices, with second row $(1,z_1,z_2)$. The limit of order 24 found there has $(\alpha,\beta)=(e^{2\pi i/6},e^{2\pi i/3})$, so the 2 generators of Theorem \ref{thm:six-blowup} are $(1,-1)$ and $(-1,-1)$, and this is the case of the cube, with the ideal being $(2)$. Thus, by Example \ref{ex:six-cube}, we have
$$c_n=\frac{3\cdot2^n+3\cdot4^n+6^n}{24}$$
which gives $1,4,18,88,456,2464,13728$, and by \eqref{eq:annular} the annular multiplicities at the levels $3,\ldots,7$ are $1,10,29,114,419$, which are the numbers found there by floating point computations.

Similarly, for $K_{3,3}$ we obtain $c_n=\frac1{18}(4\cdot3^n+6^n)$, starting with $1,4,18,90$, and annular multiplicities $1,12,39,152,559$, which are the numbers found there at the limits having $c_4=90$. Finally, the limit with $(\alpha,\beta)=(e^{2\pi i/6},e^{4\pi i/3})$ considered there has a phase group of order 12, so the ideal is $(2+2\zeta)$, as for the Nauru graph in Example \ref{ex:six-twelfth}, and this is why this limit, of order 72, has the same invariants as the generic one up to the level 7, the annular multiplicities at the levels $3,\ldots,7$ being $1,8,19,66,209$ in both cases. These computations are done by \texttt{jalabert.py} too.

\subsection{The generic law}

In the remaining case, where the ideal is zero, the law is that of the honeycomb and kagome lattices, by Proposition \ref{prop:two-rows-graph} and Theorem \ref{thm:kagome}, and it can be computed explicitly. Consider the complete elliptic integral of the first kind, with modulus $k\in[0,1)$,
$$K(k)=\int_0^{\pi/2}\frac{d\varphi}{\sqrt{1-k^2\sin^2\varphi}}$$

\begin{thm}\label{thm:density}
Let $H=F_2\otimes_QF_3$, with $p$ or $q$ of infinite order. Then
$$\mu_H=\frac23\delta_0+\frac16p_3(|x-3|)dx$$
on $[0,6]$, where $p_3$ is the density of the length of the 3-step planar random walk, given by
$$p_3(r)=\frac{4r}{\pi^2\sqrt{(1+r)^3(3-r)}}K\left(\sqrt{\frac{16r}{(1+r)^3(3-r)}}\right),\quad
p_3(r)=\frac{\sqrt r}{\pi^2}K\left(\sqrt{\frac{(1+r)^3(3-r)}{16r}}\right)$$
for $0<r<1$, and for $1<r<3$. Thus the density of $\mu_H$ on $(0,6)$ is continuous, except for logarithmic singularities at $x=2$ and $x=4$, it vanishes linearly at $x=3$, and it tends to $\frac{\sqrt3}{12\pi}$ at the endpoints $0$ and $6$.
\end{thm}

\begin{proof}
By Theorem \ref{thm:six-blowup} we have $K_Q=\mathbb T^2$, and $\mu_H=\frac23\delta_0+\frac13 law(3+\varepsilon|S|)$ with $S=1+U+V$, where $U,V$ are independent and uniform on $\mathbb T$. Thus the density of $\mu_H$ on $(0,6)$ is $\frac16p_3(|x-3|)$, where $p_3$ is the density of $|S|$. The number $|S|$ is the length of a 3-step planar random walk with unit steps, and also the absolute value of the energy in the nearest neighbor model on the honeycomb lattice, and its density is known. The above formulas are the classical ones for the density of states of the honeycomb lattice, due to Hobson and Nierenberg \cite{HN53}, see also \cite{CGPNG09}, and an equivalent form, due to Borwein, Straub, Wan and Zudilin \cite{BSWZ12}, is
$$p_3(r)=\frac{2\sqrt3r}{\pi(3+r^2)}{}_2F_1\left(\frac13,\frac23;1;\frac{r^2(9-r^2)^2}{(3+r^2)^3}\right)$$
The 2 moduli in the statement are inverse to each other, and equal 1 at $r=1$, with
$$1-\frac{16r}{(1+r)^3(3-r)}=\frac{(1-r)^3(3+r)}{(1+r)^3(3-r)}$$
so the logarithmic singularity of $K$ at $k=1$ gives a logarithmic singularity of $p_3$ at $r=1$. At $r=0$ and $r=3$ the moduli vanish, and $K(0)=\frac\pi2$ gives $p_3(r)\sim\frac{2r}{\pi\sqrt3}$ and $p_3(3)=\frac{\sqrt3}{2\pi}$. The program \texttt{jalabert.py} checks numerically that the 2 forms of $p_3$ agree, and that the moments of $p_3$ and of $\mu_H$ are the numbers $W_3(2j)$ and $c_n$ of Proposition \ref{prop:two-rows}.
\end{proof}

In the language of the honeycomb lattice, the point $x=3$ corresponds to the Dirac points, the singularities at $x=2,4$ are the van Hove singularities, and the atom at $0$ is the flat band of the kagome lattice. Observe that the singularities are located at $x=2$ and $x=4$, which are values of the Jones index. The above density makes the link with the theory at index $\leq4$, where the circular measures are given by trigonometric densities: at index 6 the trigonometric functions are replaced by elliptic ones.

\section{Growth exponents, indices 8 and 9}\label{sec:growth}

We first find the asymptotic behavior of the moments $c_n$, for any parameter matrix, and then we discuss the deformations at indices 8 and 9, where new phenomena appear.

\subsection{The growth of the moments}

Let $Q\in\mathbb T^{X\times Y}$ be dephased, and consider the following objects:
\begin{enumerate}
\item The dimension $d=\dim K_Q$, which equals $(M-1)(N-1)-rank(N_Q)$.

\item The number $\kappa$ of connected components of $K_Q$, which is the order of the torsion subgroup of $T_Q$.

\item The lattice $\Lambda^*_Q$ of the integer matrices $m=(m_{ij})_{i,j\neq0}$ which are orthogonal to $N_Q$, in the sense that $\sum_{i,j}m_{ij}r_{ij}=0$ whenever $\sum_{i,j}r_{ij}t_{i,-j}\in N_Q$. This lattice has rank $d$.

\item The quadratic form $q(\theta)=\sum_{i\in X,j\in Y}(\theta_{ij}-\theta_{i\cdot}-\theta_{\cdot j}+\theta_{\cdot\cdot})^2$ on the real matrices $\theta\in\mathbb R^{X\times Y}$ whose first row and column vanish, where $\theta_{i\cdot},\theta_{\cdot j},\theta_{\cdot\cdot}$ are the averages on the rows, on the columns, and the total average. In terms of the free entries, this is the quadratic form with matrix $(1-\frac1M\mathbb I_{M-1})\otimes(1-\frac1N\mathbb I_{N-1})$, with $\mathbb I$ being the all-one matrices.

\item The determinant $\Delta_Q$ of the Gram matrix of $q$, with respect to a basis of $\Lambda_Q^*$.
\end{enumerate}

\begin{thm}\label{thm:growth}
For any dephased $Q\in\mathbb T^{X\times Y}$ we have
$$c_n\sim C(MN)^nn^{-d/2},\quad C=\frac1{MN\kappa\sqrt{\Delta_Q}}\left(\frac{MN}{4\pi}\right)^{d/2}$$
In particular, for generic $Q$ we have
$$c_n^{gen}\sim\frac{M^{(N-1)/2}N^{(M-1)/2}}{MN}\left(\frac{MN}{4\pi n}\right)^{d/2}(MN)^n$$
as $n\to\infty$, with $d=(M-1)(N-1)$.
\end{thm}

\begin{proof}
We use Theorem \ref{thm:gram-all}. Since $Tr(ZZ^*)=MN$, the largest eigenvalue $\lambda(Z)$ of $ZZ^*$ satisfies $\lambda(Z)\leq MN$, with equality precisely when $Z$ has rank 1, which for a matrix with first row and column equal to 1 means $Z=1$. The other eigenvalues have sum $MN-\lambda(Z)$, so
\begin{equation}\label{eq:growth-trace}
\lambda(Z)^n\leq Tr((ZZ^*)^n)\leq\lambda(Z)^n+(M-1)(MN-\lambda(Z))^n
\end{equation}

We first study $\lambda$ near $Z=1$. Write $Z=e^{i\theta}$, entrywise, with $\theta$ small, having vanishing first row and column, and let $u\in\mathbb C^X$ be the normalized all-one vector, which is the eigenvector of $ZZ^*$ at $\theta=0$, the other eigenvalues being then 0. The second order perturbation formula gives
$$\lambda(Z)=\langle ZZ^*u,u\rangle+\frac{\|(1-P)B_1u\|^2}{MN}+O(|\theta|^3)$$
where $P$ is the projection on $u$, and $B_1=i(\theta\mathbb I^t-\mathbb I\theta^t)$ is the first order term of $ZZ^*$, with $\mathbb I$ being the all-one $M\times N$ matrix. With $\theta_{i\cdot},\theta_{\cdot j},\theta_{\cdot\cdot}$ as above, we have
$$\langle ZZ^*u,u\rangle=\frac1M\sum_j\Big|\sum_iZ_{ij}\Big|^2=MN-\sum_{i,j}(\theta_{ij}-\theta_{\cdot j})^2+O(|\theta|^4)$$
and $\|(1-P)B_1u\|^2=MN^2\sum_i(\theta_{i\cdot}-\theta_{\cdot\cdot})^2$. Since $\sum_{i,j}(\theta_{ij}-\theta_{\cdot j})^2-N\sum_i(\theta_{i\cdot}-\theta_{\cdot\cdot})^2=q(\theta)$, we obtain
$$\lambda(Z)=MN-q(\theta)+O(|\theta|^3)$$
The form $q$ is positive definite, because a matrix killed by the double centering is the sum of a function of $i$ and of a function of $j$, and so vanishes, its first row and column being zero.

Now let $V=\Lambda_Q^*\otimes\mathbb R$, which is the space of the real matrices which are orthogonal to $N_Q$. The map $\theta\to e^{i\theta}$ is a surjective morphism from $V$ onto the connected component $K_Q^\circ$ of the identity, with kernel $2\pi\Lambda_Q^*$, and the Haar measure of $K_Q$, restricted to $K_Q^\circ$, corresponds to the Lebesgue measure of $V/2\pi\Lambda_Q^*$, normalized as to have mass $\frac1\kappa$.

Fix a small ball $B\subset V$ around 0. Outside the image of $B$ we have $\lambda(Z)\leq MN-\varepsilon$ with $\varepsilon>0$, by compactness, and on the image of $B$ we have $MN-\lambda(Z)\leq\frac{MN}2$. By \eqref{eq:growth-trace}, the contributions to $MNc_n$ of the complement of the image of $B$, and of the second term of \eqref{eq:growth-trace} on the image of $B$, are $O((MN-\varepsilon)^n)$ and $O((MN/2)^n)$. Thus, up to exponentially small terms,
$$MNc_n\simeq\frac{(MN)^n}{\kappa\cdot vol(V/2\pi\Lambda_Q^*)}\int_B\exp\left(n\log\left(1-\frac{q(\theta)+O(|\theta|^3)}{MN}\right)\right)d\theta$$
and the Laplace method gives the equivalent
$$\frac{(MN)^n}{\kappa\cdot vol(V/2\pi\Lambda_Q^*)}\left(\frac{\pi MN}n\right)^{d/2}\cdot\frac1{\sqrt{\det(q|_V)}}=\frac{(MN)^n}{\kappa\sqrt{\Delta_Q}}\left(\frac{MN}{4\pi n}\right)^{d/2}$$
since $vol(V/2\pi\Lambda_Q^*)^2\det(q|_V)=(2\pi)^{2d}\Delta_Q$, for any choice of the Lebesgue measure on $V$. For generic $Q$ we have $\Lambda_Q^*=\mathbb Z^{(M-1)(N-1)}$ and $\kappa=1$, and $\Delta_Q$ is the determinant of the matrix in (4) above, which is $M^{-(N-1)}N^{-(M-1)}$, since $\det(1-\frac1M\mathbb I_{M-1})=\frac1M$.
\end{proof}

The above result is a local limit theorem for the random walk on the group $\Gamma_Q$, which contains $\mathbb Z^d$ as a subgroup of finite index. When $K_Q$ is finite we have $d=0$ and $\Delta_Q=1$, and we obtain $c_n\sim(MN)^n/|G_H|$, which corresponds to the fact that the atom of $\mu_H$ at $MN$ has mass $1/|G_H|$. In general, the theorem shows that the rooted law determines the number $d$, which is the corank of the lattice of relations $N_Q$, and also the number $\kappa\sqrt{\Delta_Q}$, which is a kind of volume. This completes Theorem \ref{thm:rigidity}, which shows that the rooted law determines the length of the shortest relation.

The program \texttt{zulle.py} computes $d,\kappa,\Delta_Q$ and $C$ from the exponents of the entries of $Q$, in exact arithmetic, and compares the result with the moments, computed exactly for $n$ up to 1600 when a closed formula is available, and by integration on a grid otherwise, as follows:
$$\begin{array}{|l|c|c|c|c|}\hline
{\rm matrix}&d&\kappa&\Delta_Q&C\\\hline
F_2\otimes_QF_N,\ {\rm generic}&N-1&1&2^{-(N-1)}N^{-1}&(N/\pi)^{(N-1)/2}/(2\sqrt N)\\\hline
F_3\otimes_QF_3,\ {\rm generic}&4&1&1/81&81/(16\pi^2)\\\hline
F_2\otimes_QF_4,\ {\rm row}\ (1,a,a,c),\ c\ {\rm of\ order}\ 4m&2&m&1/4&1/(2\pi m)\\\hline
F_2\otimes_QF_4,\ {\rm row}\ (1,q,1,q)&1&1&1/2&1/(4\sqrt\pi)\\\hline
F_2\otimes_QF_4,\ {\rm row}\ (1,q,i,q)&1&2&1/2&1/(8\sqrt\pi)\\\hline
F_2\otimes_QF_Y,\ Y=\mathbb Z_2^2,\ {\rm row}\ (1,p,q,pq)&2&1&1/4&1/(2\pi)\\\hline
F_3\otimes_QF_3,\ Q_{11}=q,\ Q_{22}=q^{-1}&2&1&1/3&\sqrt3/(4\pi)\\\hline
\end{array}$$
Here the parameters $a,p,q$ have infinite order, and $p,q$ are independent in the Klein case. At index 6 we obtain $c_n^{gen}\sim\frac{\sqrt3}{2\pi n}6^n$, in agreement with Theorem \ref{thm:density}, the density of $\mu_H$ at $x=6$ being $\frac{\sqrt3}{12\pi}$. The non-generic matrices in the above table are those studied below.

\subsection{The index 8}

For $Y=\mathbb Z_4$ we write the dephased second row as $(1,a,b,c)$ and we put
$$r=\frac{ac}b,\quad \alpha=\Big(\frac ca\Big)^2,\quad \beta=b^2$$
In all cases Corollary \ref{cor:two-rows-abelian} gives $\mu_H=\frac34\delta_0+\frac14 law(4+\varepsilon|S|)$ with $S=1+Z_1+Z_2+Z_3$. If $r$, $\alpha$ and $\beta$ are roots of unity, then so are $b$, $\frac ca$ and $ac$, hence $c^2=ac\cdot\frac ca$, and finally $a$, $b$ and $c$. In that case $K_Q$ is finite and given by Corollary \ref{cor:two-rows-abelian}. Thus, when looking for the infinite phase groups, we can assume that one of the numbers $r,\alpha,\beta$ has infinite order, or equivalently, that one of the numbers $a,b,c$ has infinite order.

In the generic case the coset graph is the diamond lattice, by Proposition \ref{prop:two-rows-graph}. In general, the situation is more complicated than at index 6, because the group $Y=\mathbb Z_4$ is not of prime order, and the phase group can be infinite without being the whole torus $\mathbb T^3$. In order to discuss this, let $U,V,W$ be independent random variables with Haar distribution on $\mathbb T$. We have the following result, which is the dual form of the trichotomy in \cite[Theorem 4.7]{BO26}, where the groups $N_Q$ are computed:

\begin{prop}[Cyclic phase groups]\label{prop:eight-cyclic}
Assume that $a$, $b$ or $c$ has infinite order. Then exactly one of the following cases occurs, and the random row $(1,Z_1,Z_2,Z_3)$ is given as follows.
\begin{enumerate}
\item If $r$ has infinite order and $\alpha$ or $\beta$ has infinite order, the row is $(1,U,V,W)$.
\item If $r$ is a root of unity, put $m=ord(r^4)$. The row is $(1,U,EUV,V)$, where $E$ is uniform on the $m$-th roots of unity and independent of $U,V$.
\item If $\alpha$ and $\beta$ are roots of unity, the row is $(1,UA,B,U)$, where $(A,B)$ is uniformly distributed on the finite group generated by $(\alpha\beta,\alpha\beta^{-1})$ and $(\alpha\beta^{-1},\alpha^{-1}\beta^{-1})$ and independent of $U$.
\end{enumerate}
\end{prop}

\begin{proof}
By the discussion before the proposition, $r$, $\alpha$ and $\beta$ are not all roots of unity. Thus the 3 cases are exhaustive, and they are mutually exclusive. The generators of Corollary \ref{cor:two-rows-abelian} are $g_0=1$ and
$$g_1=\Big(\frac\alpha{r^2},\frac\alpha\beta,\frac1{r^2\beta}\Big),\quad
g_2=\Big(\frac\alpha\beta,\frac1{\beta^2},\frac1{\alpha\beta}\Big),\quad
g_3=\Big(\frac1{r^2\beta},\frac1{\alpha\beta},\frac1{r^2\alpha}\Big)$$
where we used $a^2c^2=r^2\beta$ and $c^2a^{-2}=\alpha$, so that $b^2a^{-4}=\alpha r^{-2}$, $c^2a^{-2}b^{-2}=\alpha\beta^{-1}$ and $a^{-2}c^{-2}=r^{-2}\beta^{-1}$. For $l=(l_1,l_2,l_3)\in\mathbb Z^3$, put $P=l_1-l_3$ and $P'=l_1+2l_2+l_3$. A computation gives
\begin{equation}\label{eq:eight-relations}
g_2^l=\alpha^P\beta^{-P'},\quad
(g_1g_3^{-1})^l=\alpha^{P'}\beta^P,\quad
g_1^l=r^{-2(l_1+l_3)}\alpha^{l_1+l_2}\beta^{-l_2-l_3}
\end{equation}
where $g^l$ denotes $\prod_jg(j)^{l_j}$. The character $z\to z^l$ is trivial on $K_Q$ if and only if these 3 numbers are equal to 1.

Assume that $\alpha$ or $\beta$ has infinite order. If $(P,P')\neq(0,0)$, the first 2 relations give $\alpha^{P^2+P'^2}=(\alpha^P\beta^{-P'})^P(\alpha^{P'}\beta^P)^{P'}=1$ and $\beta^{P^2+P'^2}=(\alpha^P\beta^{-P'})^{-P'}(\alpha^{P'}\beta^P)^P=1$, which is impossible. Thus $P=P'=0$, that is $l=l_1(1,-1,1)$, and the third relation becomes $r^{-4l_1}=1$. If $r$ has infinite order, no nontrivial character is trivial on $K_Q$, hence $K_Q=\mathbb T^3$, which is case 1. If $r$ is a root of unity, the annihilator of $K_Q$ is generated by $m(1,-1,1)$, so that $K_Q=\{z:(z_1z_3z_2^{-1})^m=1\}$. The map $(u,v,\lambda)\to(u,\lambda uv,v)$ is an isomorphism from $\mathbb T^2\times\mu_m$ onto this group, where $\mu_m$ is the group of $m$-th roots of unity. This gives case 2.

Assume now that $\alpha$ and $\beta$ are roots of unity. Then $r$ has infinite order. Since $\alpha^{l_1+l_2}\beta^{-l_2-l_3}$ is a root of unity, the third relation forces $l_1+l_3=0$. Write $l=(h,k,-h)$. The character $z\to z^l$ is then $z\to(z_1z_3^{-1})^hz_2^k$, and by \eqref{eq:eight-relations} it is trivial on $K_Q$ if and only if
$$\alpha^{2h}\beta^{-2k}=1,\quad\alpha^{2k}\beta^{2h}=1,\quad\alpha^{h+k}\beta^{h-k}=1$$
The images of $g_1$, $g_2$ and $g_3$ under the homomorphism $z\to(z_1z_3^{-1},z_2)$ are
$$(\alpha\beta,\alpha\beta^{-1}),\quad(\alpha^2,\beta^{-2}),\quad(\alpha\beta^{-1},\alpha^{-1}\beta^{-1})$$
and the second one is the product of the other 2. Since $g_1g_3^{-1}$, $g_2$ and $g_1$ generate the same group as $g_1$, $g_2$ and $g_3$, the 3 displayed conditions mean that the character $(A,B)\to A^hB^k$ of $\mathbb T^2$ is trivial on the finite group $K'$ generated by these images. Hence the annihilator of $K_Q$ is the set of the vectors $(h,k,-h)$ such that $(h,k)$ annihilates $K'$. Taking annihilators again, we obtain $K_Q=\{z\in\mathbb T^3:(z_1z_3^{-1},z_2)\in K'\}$. The map $(u,A,B)\to(uA,B,u)$ is an isomorphism from $\mathbb T\times K'$ onto $K_Q$, which gives case 3.
\end{proof}

In \eqref{eq:eight-relations}, the exponent vectors $(P,-P')$ and $(P',P)$ of $(\alpha,\beta)$ are the rows of the matrix of the multiplication by the Gaussian integer $P+i P'$ on $\mathbb Z[i]\simeq\mathbb Z^2$. This reflects the decomposition of $I_Y\otimes\mathbb Q$ according to the characters $y\to\pm i$ and $y\to-1$ of $Y$, used in the proof of \cite[Theorem 4.7]{BO26}.

By Theorem \ref{thm:growth}, in the 3 cases of Proposition \ref{prop:eight-cyclic} the moments $c_n$ grow as $8^nn^{-3/2}$, as $8^nn^{-1}$, and as $8^nn^{-1/2}$, the number of connected components of the phase group being $1$, $m$, and the order of the finite group in case 3. In case 2 the law differs from the generic one, but the difference is seen quite late:

\begin{prop}[The first exceptional moment]\label{prop:eight-first-moment}
We have
$$c_n=c_n^{gen}\quad{\rm for}\ 1\leq n<4m,
\quad
c_{4m}-c_{4m}^{gen}=\frac12\binom{2m}m^2$$
in case 2 of Proposition \ref{prop:eight-cyclic}.
\end{prop}

\begin{proof}
In $|S|^{2j}$, each monomial comes from 2 words of length $j$ in $1,U,EUV,V$. Write $\kappa,\kappa'$ for their contents and $\delta=\kappa-\kappa'$, so that $\sum_i\delta_i=0$. Haar integration keeps the monomial exactly when $\delta_1+\delta_2=\delta_2+\delta_3=0$ and $m$ divides $\delta_2$. Together these conditions mean
$$\delta\in m\mathbb Z(1,-1,1,-1)$$
whereas for generic parameters the expectation keeps exactly the pairs with $\delta=0$.

The positive part and the negative part of $\delta=\pm m(1,-1,1,-1)$ have total size $2m$, so the first additional pairs occur for $j=2m$. In one orientation, the first word contains $m$ copies of $1$ and $m$ copies of $EUV$, and the second word contains $m$ copies of $U$ and $m$ copies of $V$. There are $\binom{2m}m^2$ such pairs, and the same number in the other orientation.

Finally, in $\mathbb E(4+\varepsilon|S|)^n$ only the even powers of $|S|$ survive, and $|S|^{4m}$ first occurs for $n=4m$, with coefficient one. The weight $\frac14$ of the second term of $\mu_H$ gives the result.
\end{proof}

Every $m\geq1$ occurs, for instance with $a=b$ of infinite order and $c$ a primitive $4m$-th root of unity. For $m=1$ we have $S=(1+U)(1+V)$, so that $\mathbb E|S|^{2j}=\binom{2j}j^2$ by taking the constant terms in $U$ and in $V$. This gives $c_4=169$ instead of the generic value $167$, in agreement with \cite[Section 4.3]{BO26}. The program \texttt{wordgram.py} checks the formula of Proposition \ref{prop:eight-first-moment} for $m=1,2,3,4$.

\begin{example}[An atom at the middle of the spectrum]
Let $q$ have infinite order. For the row $(1,q,1,q)$ we have $\alpha=\beta=1$, so that $K'$ is trivial in case 3 of Proposition \ref{prop:eight-cyclic}. The random row is $(1,U,1,U)$ and $S=2+2U$. With $t$ uniform on $[0,2\pi]$, the rooted law is
$$\frac34\delta_0+\frac14 law(4+4\cos t)$$
whose second moment is $6$, whereas the generic second moment is $5$. For the row $(1,q,i,q)$ we have $\alpha=1$ and $\beta=-1$, so that $(A,B)$ is equally likely to be $(1,1)$ or $(-1,-1)$. The corresponding sums are $2+2U$ and $0$, and the rooted law is
$$\frac34\delta_0+\frac18\delta_4+\frac18 law(4+4\cos t)$$
\end{example}

Here is now an example with $Y=\mathbb Z_2^2$ instead of $\mathbb Z_4$, which is never generic:

\begin{example}[A Klein four deformation]
Take $Y=\mathbb Z_2^2$ and the dephased row $(1,p,q,pq)$ in the order $00,10,01,11$. Corollary \ref{cor:two-rows-abelian} gives
$$g_{10}=(1,p^{-4},1,p^{-4}),\quad
g_{01}=(1,1,q^{-4},q^{-4}),\quad g_{11}=g_{10}g_{01}$$
Hence the random row is $(1,U,V,UV)$, where $U$ and $V$ are independent and Haar distributed on $\langle p^4\rangle$ and on $\langle q^4\rangle$. Its sum is $(1+U)(1+V)$. Even when $p=q$ has infinite order, $U$ and $V$ are independent uniform circle variables, and the fourth moment is again $169$. In particular this family is never generic, since $Z_{11}=Z_{10}Z_{01}$ on $K_Q$.

More precisely, with $u_j=\mathbb E|1+U|^{2j}$ and $v_j=\mathbb E|1+V|^{2j}$, Corollary \ref{cor:two-rows-abelian} gives
$$c_2=\frac{16+u_1v_1}4,\quad c_3=\frac{64+12u_1v_1}4,\quad c_4=\frac{256+96u_1v_1+u_2v_2}4$$
Here $u_1=2$ if $p^4\neq1$, and $u_1=4$ otherwise. In the first case $\mathbb E(U)=0$ and $u_2=6+2Re(\mathbb E(U^2))\geq6$, since $\mathbb E(U^2)\in\{0,1\}$; similarly for $V$. Thus $p^4,q^4\neq1$ give $c_2=5$, $c_3=28$, $c_4\geq169$, with the first nongeneric moment at $n=4$. Exactly one of $p^4,q^4$ equal to 1 gives $c_2=6$, and both give $c_2=8$. The script \texttt{klein\_moments.py} checks the first 4 moments for finite cyclic and circle choices.
\end{example}

\subsection{The index 9}

We now take $X=Y=\mathbb Z_3$. On one exceptional family the moments $1,5,31,213$ were computed in \cite[Section 4.3]{BO26}, compared with $1,5,29,181$ at a generic parameter. We compute here the whole measure, on the slice
$$Q
=\begin{pmatrix}1&1&1\\1&q&1\\1&1&q^{-1}\end{pmatrix},
\quad q\in\mathbb T$$
The squared length of the 3-step walk of index 6 appears again.

\begin{thm}[The squared 3-step walk]
For $Q$ as above, put
$$P_q
=\{(u,v)\in\mathbb T^2:vu^{-2}\in\langle q^3\rangle,\ uv^{-2}\in\langle q^3\rangle\}$$
and let $(U,V)$ be Haar distributed on $P_q$. Then
$$\mu_H=\frac23\delta_0+\frac13 law(|1+U+V|^2)$$
If $q^3$ has finite order $l$, then $P_q$ has $3l^2$ elements. If $q$ has infinite order, then $P_q=\mathbb T^2$, we have $\mu_H=\frac23\delta_0+\frac13 law(3+A)$, where $A$ is the adjacency operator of the triangular lattice, and for any integer $n\geq1$,
$$c_n=\frac13W_3(2n)=\frac13\sum_{a+b+c=n}\binom n{a,b,c}^2$$
where $a,b,c\geq0$.
\end{thm}

\begin{proof}
We first compute $K_Q$. Write $Q_{ij}=q^{\nu_{ij}}$ with $\nu_{11}=1$, $\nu_{22}=-1$ and $\nu_{ij}=0$ otherwise. A direct computation from \eqref{eq:gram-generators} shows that $g_{k,e}=1$ if $k=0$ or $e=0$, and that, in the coordinates $(Z_{11},Z_{12},Z_{21},Z_{22})$, the exponents of $q$ in the remaining generators are
$$g_{1,1},g_{2,2}:3(-1,0,0,1),\quad g_{1,2}:3(0,1,1,1),\quad g_{2,1}:-3(1,1,1,0)$$
The first vector is the sum of the other 2. Hence $K_Q$ is the closed subgroup of the matrices
$$Z=\begin{pmatrix}1&1&1\\1&A&AB\\1&AB&B\end{pmatrix},
\quad A,B\in\langle q^3\rangle$$
with $(A,B)$ Haar distributed on $\langle q^3\rangle^2$.

Consider the phase matrix
$$R=\begin{pmatrix}1&U&V\\U&V&1\\V&1&U\end{pmatrix}$$
Dividing its second and third columns by $U$ and $V$, and then its second and third rows by $U$ and $V$, gives the matrix $Z$ above with $A=VU^{-2}$ and $B=UV^{-2}$, since $AB=(UV)^{-1}$. The homomorphism $\psi(u,v)=(vu^{-2},uv^{-2})$ of $\mathbb T^2$ has exponent matrix of determinant $3$. Hence it is surjective, with kernel $\{(\omega^j,\omega^{2j}):j=0,1,2\}$, where $\omega=e^{2\pi i/3}$. By definition $P_q=\psi^{-1}(\langle q^3\rangle^2)$, so that $\psi$ restricts to a surjective homomorphism $P_q\to\langle q^3\rangle^2$, which carries the Haar measure of $P_q$ to the Haar measure of $\langle q^3\rangle^2$.

Multiplying rows and columns by phases does not change the eigenvalues of $RR^*$. Exchanging the last 2 rows of $R$ gives the circulant matrix with first row $(1,U,V)$, which is normal with eigenvalues $1+\omega^jU+\omega^{2j}V$, $j=0,1,2$. Thus the eigenvalues of $ZZ^*$ are
$$|1+\omega^jU+\omega^{2j}V|^2,\quad j=0,1,2$$
The translation by $(\omega,\omega^2)$ preserves the Haar measure of $P_q$, so the 3 eigenvalues have the same law. Theorem \ref{thm:gram-all} gives the formula for $\mu_H$. If $q^3$ has order $l$, then $\langle q^3\rangle^2$ has $l^2$ elements and each fiber of $\psi$ has 3 elements. If $q$ has infinite order, then $\langle q^3\rangle=\mathbb T$ and $P_q=\mathbb T^2$, so that $U$ and $V$ are independent and uniform on $\mathbb T$. We have
$$|1+U+V|^2=3+U+\bar U+V+\bar V+U\bar V+\bar UV$$
and the last 6 terms are the Fourier transform of the adjacency operator of the triangular lattice, which gives the formula with $A$. Finally, expand $(1+U+V)^n$ and its conjugate. The expectation keeps the pairs of words with the same numbers of letters $U$ and $V$, hence with the same content $(a,b,c)$, and each content contributes the square of its multinomial coefficient.
\end{proof}

The first 7 values are $c_n=1,5,31,213,1551,11723,90945$, and on this family the moments $1,5,31,213$ were found in \cite{BO26}. These 7 values are confirmed by the program \texttt{wordgram.py}. If $q$ has infinite order, the shortest nontrivial element of $N_Q$ has length $3$ by Theorem \ref{thm:rigidity}, since $c_2=5$ is the generic value and $c_3=31$ exceeds the generic value $29$.

Also, Theorem \ref{thm:growth} gives $c_n\sim\frac{\sqrt3}{4\pi n}9^n$, to be compared with the generic growth, which is in $9^nn^{-2}$. Observe that the triangular lattice is the lattice of the second neighbors on the honeycomb lattice, in agreement with the fact that the 3-step walk of the index 6 appears here through its squared length. The 3 lattices which appear at indices 6 and 9, namely the honeycomb, kagome and triangular lattices, are thus closely related.

\section{Circulant families and Jacobi parameters}\label{sec:circulant-jacobi}

The circulant matrix used at index 9 extends to every square index. We construct a family with one parameter whose phase group has dimension $N-1$, and whose law is the squared length of an $N$-step planar walk. We then express the first-relation theorem in the Hankel and Jacobi language of \cite{BP26}, applying it both to this family and to the delayed exceptional moments at index 8.

\subsection{A circulant family at square index}

Let $G$ be a finite abelian group of order $N\geq2$, written additively, and let $q\in\mathbb T$. Consider
$$Q^{(q)}_{ij}=q^{1+\delta_{i+j,0}-\delta_{i,0}-\delta_{j,0}},\quad H_q=F_G\otimes_{Q^{(q)}}F_G$$
The first row and column of $Q^{(q)}$ are 1. Its other entries are $q^2$ when $i+j=0$, and $q$ otherwise. This is the dephasing of the matrix $C_{ij}=q^{\delta_{i+j,0}}$.

\begin{thm}\label{thm:circulant-family}
If $q$ has infinite order, the phase group of $H_q$ is the connected torus
$$K_{Q^{(q)}}=B_G=\{D(z):z\in\mathbb T^G\},\quad D(z)_{ij}=\frac{z_{i+j}z_0}{z_iz_j},\quad \dim B_G=N-1$$
The rooted law is
$$\mu_{H_q}=\left(1-\frac1N\right)\delta_0+\frac1N law(|1+U_1+\cdots+U_{N-1}|^2)$$
where the $U_j$ are independent uniform circle variables. In particular this law depends only on $N$, and not on $q$ or on the abelian group structure of $G$.
\end{thm}
\begin{proof}
The map $D$ is a homomorphism for coordinatewise multiplication. Its kernel consists of the vectors $z_i=a\chi(i)$, with $a\in\mathbb T$ and $\chi\in\widehat G$: indeed $D(z)=1$ says that $i\to z_i/z_0$ is a character. Thus $B_G$ is a connected torus of dimension $N-1$.

We compute the phase generators on the undephased matrix $C$, as permitted by Proposition \ref{prop:gram-generators}. Let $a_t$, $t\in G$, denote the coordinate vectors in $\mathbb R^G$, and put
$$v_{k,e}=a_{-e}+a_k-a_{k-e}-a_0$$
Substitution in the generator formula gives $g_{k,e}=D(q^{v_{k,e}})$, where exponentiation is coordinatewise. Hence $K_{Q^{(q)}}\subset B_G$. To prove the reverse inclusion, sum over $k$:
$$\sum_{k\in G}v_{k,e}=N(a_{-e}-a_0)$$
Therefore the products of the generators include $D(q^{N(a_t-a_0)})$ for every $t\in G$. Since $q^N$ has infinite order, the closure of the subgroup generated by the vectors $q^{N(a_t-a_0)}$, with $t\neq0$, is
$$S_G=\{z\in\mathbb T^G:\prod_{t\in G}z_t=1\}$$
Each coordinate other than the zeroth one can be varied independently along a dense circle, and the zeroth coordinate is the inverse of their product. Moreover $D(S_G)=B_G$, because multiplying any $z$ by a common phase can make its product equal to 1, without changing $D(z)$. This proves the phase-group assertion.

Push Haar measure on $\mathbb T^G$ through $D$ to obtain Haar measure on $B_G$. The matrices $D(z)$ and $R(z)=(z_{i+j})_{ij}$ differ by diagonal unitary row and column factors, so their Gram eigenvalues agree. Negating the row index turns $R(z)$ into the convolution matrix $(z_{j-i})_{ij}$. Its Fourier eigenvectors give Gram eigenvalues
$$\left|\sum_{t\in G}z_t\chi(t)\right|^2,\quad \chi\in\widehat G$$
For independent Haar coordinates $z_t$, all these eigenvalues have the same distribution. Theorem \ref{thm:gram-all} therefore gives the stated law, with $N$ independent phases in the sum. Dividing them by $z_0$ leaves the modulus unchanged and makes the first phase equal to 1, giving the displayed version.
\end{proof}

For $G=\mathbb Z_4$, with the ordering $0,1,2,3$, the parameter matrix is
$$Q^{(q)}=\begin{pmatrix}1&1&1&1\\1&q&q&q^2\\1&q&q^2&q\\1&q^2&q&q\end{pmatrix}$$
It gives the law $\frac34\delta_0+\frac14 law(|1+U+V+W|^2)$ at index 16. The same law is obtained with $G=\mathbb Z_2^2$. At $N=3$ we recover the law of the infinite-order slice in Section \ref{sec:growth}, through a different parameter matrix. These assertions concern the laws, without asserting equivalence of the Hadamard matrices or isomorphism of their quantum groups.

The scalar law has a lattice interpretation, which extends the one found at index 9. For independent uniform phases $z_t$, with $t\in G$, we have
$$\Big|\sum_{t\in G}z_t\Big|^2=N+\sum_{s\neq t}z_s\bar z_t$$
and the last $N(N-1)$ terms are the Fourier transform of the adjacency operator $A$ of the following graph $\Lambda_G$. The vertices are the vectors $m\in\mathbb Z^G$ with $\sum_tm_t=0$, and $m$ is joined to the vectors $m+a_s-a_t$, with $s\neq t$, where $a_t$ are the coordinate vectors, as before. This is the root lattice of type $A_{N-1}$, with the roots as steps. Thus
$$\mu_{H_q}=\left(1-\frac1N\right)\delta_0+\frac1N law(N+A)$$
with the law of $A$ being taken at any vertex. At $N=3$ we recover the triangular lattice of Section \ref{sec:growth}, and at $N=4$ we obtain the face-centered cubic lattice, where each vertex has 12 neighbors. Observe also that $\Lambda_G$ is the lattice of the second neighbors on the generic coset graph of Proposition \ref{prop:two-rows-graph}, which is the honeycomb lattice at $N=3$, and the diamond lattice at $N=4$.

\subsection{Moments and growth}

The scalar formula makes every moment a finite combinatorial sum. It also identifies the first relation which separates this family from the fully generic parameters.

\begin{prop}\label{prop:circulant-moments}
For the family of Theorem \ref{thm:circulant-family}, we have $c_0=1$ and
$$c_n=\frac1NW_N(2n)=\frac1N\sum_{a_1+\cdots+a_N=n}\binom n{a_1,\ldots,a_N}^2$$
for $n\geq1$, with $a_j\geq0$. In particular,
$$c_1=1,\quad c_2=2N-1,\quad c_3=6N^2-9N+4$$
For $N\geq3$ the first nongeneric moment is $c_3$, with defect $(N-1)(N-2)$, and the shortest nontrivial relation has length 3. At $N=2$ the law is generic.
\end{prop}
\begin{proof}
Expand the $n$-th power of the sum of $N$ independent phases and its conjugate. Integration keeps exactly the pairs of words with the same content $(a_1,\ldots,a_N)$, giving the squared multinomial sum. At degree 2 the contents have types $(2)$ and $(1,1)$, and at degree 3 they have types $(3)$, $(2,1)$ and $(1,1,1)$. The sums are
$$N+4\binom N2=N(2N-1),\quad N+9N(N-1)+36\binom N3=N(6N^2-9N+4)$$
Division by $N$ gives the first moments.

For generic square parameters, the Gram model instead has independent matrix entries. A term in the third Gram trace survives when each entry phase has balanced exponents. Either all row indices coincide, all column indices coincide, or there are exactly 2 distinct rows and 2 distinct columns. The first 2 cases contribute $N^4$ each, with an overlap of $N^2$ terms. In the last case the singleton row has 3 possible positions, and the singleton column must occupy the opposite position, giving $3N^2(N-1)^2$ terms. After division by $N^2$ we obtain
$$c_1^{gen}=1,\quad c_2^{gen}=2N-1,\quad c_3^{gen}=2N^2-1+3(N-1)^2=5N^2-6N+2$$
Thus the first defect is $(N-1)(N-2)$ when $N\geq3$, and Theorem \ref{thm:rigidity} identifies its relation length. At $N=2$, the torus $B_G$ has the full dimension $(N-1)^2=1$ of the dephased phase torus, so the law is generic.
\end{proof}

By Proposition \ref{prop:two-rows-graph}, the number $Nc_n$ counts the closed walks of length $2n$ on the generic coset graph at index $2N$. In other words, if $|S|$ denotes the length of the $N$-step planar walk, the generic law at index $2N$ involves the variable $N+\varepsilon|S|$, while the present law at index $N^2$ involves $|S|^2$. In terms of the density $p_N$ of $|S|$, studied in \cite{BSWZ12}, the continuous part of $\mu_{H_q}$ is $\frac1{2N\sqrt x}p_N(\sqrt x)dx$, on $[0,N^2]$.

At index 16 we obtain the Domb numbers divided by 4, namely $c_n=1,7,64,679,7876,96784$ for $1\leq n\leq6$. The program \texttt{coppi.py} checks the phase-generator identities on cyclic and noncyclic groups, and independently counts return words using the block action. These counts agree with the multinomial formula through degree 6 for $N=2,3$, through degree 5 for both groups of order 4, and through degree 3 for $N=5$. It also checks the generic third moments directly by balancing Gram monomials.

The phase-group dimension $N-1$ is smaller than the generic value $(N-1)^2$ when $N\geq3$. We can compute the resulting growth constant directly from the scalar law.

\begin{cor}\label{cor:circulant-growth}
For the circulant family with $q$ of infinite order, we have
$$c_n\sim\frac{N^{(N-2)/2}}{2^{N-1}\pi^{(N-1)/2}}(N^2)^nn^{-(N-1)/2}$$
as $n\to\infty$. Thus the power of $n$ agrees with Theorem \ref{thm:growth} at $d=N-1$.
\end{cor}
\begin{proof}
Fix the zeroth phase to be 1 and write the others as $e^{i\theta_j}$. The unique maximum of the squared modulus, modulo $2\pi$, occurs when all the phases are 1. Near this maximum,
$$\left|1+\sum_{j=1}^{N-1}e^{i\theta_j}\right|^2=N^2-\theta^t(NId-\mathbb I_{N-1})\theta+O(|\theta|^4)$$
The quadratic form has eigenvalue 1 on the all-one vector and eigenvalue $N$ on its orthogonal complement, so its determinant is $N^{N-2}$. Outside any fixed neighborhood of the maximum, the contribution is exponentially smaller than $(N^2)^n$. Laplace integration with Haar measure $d\theta/(2\pi)^{N-1}$ therefore gives
$$c_n\sim\frac{(N^2)^n}{N(2\pi)^{N-1}}\left(\frac{\pi N^2}n\right)^{(N-1)/2}\frac1{N^{(N-2)/2}}$$
which simplifies to the stated constant. At $N=3$ it is $\sqrt3/(4\pi)$, agreeing with the index-9 computation, and at $N=4$ it is $1/(2\pi^{3/2})$.
\end{proof}

Comparing with Theorem \ref{thm:growth}, where $\kappa=1$ since $B_G$ is connected, we obtain $\Delta_Q=N^{N-4}$. The program \texttt{zulle.py} confirms this, from the exponents of the entries of $Q^{(q)}$, for $\mathbb Z_3,\mathbb Z_4,\mathbb Z_2^2,\mathbb Z_5,\mathbb Z_6$.

\subsection{The first Jacobi defect}

The invariant-dimension viewpoint of \cite[Sections 10 and 11]{BP26} suggests expressing our moment comparisons through orthogonal polynomials. For a probability law with moments $c_n$, let $p_j$ be its monic orthogonal polynomials and put
$$h_j=\int p_j(x)^2d\mu(x),\quad D_k=\det(c_{i+j})_{0\leq i,j<k},\quad D_0=1$$
As long as the relevant norms are positive, the monic recurrence and its coefficients are
$$xp_j=p_{j+1}+\beta_jp_j+\lambda_jp_{j-1},\quad \beta_j=\frac1{h_j}\int xp_j(x)^2d\mu(x),\quad \lambda_j=\frac{h_j}{h_{j-1}}$$
Here the last formula is for $j\geq1$, while $p_0=1$, $p_{-1}=0$ and $\lambda_0=0$. Also $h_j=D_{j+1}/D_j$. In the symmetric Jacobi matrix the off-diagonal entries are $\sqrt{\lambda_j}$; we use the monic coefficients $\lambda_j$ throughout.

\begin{prop}\label{prop:jacobi-defect}
Let 2 compactly supported probability laws first differ in moment degree $L\geq2$, with $\delta=c_L-c_L^0>0$, and suppose the reference law has infinite support. Their first changed Jacobi coefficient is
$$\lambda_r-\lambda_r^0=\frac\delta{h_{r-1}^0}\quad(L=2r),\quad \beta_r-\beta_r^0=\frac\delta{h_r^0}\quad(L=2r+1)$$
In the order $\beta_0,\lambda_1,\beta_1,\lambda_2,\ldots$, all preceding coefficients agree. In the even case we also have
$$D_k=D_k^0\quad(k\leq r),\quad D_{r+1}-D_{r+1}^0=\delta D_r^0$$
so the first changed leading Hankel determinant is determined by the same defect.
\end{prop}
\begin{proof}
The moments through degree $2r-1$ determine $p_r$, since its coefficients solve the orthogonality equations against $1,x,\ldots,x^{r-1}$. Thus, for $L=2r$, both laws have the same monic $p_r$, and the difference of its squared norms is $\delta$, because $p_r^2$ is monic of degree $2r$. All lower norms agree. We obtain $\lambda_r-\lambda_r^0=\delta/h_{r-1}^0$, and the new norm is positive since $h_r=h_r^0+\delta>0$. In the Hankel matrix of size $r+1$, only the bottom-right entry changes, by $\delta$; its cofactor is $D_r^0$.

For $L=2r+1$, the polynomial $p_r$ and its norm $h_r$ both agree with the reference ones. The polynomial $xp_r^2$ is monic of degree $2r+1$, so its integral changes by $\delta$. Division by $h_r^0$ gives the second formula. The preceding recurrence coefficients use only lower moments, proving the remaining assertion.
\end{proof}

For a nongeneric Fourier deformation, Theorem \ref{thm:rigidity} makes $L$ the shortest nontrivial relation length and makes $\delta$ positive. The generic law has infinite support, as follows for instance from Theorem \ref{thm:growth} with $d=(M-1)(N-1)>0$. The proposition therefore gives a Jacobi form of the first-relation theorem: an even length first changes a coefficient $\lambda_r$, and an odd length first changes a coefficient $\beta_r$.

For the circulant family, Proposition \ref{prop:circulant-moments} gives $L=3$ for $N\geq3$. Since $p_1=x-1$, we obtain
$$\beta_0=1,\quad \lambda_1=2(N-1),\quad \beta_1=\frac{6N-7}2,\quad \beta_1^{gen}=\frac{5(N-1)}2$$
In particular $\beta_1-\beta_1^{gen}=(N-2)/2$, so the first extra relations are visible at the first nontrivial diagonal coefficient.

The delayed index-8 family exhibits the even case. It also shows how far out in the Jacobi recurrence the first difference can occur.

\begin{prop}\label{prop:delayed-hankel}
For the family of Proposition \ref{prop:eight-first-moment}, with $m$ connected components, we have
$$D_k=D_k^{gen}\quad(k\leq2m),\quad D_{2m+1}-D_{2m+1}^{gen}=\frac12\binom{2m}m^2D_{2m}^{gen}$$
The first changed Jacobi coefficient satisfies
$$\lambda_{2m}-\lambda_{2m}^{gen}=\frac12\binom{2m}m^2\frac{D_{2m-1}^{gen}}{D_{2m}^{gen}}$$
while $\beta_j=\beta_j^{gen}$ for $j\leq2m-1$ and $\lambda_j=\lambda_j^{gen}$ for $j<2m$.
\end{prop}
\begin{proof}
The first moment defect is at degree $4m$, with value $\frac12\binom{2m}m^2$, by Proposition \ref{prop:eight-first-moment}. Apply Proposition \ref{prop:jacobi-defect} with $r=2m$ and use $h_{2m-1}^{gen}=D_{2m}^{gen}/D_{2m-1}^{gen}$.
\end{proof}

For $m=1$, the moments through degree 3 are $1,1,5,28$ in both cases, while the fourth moments are 167 and 169. Hence
$$D_3^{gen}=39,\quad D_3=47,\quad \lambda_2^{gen}=\frac{39}{16},\quad \lambda_2=\frac{47}{16}$$
The program \texttt{coppi.py} verifies the determinant and Jacobi formulas for $m=1,\ldots,5$, using exact character counts on $\mathbb T^2$ times the group of $m$-th roots of unity. Since $m$ can be arbitrarily large, no fixed finite initial segment of moments, Hankel determinants or Jacobi parameters detects genericity uniformly at index 8. Here the comparison is between genuine positive subfactor laws; the long positive truncations without a full law in \cite{BP26} concern a different obstruction.

\subsection{Finite-depth laws with long generic truncations}

The same question can be asked at every product index, with finite-depth laws in place of the index-8 family. For this subsection let $X,Y$ have orders $M,N\geq2$, and put
$$Q_{ij}=q^{(1-\delta_{i,0})(1-\delta_{j,0})},\quad H_q=F_X\otimes_QF_Y,\quad d=(M-1)(N-1)$$
Thus the first row and column of $Q$ are 1 and every other entry is $q$. Here $d$ denotes the dimension of the generic phase torus.

\begin{thm}\label{thm:finite-prefix}
Suppose that $q$ has order $l\geq2$ and $\gcd(l,MN)=1$. The phase group is
$$K_Q=\{Z\in\mathbb T^{X\times Y}:Z_{0j}=Z_{i0}=1,\ Z_{ij}^l=1\},\quad |K_Q|=l^d$$
The subfactor has finite depth, and its moments satisfy
$$c_n=c_n^{gen}\quad(0\leq n<2l),\quad c_{2l}-c_{2l}^{gen}=d$$
with $c_n^{gen}$ as in Theorem \ref{thm:gram}. In particular the first extra relation has length $2l$.
\end{thm}
\begin{proof}
The matrix $Q$ is the dephasing of $C_{ij}=q^{\delta_{i,0}\delta_{j,0}}$. Substitution in Proposition \ref{prop:gram-generators} gives, for nonzero $i,j,k,e$,
$$g_{k,e}(i,j)=q^{-(1+\delta_{i,k})(1+\delta_{j,-e})}$$
The generators with $k=0$ or $e=0$ are trivial. After reindexing $e$ by $-e$, their exponent matrix is $-A\otimes B$, where
$$A=Id+\mathbb I_{M-1},\quad B=Id+\mathbb I_{N-1},\quad |\det(-A\otimes B)|=M^{N-1}N^{M-1}$$
Here $\mathbb I$ denotes the all-one matrix. Indeed $A$ has eigenvalues $M,1,\ldots,1$, and similarly for $B$. The determinant is invertible modulo $l$, so the generators fill the entire grid of $l$-th roots in the $d$ free coordinates. Explicitly, we have
$$A^{-1}=Id-\frac1M\mathbb I_{M-1},\quad B^{-1}=Id-\frac1N\mathbb I_{N-1}$$
These formulas remain valid modulo $l$, with $M^{-1},N^{-1}$ denoting their modular inverses. They allow any prescribed vector of free phase exponents to be expressed as a combination of the generator vectors.

The dual group $T_Q$ therefore has $l^d$ elements, and $\Gamma_Q=T_Q\rtimes Y$ has $Nl^d$ elements. The quotient description of Theorem \ref{thm:dita-all} gives $\dim C(G_{H_q})=MNl^d$. There are only finitely many irreducible representations, so the associated Hadamard subfactor has finite depth.

Multiplying rows and columns by phases preserves Gram spectra. Hence we may compute the finite-grid integral using independent uniform $l$-th roots in all $MN$ positions: dephasing this matrix pushes its uniform measure to Haar measure on $K_Q$. With $i_a\in X$, $j_a\in Y$ and $i_{n+1}=i_1$, expand
$$Tr((ZZ^*)^n)=\sum_{i_1,\ldots,i_n}\sum_{j_1,\ldots,j_n}\prod_{a=1}^nZ_{i_aj_a}\bar Z_{i_{a+1}j_a}$$
Let $E=(E_{ij})$ be the exponent matrix of a term. Every row sum and column sum of $E$ is zero, and the sum of its positive entries is at most $n$. Generic integration keeps exactly $E=0$, while finite-grid integration keeps exactly the matrices whose entries are divisible by $l$.

A nonzero such matrix has positive mass at least $2l$. To see this, choose a positive entry. A negative entry in the same row forces a second positive entry in another row, since its column sum is zero. Both positive entries are at least $l$. Thus no extra term occurs for $n<2l$.

At $n=2l$, equality forces exactly 2 positive and 2 negative entries, equal to $l$ and $-l$. The zero row and column sums place them at the corners of a rectangle. There is no cancellation between the positive and negative occurrences in the monomial, since the positive mass has already reached $n$. For distinct rows $i,i'$ and columns $j,j'$, the extra monomials are therefore
$$\left(Z_{ij}\bar Z_{i'j}Z_{i'j'}\bar Z_{ij'}\right)^l$$
and their conjugates. Once the initial row and column have been chosen, the next row must be the other row of the rectangle, and the next column must be the other column. Repeating this gives the whole cyclic sequence. Thus there are exactly 4 terms for each unordered pair of rows and unordered pair of columns. No extra factor $l$ appears: shifting by a full period repeats the same index sequence in the trace expansion. Using the normalization in Theorem \ref{thm:gram-all}, we obtain
$$c_{2l}-c_{2l}^{gen}=\frac4{MN}\binom M2\binom N2=(M-1)(N-1)$$
The assertion about relation length now follows from Theorem \ref{thm:rigidity}.
\end{proof}

The determinant calculation also shows that an infinite-order $q$ in this same family gives the full generic phase torus. Indeed its exponent lattice has finite index in $\mathbb Z^d$, and the powers of $q$ along any nonzero integer multiple of a coordinate axis are dense in the circle. The restriction on $l$ in the theorem makes the finite case just as explicit.

In the covering picture of Theorem \ref{thm:covering}, the finite case corresponds to $N_Q=lT$, which is the only subgroup of $T$ with quotient isomorphic to $\mathbb Z_l^d$. Thus $B_Q$ is the covering of $K_{X,Y}$ associated to its homology with coefficients in $\mathbb Z_l$, and the extra closed walks found in the above proof are the 4-cycles of $K_{X,Y}$, traversed $l$ times. There are $M(M-1)(N-1)$ such walks based at a given white vertex, in agreement with the formula $c_n=w_n/M$ of Proposition \ref{prop:covering-walks}. As for the assumption on $l$, this cannot be dropped. At $M=N=2$ the above determinant is 4, we recover the group $K_Q=\langle q^4\rangle$ found in Section \ref{sec:lattices}, and for $q$ of order 2 or 4 the law is the Fourier one. The program \texttt{zulle.py} checks the order $l^d$ of the group $K_Q$ in 7 cases, and this latter phenomenon as well.

\begin{cor}\label{cor:finite-jacobi}
For the family of Theorem \ref{thm:finite-prefix}, the first changed Jacobi coefficient and leading Hankel determinant satisfy
$$\lambda_l-\lambda_l^{gen}=\frac d{h_{l-1}^{gen}},\quad D_{l+1}-D_{l+1}^{gen}=dD_l^{gen}$$
All preceding Jacobi coefficients agree, and $D_k=D_k^{gen}$ for $k\leq l$. As $l\to\infty$ through orders coprime to $MN$, these finite-depth laws converge weakly to the generic law. At fixed $X,Y$, no finite initial segment of moments or Jacobi coefficients distinguishes finite depth from the generic infinite-depth case.
\end{cor}
\begin{proof}
Apply Proposition \ref{prop:jacobi-defect} with $L=2l$ and $\delta=d$. All the laws are supported on $[0,MN]$, and each fixed moment is eventually equal to its generic value. Polynomial approximation of continuous functions on this interval gives weak convergence. The generic law has infinite support by Theorem \ref{thm:growth}, so its subfactor has infinite depth. Integers $l$ coprime to $MN$ can be arbitrarily large, proving the last assertion.
\end{proof}

For example, take $X=Y=\mathbb Z_3$ and $q=-1$, so that $l=2$ and $d=4$. We obtain
$$Q=\begin{pmatrix}1&1&1\\1&-1&-1\\1&-1&-1\end{pmatrix},\quad(c_1,c_2,c_3,c_4)=(1,5,29,185)$$
Here $|K_Q|=16$ and $\dim C(G_{H_q})=144$. The generic fourth moment is 181, and the corresponding Hankel and Jacobi values are
$$D_3^{gen}=48,\quad D_3=64,\quad\lambda_2^{gen}=3,\quad\lambda_2=4$$
The program \texttt{coppi.py} checks the generator matrix on cyclic and noncyclic groups, the first moment and Jacobi defects for $M=2$, $2\leq N\leq5$ and coprime orders $2\leq l\leq7$, and 3-row examples by direct Gram counts. Independent block-action counts also verify the finite-grid laws, including the displayed index-9 example.

\begin{thebibliography}{99}
\hypersetup{urlcolor=black}
\bibitem{AH99}
M. Asaeda and U. Haagerup, Exotic subfactors of finite depth with Jones indices $(5+\sqrt{13})/2$ and $(5+\sqrt{17})/2$, {\em Comm. Math. Phys.} {\bf 202} (1999), 1--63.

\bibitem{ABF26}
Astra, Banica and Fable, Karlsson matrices and the kagome lattice, \href{https://zenodo.org/records/24221202}{Zenodo 24221202} (2026).

\bibitem{Ban09}
T. Banica, Cyclotomic expansion of exceptional spectral measures, {\em Internat. J. Math.} {\bf 20} (2009), 275--297.

\bibitem{BBi15}
T. Banica and J. Bichon, Random walk questions for linear quantum groups, {\em Int. Math. Res. Not. IMRN} (2015), no. 24, 13406--13436.

\bibitem{BBS09}
T. Banica, J. Bichon and J.-M. Schlenker, Representations of quantum permutation algebras, {\em J. Funct. Anal.} {\bf 257} (2009), 2864--2910.

\bibitem{BB07}
T. Banica and D. Bisch, Spectral measures of small index principal graphs, {\em Comm. Math. Phys.} {\bf 269} (2007), 259--281.

\bibitem{BCZ09}
T. Banica, B. Collins and P. Zinn-Justin, Spectral analysis of the free orthogonal matrix, {\em Int. Math. Res. Not. IMRN} (2009), no. 17, 3286--3309.

\bibitem{BO26}
T. Banica and C. Opus, Higher invariants of complex Hadamard matrices, \href{https://hexagonmath.org/2611.00418v1}{Hexagon 2611.00418v1} (2026).

\bibitem{BP26}
T. Banica and S. Palcoux, Invariant tensor dimensions as Hamburger moments, \href{https://zenodo.org/records/24348893}{Zenodo 24348893} (2026).

\bibitem{Bic04}
J. Bichon, Free wreath product by the quantum permutation group, {\em Algebr. Represent. Theory} {\bf 7} (2004), 343--362.

\bibitem{Bic15}
J. Bichon, Quotients and Hopf images of a smash coproduct, {\em Tsukuba J. Math.} {\bf 39} (2015), 285--310.

\bibitem{BSWZ12}
J. M. Borwein, A. Straub, J. Wan and W. Zudilin, Densities of short uniform random walks, with an appendix by D. Zagier, {\em Canad. J. Math.} {\bf 64} (2012), 961--990.

\bibitem{Bur15}
R. D. Burstein, Group-type subfactors and Hadamard matrices, {\em Trans. Amer. Math. Soc.} {\bf 367} (2015), 6783--6807.

\bibitem{CGPNG09}
A. H. Castro Neto, F. Guinea, N. M. R. Peres, K. S. Novoselov and A. K. Geim, The electronic properties of graphene, {\em Rev. Mod. Phys.} {\bf 81} (2009), 109--162.

\bibitem{CM80}
H. S. M. Coxeter and W. O. J. Moser, Generators and relations for discrete groups, 4th ed., {\em Ergeb. Math. Grenzgeb.} {\bf 14}, Springer, 1980.

\bibitem{Del96}
P. Deligne, La s\'erie exceptionnelle de groupes de Lie, {\em C. R. Acad. Sci. Paris S\'er. I Math.} {\bf 322} (1996), 321--326.

\bibitem{DDM96}
P. Deligne and R. de Man, La s\'erie exceptionnelle de groupes de Lie II, {\em C. R. Acad. Sci. Paris S\'er. I Math.} {\bf 323} (1996), 577--582.

\bibitem{DK94}
M. S. Dijkhuizen and T. H. Koornwinder, CQG algebras: a direct algebraic approach to compact quantum groups, {\em Lett. Math. Phys.} {\bf 32} (1994), 315--330.

\bibitem{Dit04}
P. Di\c t\u a, Some results on the parametrization of complex Hadamard matrices, {\em J. Phys. A} {\bf 37} (2004), 5355--5374.

\bibitem{EP10}
D. E. Evans and M. Pugh, Spectral measures and generating series for nimrep graphs in subfactor theory, {\em Comm. Math. Phys.} {\bf 295} (2010), 363--413.

\bibitem{EP11}
D. E. Evans and M. Pugh, Spectral measures and generating series for nimrep graphs in subfactor theory II: $SU(3)$, {\em Comm. Math. Phys.} {\bf 301} (2011), 771--809.

\bibitem{EP15}
D. E. Evans and M. Pugh, Spectral measures for $G_2$, {\em Comm. Math. Phys.} {\bf 337} (2015), 1161--1197.

\bibitem{GHJ89}
F. M. Goodman, P. de la Harpe and V. F. R. Jones, Coxeter graphs and towers of algebras, {\em Math. Sci. Res. Inst. Publ.} {\bf 14}, Springer, 1989.

\bibitem{Haa94}
U. Haagerup, Principal graphs of subfactors in the index range $4<[M:N]<3+\sqrt2$, in {\em Subfactors (Kyuzeso, 1993)}, World Sci. Publ., 1994, 1--38.

\bibitem{HN53}
J. P. Hobson and W. A. Nierenberg, The statistics of a two-dimensional, hexagonal net, {\em Phys. Rev.} {\bf 89} (1953), 662.

\bibitem{Izu01}
M. Izumi, The structure of sectors associated with Longo--Rehren inclusions. II. Examples, {\em Rev. Math. Phys.} {\bf 13} (2001), 603--674.

\bibitem{Jon83}
V. F. R. Jones, Index for subfactors, {\em Invent. Math.} {\bf 72} (1983), 1--25.

\bibitem{Jon01}
V. F. R. Jones, The annular structure of subfactors, in {\em Essays on Geometry and Related Topics}, {\em Monogr. Enseign. Math.} {\bf 38} (2001), 401--463.

\bibitem{Jon99}
V. F. R. Jones, Planar algebras, I, {\em New Zealand J. Math.} {\bf 52} (2021), 1--107.

\bibitem{JMS14}
V. F. R. Jones, S. Morrison and N. Snyder, The classification of subfactors of index at most $5$, {\em Bull. Amer. Math. Soc. (N.S.)} {\bf 51} (2014), 277--327.

\bibitem{JS97}
V. Jones and V. S. Sunder, Introduction to subfactors, {\em London Math. Soc. Lecture Note Ser.} {\bf 234}, Cambridge University Press, 1997.

\bibitem{KS01}
M. Kotani and T. Sunada, Standard realizations of crystal lattices via harmonic maps, {\em Trans. Amer. Math. Soc.} {\bf 353} (2001), 1--20.

\bibitem{LM02}
J. M. Landsberg and L. Manivel, Triality, exceptional Lie algebras and Deligne dimension formulas, {\em Adv. Math.} {\bf 171} (2002), 59--85.

\bibitem{Mon23}
M. Montgomery, Colored planar algebras for commuting squares and applications to Hadamard subfactors, {\em J. Funct. Anal.} {\bf 285} (2023), 110010.

\bibitem{NT13}
S. Neshveyev and L. Tuset, Compact quantum groups and their representation categories, {\em Cours Sp\'ecialis\'es} {\bf 20}, Soci\'et\'e Math\'ematique de France, 2013.

\bibitem{Nic06}
R. Nicoara, A finiteness result for commuting squares of matrix algebras, {\em J. Operator Theory} {\bf 55} (2006), 295--310.

\bibitem{Nic10}
R. Nicoara, Subfactors and Hadamard matrices, {\em J. Operator Theory} {\bf 64} (2010), 453--468.

\bibitem{Sun13}
T. Sunada, Topological crystallography. With a view towards discrete geometric analysis, {\em Surv. Tutor. Appl. Math. Sci.} {\bf 6}, Springer, 2013.

\end{thebibliography}
\end{document}
